% ARCHIVO GENERADO por generate_software_doc.py -- no editar a mano.
\chapter{Software}\label{chap:software}
\lstset{language=Python,basicstyle=\ttfamily\scriptsize,numbers=left,numberstyle=\tiny\color{gray},breaklines=true,frame=single,framerule=0.3pt,rulecolor=\color{gray!50},keywordstyle=\color{blue!70!black},commentstyle=\color{green!40!black},stringstyle=\color{red!60!black},showstringspaces=false,columns=fullflexible,keepspaces=true,inputencoding=utf8,extendedchars=true,literate={°}{{\textdegree{}}}1 {±}{{\(\pm\)}}1 {·}{{\(\cdot\)}}1 {¿}{{?`}}1 {Á}{{\'A}}1 {Í}{{\'I}}1 {Ó}{{\'O}}1 {Ú}{{\'U}}1 {á}{{\'a}}1 {é}{{\'e}}1 {í}{{\'i}}1 {ñ}{{\~n}}1 {ó}{{\'o}}1 {ú}{{\'u}}1 {↔}{{\(\leftrightarrow\)}}1 {↕}{{\(\updownarrow\)}}1}
Este capítulo documenta el software del robot Rover X (mBot2 con kit Rover, MicroPython en el CyberPi) del equipo Robótica 76 para MakeX 2026. Se genera automáticamente a partir de los comentarios Doxygen del código (carpeta software/) y de los datos del propio programa: misiones, medidas de la pista y parámetros ajustables.

\section{Descripci\'on general}
El robot usa \textbf{dos programas}, uno para cada parte de la competencia. Cada uno se carga en el CyberPi como \textbf{un solo archivo} (carpeta \texttt{mblock/}), que \texttt{tools/scripts/bundle\_for\_mblock.py} arma juntando los m\'odulos de \texttt{software/}. El archivo empaquetado solo importa m\'odulos de Makeblock (\texttt{cyberpi}, \texttt{mbot2}, \texttt{mbuild}, \texttt{gamepad}) y de MicroPython (\texttt{time}, \texttt{math}, \texttt{json}).

\begin{longtable}{@{}p{0.27\linewidth}p{0.67\linewidth}@{}}\toprule
\textbf{Archivo} & \textbf{Qu\'e hace} \\\midrule\endhead
\multicolumn{2}{@{}l}{\emph{Programas para el robot} (\texttt{software/robot})} \\
\hyperref[mod:rover-x-auto]{\texttt{rover\_\allowbreak{}x\_\allowbreak{}auto.py}} & Programa AUTÓNOMO del robot (se carga en el CyberPi como mblock/rover\_x\_auto.py). \\
\hyperref[mod:rover-x-manual]{\texttt{rover\_\allowbreak{}x\_\allowbreak{}manual.py}} & Programa MANUAL del robot (se carga en el CyberPi como mblock/rover\_x\_manual.py). \\
\multicolumn{2}{@{}l}{\emph{Autónomo} (\texttt{software/autonomous})} \\
\hyperref[mod:autonomous]{\texttt{autonomous.py}} & Ejecutor del área automática: corre las misiones de missions.py. \\
\hyperref[mod:missions]{\texttt{missions.py}} & Misiones del área automática, escritas como listas de pasos (la RUTA). \\
\multicolumn{2}{@{}l}{\emph{Control manual} (\texttt{software/manual})} \\
\hyperref[mod:manual]{\texttt{manual.py}} & Control manual con el Bluetooth Controller de Makeblock (misiones de alianza). \\
\multicolumn{2}{@{}l}{\emph{Control} (\texttt{software/controls})} \\
\hyperref[mod:motion]{\texttt{motion.py}} & Movimientos con giroscopio y encoders (sin seguir la línea). \\
\hyperref[mod:linefollow]{\texttt{linefollow.py}} & Seguidor de línea PID con detección de cruces robusta para línea delgada. \\
\hyperref[mod:odometry]{\texttt{odometry.py}} & Odometría de posición (x, y, rumbo) con encoders + giroscopio, robusta al patinaje. \\
\hyperref[mod:gripper]{\texttt{gripper.py}} & Garra (servo de cierre) y brazo (servo de altura). \\
\multicolumn{2}{@{}l}{\emph{Base (core)} (\texttt{software/core})} \\
\hyperref[mod:hal]{\texttt{hal.py}} & Capa de hardware (HAL): ÚNICO lugar que llama a la API del mBot2. \\
\hyperref[mod:config]{\texttt{config.py}} & Configuración del robot (mBot2 + orugas + garra de dos servos): puertos, signos y parámetros ajustables cfg.P. \\
\hyperref[mod:pista]{\texttt{pista.py}} & MEDIDAS DE LA PISTA que usa el programa AUTÓNOMO, por nivel. \\
\hyperref[mod:geometry]{\texttt{geometry.py}} & Cuentas que dependen de la geometría del robot (cfg.P). \\
\hyperref[mod:colors]{\texttt{colors.py}} & Reconocimiento de colores: zonas de inicio, marcas verdes y tarjetas. \\
\hyperref[mod:settings]{\texttt{settings.py}} & Menú de ajustes en el CyberPi, calibración de línea y guardado de parámetros. \\
\bottomrule\end{longtable}
\begin{figure}[H]\centering\resizebox{\linewidth}{!}{%
\begin{tikzpicture}[>=Stealth,
  md/.style={draw,rounded corners,fill=blue!8,font=\footnotesize\ttfamily,minimum height=0.65cm,minimum width=2.1cm,inner sep=3pt},
  ly/.style={font=\footnotesize\bfseries,anchor=east}]
\node[ly] at (-1.6,0.00) {Programas};
\node[md] (roverxauto) at (5.00,0.00) {rover\_x\_auto};
\node[md] (roverxmanual) at (7.50,0.00) {rover\_x\_manual};
\node[ly] at (-1.6,-1.70) {Autónomo / manual};
\node[md] (autonomous) at (3.75,-1.70) {autonomous};
\node[md] (missions) at (6.25,-1.70) {missions};
\node[md] (manual) at (8.75,-1.70) {manual};
\node[ly] at (-1.6,-3.40) {Control};
\node[md] (motion) at (2.50,-3.40) {motion};
\node[md] (linefollow) at (5.00,-3.40) {linefollow};
\node[md] (odometry) at (7.50,-3.40) {odometry};
\node[md] (gripper) at (10.00,-3.40) {gripper};
\node[ly] at (-1.6,-5.10) {Base};
\node[md] (settings) at (0.00,-5.10) {settings};
\node[md] (colors) at (2.50,-5.10) {colors};
\node[md] (geometry) at (5.00,-5.10) {geometry};
\node[md] (pista) at (7.50,-5.10) {pista};
\node[md] (config) at (10.00,-5.10) {config};
\node[md] (hal) at (12.50,-5.10) {hal};
\draw[->,gray!70] (roverxauto) to (autonomous);
\draw[->,gray!70] (roverxauto) to (config);
\draw[->,gray!70] (roverxauto) to (gripper);
\draw[->,gray!70] (roverxauto) to (hal);
\draw[->,gray!70] (roverxauto) to (settings);
\draw[->,gray!70] (roverxmanual) to (config);
\draw[->,gray!70] (roverxmanual) to (gripper);
\draw[->,gray!70] (roverxmanual) to (hal);
\draw[->,gray!70] (roverxmanual) to (manual);
\draw[->,gray!70] (roverxmanual) to (settings);
\draw[->,gray!70] (autonomous) to (colors);
\draw[->,gray!70] (autonomous) to (geometry);
\draw[->,gray!70] (autonomous) to (gripper);
\draw[->,gray!70] (autonomous) to (linefollow);
\draw[->,gray!70,bend left=25] (autonomous) to (missions);
\draw[->,gray!70] (autonomous) to (motion);
\draw[->,gray!70] (autonomous) to (odometry);
\draw[->,gray!70] (missions) to (pista);
\draw[->,gray!70] (manual) to (geometry);
\draw[->,gray!70] (motion) to (colors);
\draw[->,gray!70] (motion) to (geometry);
\draw[->,gray!70] (linefollow) to (colors);
\draw[->,gray!70] (linefollow) to (geometry);
\draw[->,gray!70,bend left=25] (linefollow) to (motion);
\draw[->,gray!70] (odometry) to (geometry);
\draw[->,gray!70,bend left=25] (odometry) to (motion);
\draw[->,gray!70] (gripper) to (geometry);
\draw[->,gray!70,bend left=25] (hal) to (geometry);
\draw[->,gray!70,bend left=25] (config) to (pista);
\draw[->,gray!70] (settings) to (gripper);
\draw[->,gray!70] (settings) to (motion);
\end{tikzpicture}
}\caption{Dependencias entre m\'odulos, generado del c\'odigo: una flecha va del m\'odulo que importa al importado. Casi todos usan \texttt{config} y \texttt{hal}; esas flechas solo se dibujan desde los programas.}\end{figure}
\section{Programas para el robot y botones}
\subsection{Autónomo: \texttt{mblock/rover\_x\_auto.py}}
\begin{center}\begin{tabular}{@{}ll@{}}\toprule
\textbf{Bot\'on del CyberPi} & \textbf{Acci\'on} \\\midrule
Joystick arriba & elegir el nivel 1 \\
Joystick abajo & elegir el nivel 2 \\
Joystick izquierda & elegir el nivel 3 \\
Joystick derecha & elegir el nivel 4 \\
B & empezar el tramo del nivel elegido \\
A & parar el robot (para pedir reinicio) \\
Joystick centro & men\'u de ajustes y calibraci\'on (en la pantalla de inicio) \\
\bottomrule\end{tabular}\end{center}
\subsection{Manual: \texttt{mblock/rover\_x\_manual.py}}
\begin{center}\begin{tabular}{@{}ll@{}}\toprule
\textbf{Bot\'on del CyberPi} & \textbf{Acci\'on} \\\midrule
B & empezar el control con el gamepad \\
A & parar (el robot se detiene y espera otra vez B) \\
\bottomrule\end{tabular}\end{center}
Mandos del gamepad (de \texttt{manual.py}):

Misiones manuales (reglamento): cubos M-A-K-E-X en la barra del nivel 2 (en orden y de frente), activar las rampas (el primer acceso es por la ESCALERA) y 4 cilindros (5 x 15 cm) parados en el área negra.

Mapeo (ver docs/control\_manual.md):

{\small\begin{alltt}
Joystick izq. \(\updownarrow\) (Ly)  avanzar / retroceder      Joystick der. \(\leftrightarrow\) (Rx)  girar
Cruceta               movimientos finos lentos (si los joysticks están en el centro)
L1 / L2               brazo arriba / abajo (mantener)
R1 / R2               garra abrir / cerrar (mantener; entre claw\_open y claw\_cube - 20\textdegree{})
N1                    modo rápido / preciso (LEDs azul / naranja)
N2                    brazo a la altura de la barra (cubos MakeX)
N3 (mantener)         escalera: recto con el giroscopio, potencia de escalera
N4                    brazo arriba (garra sobre la nariz): viajar, girar cerca de cosas
Joystick der. \(\updownarrow\) (Ry)  brazo fino (además de L1/L2)
\end{alltt}}
Ayudas:

\begin{itemize}\setlength\itemsep{0pt}
\item Al avanzar sin girar, el giroscopio SOSTIENE EL RUMBO (rampas, recortes, la oruga izquierda que avanza 2.3 \% menos). Al soltar el giro se toma el rumbo nuevo.
\item hal.drive limita la aceleración mientras va inclinado (no se vuelca en rampas/escalera).
\item El brazo se mueve por ÁNGULO DEL SERVO (0..90\textdegree{}, tabla medida en config.py): a 0\textdegree{} la garra queda baja y al frente (\textasciitilde{}8.5 cm), a 90\textdegree{} arriba sobre la nariz.
\end{itemize}
\section{Misiones del aut\'onomo}
La ruta est\'a en \texttt{missions.py} como una lista de misiones; cada misi\'on es una lista de pasos que ejecuta \texttt{autonomous.run\_step()} (secci\'on \ref{fn:autonomous:run-step}). \emph{Cr\'itica}: si falla, el tramo se detiene. \emph{Fin de tramo}: el robot se detiene ah\'i y el equipo pide reinicio en el siguiente nivel.
\begin{longtable}{@{}rrrlccr@{}}\toprule
\textbf{Misi\'on} & \textbf{Tramo} & \textbf{Pts} & \textbf{Nombre} & \textbf{Cr\'itica} & \textbf{Fin de tramo} & \textbf{Pasos} \\\midrule\endhead
1 & 1 & 50 & Seguir linea & s\'i &  & 11 \\
2 & 1 & 50 & Tomar cubo &  &  & 6 \\
3 & 1 & 50 & Subir nivel 2 & s\'i &  & 9 \\
4 & 1 & 100 & Dejar cubo &  & s\'i & 9 \\
5 & 2 & 100 & Zona obstaculos & s\'i &  & 15 \\
6 & 2 & 100 & Subir nivel 3 & s\'i & s\'i & 2 \\
7 & 3 & 100 & Mover prisma &  &  & 14 \\
9 & 3 & 150 & Escalera nivel 4 & s\'i &  & 10 \\
10 & 4 & 100 & Evitar torre & s\'i &  & 7 \\
11 & 4 & 120 & Colores &  &  & 1 \\
13 & 4 & 100 & Meta azul &  &  & 1 \\
\midrule & & 1020 & \multicolumn{4}{l}{puntos posibles de estas misiones} \\
\bottomrule\end{longtable}
\section{Medidas de la pista}\label{sec:pista}
Aquí están TODAS las medidas de la pista de las que depende la ruta, separadas por nivel. Cada nivel tiene dos partes:

\begin{itemize}\setlength\itemsep{0pt}
\item "Mediciones que necesita el robot": se miden con flexómetro en la pista real.
\item "Ajustes de la maniobra": retrocesos, holguras y rumbos; se afinan probando en la pista.
\end{itemize}
Los valores de ahora salen del plano del reglamento (simulation/arena.py) y se probaron en el simulador. Los marcados  \# MEDIR  hay que confirmarlos en la pista; al medirlos, cambiar "MEDIR" por "MEDIDO".

Después de cambiar algo:

{\small\begin{alltt}
1) python3 -m pytest              (que todo siga pasando)
2) python3 simulation/run\_auto.py (ver la ruta en el simulador)
3) python3 tools/scripts/bundle\_for\_mblock.py --all   y volver a cargar mblock/rover\_x\_auto.py
\end{alltt}}
Las medidas marcadas [MENU] también se pueden cambiar desde el menú de ajustes del CyberPi (grupo "Pista"). OJO: si ya se guardaron desde el menú, en el robot manda lo guardado y no lo que diga este archivo (para volver a este archivo: menú -> "Restaurar valores").

UNIDADES Y MAPA

\begin{itemize}\setlength\itemsep{0pt}
\item Distancias en cm, ángulos en grados.
\item Posiciones (x, y) en el mapa de la vista superior del reglamento: x hacia la DERECHA desde la orilla izquierda de la arena, y hacia ABAJO desde la orilla de arriba.
\item Rumbos: 0 = derecha (este), 90 = abajo (sur), 180 = izquierda (oeste), -90 = arriba (norte).
\item "Centro de giro" = punto medio de las orugas; las sondas de línea van 5.2 cm delante de él.
\end{itemize}
Medidas del robot que importan para las holguras: 21.35 cm de ancho x 20.3 de largo; girando en el lugar las esquinas barren 16 cm (nariz) y 13.6 cm (cola).

\subsection{Nivel 1 (Zona Roja)}
Misiones 1 a 3: seguir la línea hasta el cruce del cubo, tomar el cubo, subir la rampa gris. Se recorre POR COMANDOS (avanzar X cm con el rumbo sostenido por el giroscopio).

\paragraph{Mediciones que necesita el robot}
\begin{longtable}{@{}p{0.36\linewidth}rlp{0.42\linewidth}@{}}\toprule
\textbf{Variable} & \textbf{Valor} & \textbf{Estado} & \textbf{Descripci\'on} \\\midrule\endhead
\texttt{nivel\_\allowbreak{}1\_\allowbreak{}inicio\_\allowbreak{}x} & 12.6 & MEDIR & x del centro de giro al colocar el robot en la \\
\texttt{nivel\_\allowbreak{}1\_\allowbreak{}inicio\_\allowbreak{}y} & 12.9 & MEDIR & zona roja (sobre la X, mirando al sur) \\
\texttt{nivel\_\allowbreak{}1\_\allowbreak{}inicio\_\allowbreak{}a\_\allowbreak{}linea\_\allowbreak{}inferior\_\allowbreak{}cm} & 33.8 & MEDIR & de la zona roja (centro de giro) al cruce con la línea de abajo, yendo al sur \\
\texttt{nivel\_\allowbreak{}1\_\allowbreak{}linea\_\allowbreak{}inferior\_\allowbreak{}y} & 43.2 & MEDIR & y del centro de giro cuando las sondas VEN la línea de abajo (ahí se gira; la pared está en y 61) \\
\texttt{nivel\_\allowbreak{}1\_\allowbreak{}cruce\_\allowbreak{}a\_\allowbreak{}cruce\_\allowbreak{}del\_\allowbreak{}cubo\_\allowbreak{}cm} & 59.0 & MEDIR & del cruce de abajo hasta tener las sondas sobre el cruce frente al cubo, yendo al este \\
\texttt{nivel\_\allowbreak{}1\_\allowbreak{}cruce\_\allowbreak{}a\_\allowbreak{}pedestal\_\allowbreak{}cm} & 12.4 & MEDIR menú & de donde las sondas ven el cruce (su orilla cercana) al CENTRO del pedestal del cubo \\
\texttt{nivel\_\allowbreak{}1\_\allowbreak{}linea\_\allowbreak{}del\_\allowbreak{}cubo\_\allowbreak{}y} & 49.0 & MEDIR & y de la línea del cruce frente al cubo \\
\texttt{nivel\_\allowbreak{}1\_\allowbreak{}cubo\_\allowbreak{}a\_\allowbreak{}linea\_\allowbreak{}de\_\allowbreak{}rampa\_\allowbreak{}cm} & 18.4 & MEDIR & del cruce frente al cubo, al norte, hasta la altura de la línea que sube a la rampa \\
\texttt{nivel\_\allowbreak{}1\_\allowbreak{}linea\_\allowbreak{}a\_\allowbreak{}pie\_\allowbreak{}de\_\allowbreak{}rampa\_\allowbreak{}cm} & 18.7 & MEDIR & de ahí, al este, al pie de la rampa gris \\
\bottomrule\end{longtable}
\paragraph{Ajustes de la maniobra}
\begin{longtable}{@{}p{0.36\linewidth}rlp{0.42\linewidth}@{}}\toprule
\textbf{Variable} & \textbf{Valor} & \textbf{Estado} & \textbf{Descripci\'on} \\\midrule\endhead
\texttt{nivel\_\allowbreak{}1\_\allowbreak{}x\_\allowbreak{}de\_\allowbreak{}bajada} & 15.4 &  & x por la que baja: 2.8 cm a la derecha de la línea, porque la cola barre 13.6 al girar y la orilla está en x 0 \\
\texttt{nivel\_\allowbreak{}1\_\allowbreak{}y\_\allowbreak{}hacia\_\allowbreak{}cubo} & 47.4 &  & y de la recta hacia el cubo (pedestal en y 45.6..51.8): así el pivote de después pasa a \textasciitilde{}1.5 cm de la pared \\
\texttt{nivel\_\allowbreak{}1\_\allowbreak{}holgura\_\allowbreak{}cruce\_\allowbreak{}cm} & 6.0 &  & cada tramo se queda corto esto y termina buscando el cruce \\
\texttt{nivel\_\allowbreak{}1\_\allowbreak{}retroceso\_\allowbreak{}antes\_\allowbreak{}de\_\allowbreak{}girar\_\allowbreak{}cm} & 1.5 &  & el robot se pasa \textasciitilde{}1 cm al frenar sobre la línea \\
\texttt{nivel\_\allowbreak{}1\_\allowbreak{}retroceso\_\allowbreak{}para\_\allowbreak{}pivotear\_\allowbreak{}cm} & 11.0 &  & tras el cubo, para pivotear al norte sin barrer el pedestal \\
\texttt{nivel\_\allowbreak{}1\_\allowbreak{}retroceso\_\allowbreak{}tras\_\allowbreak{}cubo\_\allowbreak{}cm} & 4.5 &  & el retroceso que ya hizo la misión 2 (lo descuenta) \\
\bottomrule\end{longtable}
\subsection{Nivel 2 (Zona Amarilla)}
Misión 4 (dejar el cubo en el pedestal), 5 (cruzar la zona de obstáculos en diagonal, sin línea) y 6 (subir la rampa blanca al nivel 3).

\paragraph{Mediciones que necesita el robot}
\begin{longtable}{@{}p{0.36\linewidth}rlp{0.42\linewidth}@{}}\toprule
\textbf{Variable} & \textbf{Valor} & \textbf{Estado} & \textbf{Descripci\'on} \\\midrule\endhead
\texttt{nivel\_\allowbreak{}2\_\allowbreak{}t\_\allowbreak{}a\_\allowbreak{}pedestal\_\allowbreak{}cm} & 8.8 & MEDIR & línea seguida desde la T hasta dejar la garra sobre el centro del pedestal (con 9.6 la nariz quedaba a 0.3 cm) \\
\texttt{nivel\_\allowbreak{}2\_\allowbreak{}colocado\_\allowbreak{}x} & 131.9 & MEDIR & centro de giro al colocarlo en la T, retroceder y \\
\texttt{nivel\_\allowbreak{}2\_\allowbreak{}colocado\_\allowbreak{}y} & 30.4 & MEDIR & girar al sur (arranque o reinicio en el nivel 2) \\
\texttt{nivel\_\allowbreak{}2\_\allowbreak{}tras\_\allowbreak{}dejar\_\allowbreak{}cubo\_\allowbreak{}x} & 133.9 & MEDIR & centro de giro al terminar la misión 4 \\
\texttt{nivel\_\allowbreak{}2\_\allowbreak{}tras\_\allowbreak{}dejar\_\allowbreak{}cubo\_\allowbreak{}y} & 26.8 & MEDIR & (mirando al sur) \\
\texttt{nivel\_\allowbreak{}2\_\allowbreak{}giro\_\allowbreak{}al\_\allowbreak{}pasillo\_\allowbreak{}x} & 134.2 & MEDIR & punto donde gira al este para entrar al pasillo: \\
\texttt{nivel\_\allowbreak{}2\_\allowbreak{}giro\_\allowbreak{}al\_\allowbreak{}pasillo\_\allowbreak{}y} & 43.5 & MEDIR & y <= 61 - 16 - 1.5 (la nariz barre 16 cm) \\
\texttt{nivel\_\allowbreak{}2\_\allowbreak{}abertura\_\allowbreak{}izquierda\_\allowbreak{}y} & 45.0 & MEDIR & centro de la abertura del muro izquierdo (y 30.5..60.2) \\
\texttt{nivel\_\allowbreak{}2\_\allowbreak{}abertura\_\allowbreak{}derecha\_\allowbreak{}y} & 15.9 & MEDIR & centro de la abertura del muro derecho (y 0.9..30.5) \\
\texttt{nivel\_\allowbreak{}2\_\allowbreak{}diagonal\_\allowbreak{}punto\_\allowbreak{}x} & 212.7 & MEDIR & la diagonal pasa por el centro de la abertura derecha \\
\texttt{nivel\_\allowbreak{}2\_\allowbreak{}diagonal\_\allowbreak{}punto\_\allowbreak{}y} & 16.0 & MEDIR &  \\
\texttt{nivel\_\allowbreak{}2\_\allowbreak{}diagonal\_\allowbreak{}deteccion\_\allowbreak{}x} & 208.96 & MEDIR & centro de giro cuando la sonda L2 toca la línea de \\
\texttt{nivel\_\allowbreak{}2\_\allowbreak{}diagonal\_\allowbreak{}deteccion\_\allowbreak{}y} & 17.82 & MEDIR & salida yendo por la diagonal \\
\texttt{nivel\_\allowbreak{}2\_\allowbreak{}x\_\allowbreak{}al\_\allowbreak{}ver\_\allowbreak{}diagonal} & 218.6 & MEDIR & x del centro de giro cuando las sondas centrales ven la diagonal de la línea yendo al este por la abertura derecha \\
\texttt{nivel\_\allowbreak{}2\_\allowbreak{}rampa\_\allowbreak{}centro\_\allowbreak{}x} & 228.15 & MEDIR & centro de la rampa blanca (x 216.7..239.6) \\
\texttt{nivel\_\allowbreak{}2\_\allowbreak{}rampa\_\allowbreak{}linea\_\allowbreak{}x} & 228.6 & MEDIR & la línea que baja por ella (0.45 cm al este del centro) \\
\bottomrule\end{longtable}
\paragraph{Ajustes de la maniobra}
\begin{longtable}{@{}p{0.36\linewidth}rlp{0.42\linewidth}@{}}\toprule
\textbf{Variable} & \textbf{Valor} & \textbf{Estado} & \textbf{Descripci\'on} \\\midrule\endhead
\texttt{nivel\_\allowbreak{}2\_\allowbreak{}retroceso\_\allowbreak{}en\_\allowbreak{}t\_\allowbreak{}cm} & 3.0 &  & retroceso en la T antes de girar al sur (la nariz barre 16 y el muro del pasillo está a \textasciitilde{}17) \\
\texttt{nivel\_\allowbreak{}2\_\allowbreak{}retroceso\_\allowbreak{}tras\_\allowbreak{}dejar\_\allowbreak{}cubo\_\allowbreak{}cm} & 6.0 &  &  \\
\texttt{nivel\_\allowbreak{}2\_\allowbreak{}diagonal\_\allowbreak{}rumbo} & -26.0 &  & rumbo de la diagonal (noreste): una sola recta por el centro de las dos aberturas, 2.3 cm a las puntas de los muros \\
\texttt{nivel\_\allowbreak{}2\_\allowbreak{}diagonal\_\allowbreak{}inicio\_\allowbreak{}x} & 150.5 &  & x del centro de giro al pivotear para quedar sobre la diagonal \\
\texttt{nivel\_\allowbreak{}2\_\allowbreak{}diagonal\_\allowbreak{}max\_\allowbreak{}cm} & 110.0 &  & tope por encoders (sobre los recortes patina: miden de más) \\
\texttt{nivel\_\allowbreak{}2\_\allowbreak{}diagonal\_\allowbreak{}vel} & 14 &  & \% potencia: lento y con el rumbo firme sobre los recortes \\
\texttt{nivel\_\allowbreak{}2\_\allowbreak{}diagonal\_\allowbreak{}kp} & 2.5 &  & firmeza del rumbo; si oscila en el robot real, bajarlo \\
\texttt{nivel\_\allowbreak{}2\_\allowbreak{}x\_\allowbreak{}giro\_\allowbreak{}salida} & 223.0 &  & girar cuando la cola rebasa la punta del muro (x 212.7) \\
\texttt{nivel\_\allowbreak{}2\_\allowbreak{}rumbo\_\allowbreak{}hacia\_\allowbreak{}rampa} & 60 &  & al salir del pasillo: las sondas quedan junto a la línea \\
\texttt{nivel\_\allowbreak{}2\_\allowbreak{}avance\_\allowbreak{}hacia\_\allowbreak{}rampa\_\allowbreak{}cm} & 5.5 &  & en diagonal hasta x \textasciitilde{}226.6 (esquina a \textasciitilde{}1.3 cm de la orilla) \\
\texttt{nivel\_\allowbreak{}2\_\allowbreak{}rampa\_\allowbreak{}retroceso\_\allowbreak{}cm} & 8.0 &  & la nariz ya estaba sobre el pie de la rampa: centrarse en plano \\
\bottomrule\end{longtable}
\subsection{Nivel 3 (Zona Azul)}
Misión 6 (termina de entrar al nivel 3), 7 (mover el prisma a la zona punteada) y 9 (ir a la escalera). La misión 8 (pelota) está omitida.

\paragraph{Mediciones que necesita el robot}
\begin{longtable}{@{}p{0.36\linewidth}rlp{0.42\linewidth}@{}}\toprule
\textbf{Variable} & \textbf{Valor} & \textbf{Estado} & \textbf{Descripci\'on} \\\midrule\endhead
\texttt{nivel\_\allowbreak{}3\_\allowbreak{}zona\_\allowbreak{}azul\_\allowbreak{}a\_\allowbreak{}dentro\_\allowbreak{}cm} & 19.2 & MEDIR & avance desde que las sondas ven la zona azul hasta quedar completamente dentro (antes de la curva en S) \\
\texttt{nivel\_\allowbreak{}3\_\allowbreak{}orilla\_\allowbreak{}zona\_\allowbreak{}azul\_\allowbreak{}x} & 228.8 & MEDIR & centro de giro al salir de la zona azul (sondas en \\
\texttt{nivel\_\allowbreak{}3\_\allowbreak{}orilla\_\allowbreak{}zona\_\allowbreak{}azul\_\allowbreak{}y} & 70.8 & MEDIR & la orilla de la zona) \\
\texttt{nivel\_\allowbreak{}3\_\allowbreak{}orilla\_\allowbreak{}a\_\allowbreak{}prisma\_\allowbreak{}cm} & 3.6 & MEDIR & de la orilla de la zona azul a la altura del prisma \\
\texttt{nivel\_\allowbreak{}3\_\allowbreak{}rumbo\_\allowbreak{}a\_\allowbreak{}zona\_\allowbreak{}punteada} & 107 & MEDIR & rumbo del prisma hacia la zona punteada \\
\texttt{nivel\_\allowbreak{}3\_\allowbreak{}prisma\_\allowbreak{}a\_\allowbreak{}zona\_\allowbreak{}punteada\_\allowbreak{}cm} & 22.0 & MEDIR & avance con el prisma hasta dejarlo en la zona \\
\texttt{nivel\_\allowbreak{}3\_\allowbreak{}frente\_\allowbreak{}escalera\_\allowbreak{}x} & 166.1 & MEDIR & punto al sur del tramo recto de la línea, frente a \\
\texttt{nivel\_\allowbreak{}3\_\allowbreak{}frente\_\allowbreak{}escalera\_\allowbreak{}y} & 96.4 & MEDIR & la escalera \\
\bottomrule\end{longtable}
\paragraph{Ajustes de la maniobra}
\begin{longtable}{@{}p{0.36\linewidth}rlp{0.42\linewidth}@{}}\toprule
\textbf{Variable} & \textbf{Valor} & \textbf{Estado} & \textbf{Descripci\'on} \\\midrule\endhead
\texttt{nivel\_\allowbreak{}3\_\allowbreak{}retroceso\_\allowbreak{}antes\_\allowbreak{}de\_\allowbreak{}bajar\_\allowbreak{}garra\_\allowbreak{}cm} & 4.0 &  & para bajar la garra sin tocar el prisma \\
\texttt{nivel\_\allowbreak{}3\_\allowbreak{}avance\_\allowbreak{}para\_\allowbreak{}tomar\_\allowbreak{}prisma\_\allowbreak{}cm} & 6.5 &  & al oeste hasta tener el prisma entre los dedos \\
\texttt{nivel\_\allowbreak{}3\_\allowbreak{}retroceso\_\allowbreak{}tras\_\allowbreak{}soltar\_\allowbreak{}cm} & 8.0 &  &  \\
\texttt{nivel\_\allowbreak{}3\_\allowbreak{}retroceso\_\allowbreak{}antes\_\allowbreak{}de\_\allowbreak{}escalera\_\allowbreak{}cm} & 5.0 &  & para que las sondas busquen la línea hacia adelante \\
\bottomrule\end{longtable}
\paragraph{Pelota (misión 8): OMITIDA por ahora; se guardan las medidas para después}
\begin{longtable}{@{}p{0.36\linewidth}rlp{0.42\linewidth}@{}}\toprule
\textbf{Variable} & \textbf{Valor} & \textbf{Estado} & \textbf{Descripci\'on} \\\midrule\endhead
\texttt{nivel\_\allowbreak{}3\_\allowbreak{}pelota\_\allowbreak{}linea\_\allowbreak{}a\_\allowbreak{}pedestal\_\allowbreak{}cm} & 23.6 & MEDIR menú & del centro de la línea (en la marca roja) al centro del pedestal de la pelota \\
\texttt{nivel\_\allowbreak{}3\_\allowbreak{}pelota\_\allowbreak{}rumbo} & 113 & MEDIR & rumbo hacia el pedestal de la pelota \\
\texttt{nivel\_\allowbreak{}3\_\allowbreak{}pelota\_\allowbreak{}rumbo\_\allowbreak{}de\_\allowbreak{}regreso} & -152 & MEDIR & rumbo de la línea en la segunda marca \\
\bottomrule\end{longtable}
\subsection{Nivel 4 (Zona Verde)}
Misión 9 (termina de subir la escalera), 10 (rodear la torre), 11 (leer las tarjetas de colores) y 13 (meta azul). La misión 12 (pelota en la casilla roja) está omitida.

\paragraph{Mediciones que necesita el robot}
\begin{longtable}{@{}p{0.36\linewidth}rlp{0.42\linewidth}@{}}\toprule
\textbf{Variable} & \textbf{Valor} & \textbf{Estado} & \textbf{Descripci\'on} \\\midrule\endhead
\texttt{nivel\_\allowbreak{}4\_\allowbreak{}escalera\_\allowbreak{}a\_\allowbreak{}dentro\_\allowbreak{}cm} & 8.0 & MEDIR & avance desde que las sondas ven la zona verde (en la cima de la escalera); entre escalera y torre hay \textasciitilde{}20 cm \\
\texttt{nivel\_\allowbreak{}4\_\allowbreak{}torre\_\allowbreak{}lado\_\allowbreak{}cm} & 10.0 & MEDIR menú & lado de la torre \\
\texttt{nivel\_\allowbreak{}4\_\allowbreak{}zona\_\allowbreak{}verde\_\allowbreak{}a\_\allowbreak{}torre\_\allowbreak{}cm} & 7.5 & MEDIR menú & del borde de salida de la zona verde a la cara de la torre \\
\texttt{nivel\_\allowbreak{}4\_\allowbreak{}pasar\_\allowbreak{}torre\_\allowbreak{}al\_\allowbreak{}oeste\_\allowbreak{}cm} & 37.0 & MEDIR & avance al oeste, paralelo a la línea, hasta pasar la torre (vuelve a la línea cerca de la 1a tarjeta) \\
\texttt{nivel\_\allowbreak{}4\_\allowbreak{}retroceso\_\allowbreak{}para\_\allowbreak{}tarjetas\_\allowbreak{}cm} & 7.0 & MEDIR & la tarjeta más cercana a la torre queda detrás de las sondas; la cola queda a \textasciitilde{}1.5 cm de la torre \\
\bottomrule\end{longtable}
\paragraph{Ajustes de la maniobra}
\begin{longtable}{@{}p{0.36\linewidth}rlp{0.42\linewidth}@{}}\toprule
\textbf{Variable} & \textbf{Valor} & \textbf{Estado} & \textbf{Descripci\'on} \\\midrule\endhead
\texttt{nivel\_\allowbreak{}4\_\allowbreak{}holgura\_\allowbreak{}torre\_\allowbreak{}cm} & 2.0 & menú & holgura al rodear (con 4 cm volvía sobre una tarjeta) \\
\texttt{nivel\_\allowbreak{}4\_\allowbreak{}rodeo\_\allowbreak{}rumbo} & 110 &  & sur-suroeste: lejos de la orilla del nivel 4 \\
\texttt{nivel\_\allowbreak{}4\_\allowbreak{}rodeo\_\allowbreak{}cm} & 20.0 &  & al girar al oeste queda a \textasciitilde{}2 cm de la esquina de la torre (con 17 la tocaba) \\
\bottomrule\end{longtable}
\paragraph{Pelota en la barra (misión 12): OMITIDA por ahora}
\begin{longtable}{@{}p{0.36\linewidth}rlp{0.42\linewidth}@{}}\toprule
\textbf{Variable} & \textbf{Valor} & \textbf{Estado} & \textbf{Descripci\'on} \\\midrule\endhead
\texttt{nivel\_\allowbreak{}4\_\allowbreak{}barra\_\allowbreak{}lado} & 1 & menú & -1 = barra de pelotas a la izquierda de la línea, 1 = derecha \\
\texttt{nivel\_\allowbreak{}4\_\allowbreak{}barra\_\allowbreak{}avance\_\allowbreak{}cm} & 8.2 & MEDIR menú & avance desde la línea para dejar la garra sobre el agujero \\
\bottomrule\end{longtable}
\section{Par\'ametros ajustables (\texttt{cfg.P})}
Valores de \texttt{config.py}. Se pueden cambiar desde el men\'u de ajustes del CyberPi (los que est\'an en \texttt{cfg.TUNE}) y se guardan en el robot; lo guardado manda sobre el archivo. Marcas: \emph{MEDIR}/\emph{MEDIDO} (del robot real), \emph{CALIBRAR}/\emph{CALIBRADO} (con las pruebas), \emph{TODO} (pendiente).
\subsection*{Geometría del robot (MEDIR, hoja en docs/medidas\_robot.md).}
TODO(medidas): los valores marcados TODO son ESTIMADOS a ojo de la foto del diseño propuesto (docs/img/medidas\_propuestas.webp), no medidos. Reemplazarlos con la hoja. Las distancias "\_fwd" se miden hacia adelante desde el PUNTO DE REFERENCIA (punto naranja sobre la caja del mBot2). El centro de giro real de un robot de orugas es el punto medio de las orugas a lo largo; ref\_rot\_fwd dice cuánto está adelante de la referencia.

\begin{longtable}{@{}p{0.2\linewidth}p{0.17\linewidth}lp{0.47\linewidth}@{}}\toprule
\textbf{Clave} & \textbf{Valor} & \textbf{Men\'u} & \textbf{Descripci\'on} \\\midrule\endhead
\texttt{ref\_\allowbreak{}rot\_\allowbreak{}fwd} & \texttt{4.5} & Geometria & TODO(medir) referencia -> punto medio de las orugas (+ = adelante) \\
\texttt{track\_\allowbreak{}c} & \texttt{12.1} & Geometria & MEDIDO (C): eje rueda blanca -> eje 1a negra (tramo apoyado en el piso) \\
\texttt{track\_\allowbreak{}d} & \texttt{4.85} &  & MEDIDO (D): eje 1a negra -> eje 2a negra (nariz levantada) \\
\texttt{nose\_\allowbreak{}deg} & \texttt{45.0} &  & TODO(medir) inclinación de D sobre el piso (estimada de la foto) \\
\texttt{wheel\_\allowbreak{}d} & \texttt{3.95} &  & MEDIDO: diámetro de las ruedas (blanca y negras miden igual) \\
\texttt{belt\_\allowbreak{}t} & \texttt{0.4} &  & TODO(medir) grosor de la banda (de la foto) \\
\texttt{belt\_\allowbreak{}w} & \texttt{3.0} & Geometria & MEDIDO (B, vista superior): ancho de la banda de cada oruga \\
\texttt{robot\_\allowbreak{}w} & \texttt{21.35} & Geometria & MEDIDO: A 15.35 (hueco entre orugas) + 2*B; borde exterior a borde exterior \\
\texttt{sensor\_\allowbreak{}fwd} & \texttt{9.7} & Geometria & = ref\_rot\_fwd (4.5) + 5.2; si se mide ref\_rot\_fwd, ajustar \\
\texttt{probe\_\allowbreak{}pitch} & \texttt{1.25} &  & MEDIDO: entre las dos sondas centrales del Quad RGB (L1-R1) \\
\texttt{probe\_\allowbreak{}outer} & \texttt{2.4} &  & MEDIDO: de cada sonda central a la de afuera (L1-L2, R1-R2) \\
\texttt{card\_\allowbreak{}probe\_\allowbreak{}fwd} & \texttt{9.7} & Geometria & la misma barra del Quad RGB (ver sensor\_fwd) \\
\texttt{sprocket\_\allowbreak{}d} & \texttt{4.18} & Geometria & CALIBRADO: diámetro efectivo (promedio de las dos orugas) \\
\texttt{gyro\_\allowbreak{}scale} & \texttt{1.0} & Geometria & TODO(calibrar) grados reales por grado leído \\
\texttt{lr\_\allowbreak{}ratio} & \texttt{0.977} & Geometria & CALIBRADO: cm por grado izquierda / derecha (63.5 / 65) \\
\texttt{skid} & \texttt{1.3} &  & TODO(calibrar) derrape al girar: B efectivo = distancia entre orugas * skid (girar 360\textdegree{} en el lugar: grados de encoder vs giroscopio) \\
\texttt{slip\_\allowbreak{}tol} & \texttt{0.05} &  &  \\
\texttt{track\_\allowbreak{}max} & \texttt{15.0} &  & grados máximos de corrección al volver a la recta (con 30 la esquina delantera tocaba el muro del pasillo al salir) \\
\texttt{track\_\allowbreak{}k} & \texttt{4.0} &  & grados de corrección de rumbo por cm de desvío lateral (odometry.track)        \# grados por ciclo de diferencia giroscopio/encoders para sospechar patinaje \\
\texttt{slip} & \texttt{1.0} & Geometria & TODO(calibrar con "Probar recta 50") distancia real / teórica \\
\bottomrule\end{longtable}
\subsection*{Brazo (fotos con el servo a 0, 45 y 90 grados). NO es un paralelogramo: la garra gira}
al subir (a 90\textdegree{} apunta hacia arriba, encima de la nariz). Se describe con una TABLA: para cada ángulo del servo, altura del centro de la garra (entre los dedos) sobre el piso y cuánto va DELANTE de la nariz de las orugas (+ adelante). Entre ángulos se interpola. TODO(medir): las alturas y distancias son estimadas de las fotos; medir con regla.

\begin{longtable}{@{}p{0.2\linewidth}p{0.17\linewidth}lp{0.47\linewidth}@{}}\toprule
\textbf{Clave} & \textbf{Valor} & \textbf{Men\'u} & \textbf{Descripci\'on} \\\midrule\endhead
\texttt{arm\_\allowbreak{}h0} & \texttt{8.5} & Brazo & TODO(medir) servo 0\textdegree{}: brazo al frente, \textasciitilde{}10\textdegree{} bajo la horizontal \\
\texttt{arm\_\allowbreak{}x0} & \texttt{8.0} & Brazo & TODO(medir) \\
\texttt{arm\_\allowbreak{}h45} & \texttt{14.5} & Brazo & TODO(medir) servo 45\textdegree{}: garra un poco inclinada hacia arriba \\
\texttt{arm\_\allowbreak{}x45} & \texttt{6.0} & Brazo & TODO(medir) \\
\texttt{arm\_\allowbreak{}h90} & \texttt{23.0} & Brazo & TODO(medir) servo 90\textdegree{}: garra apuntando hacia arriba, sobre la nariz \\
\texttt{arm\_\allowbreak{}x90} & \texttt{0.0} & Brazo & TODO(medir) \\
\texttt{arm\_\allowbreak{}pivot\_\allowbreak{}fwd} & \texttt{10.0} &  & TODO(medir) referencia -> eje fijo de arriba del brazo (solo para dibujar) \\
\texttt{arm\_\allowbreak{}pivot\_\allowbreak{}h} & \texttt{12.2} &  & MEDIDO (G): piso -> eje fijo de arriba del brazo \\
\texttt{arm\_\allowbreak{}ground} & \texttt{5.0} &  & TODO(medir) eje fijo de arriba -> eje del servo (hacia abajo) \\
\texttt{arm\_\allowbreak{}ground\_\allowbreak{}dx} & \texttt{0.7} &  & TODO(medir) el eje del servo va así de adelante del eje fijo \\
\texttt{arm\_\allowbreak{}crank} & \texttt{7.2} &  & TODO(medir) barra de abajo: eje del servo -> articulación del acoplador \\
\texttt{arm\_\allowbreak{}rocker} & \texttt{5.75} &  & TODO(medir) barra de arriba: eje fijo -> articulación del acoplador \\
\texttt{arm\_\allowbreak{}coupler} & \texttt{7.0} &  & TODO(medir) acoplador: articulación de abajo -> de arriba \\
\texttt{arm\_\allowbreak{}claw\_\allowbreak{}u} & \texttt{1.9} &  & TODO(medir) centro de la garra: cm sobre la articulación de abajo \\
\texttt{arm\_\allowbreak{}claw\_\allowbreak{}v} & \texttt{7.9} &  & TODO(medir) centro de la garra: cm delante del acoplador \\
\texttt{claw\_\allowbreak{}tip} & \texttt{3.0} &  & TODO(medir) del centro de la garra a la punta de los dedos \\
\texttt{h\_\allowbreak{}floor} & \texttt{8.5} & Brazo & lo más bajo que llega (servo 0\textdegree{}); el brazo no alcanza el piso \\
\texttt{h\_\allowbreak{}pedestal} & \texttt{12.5} & Brazo & sobre el pedestal (10 cm + medio cubo) \\
\texttt{h\_\allowbreak{}carry} & \texttt{17.5} & Brazo & transporte: base del cubo >= 5 cm sobre el pedestal (10 + 5 + 2.5) \\
\texttt{h\_\allowbreak{}bar} & \texttt{12.0} & Brazo & barra de pelotas (10 cm) \\
\texttt{h\_\allowbreak{}prism} & \texttt{8.5} & Brazo & lo más baja (servo 0\textdegree{}), abierta: lleva el prisma como en una cuna (misión 7) \\
\bottomrule\end{longtable}
\subsection*{Garra (CALIBRAR)}
\begin{longtable}{@{}p{0.2\linewidth}p{0.17\linewidth}lp{0.47\linewidth}@{}}\toprule
\textbf{Clave} & \textbf{Valor} & \textbf{Men\'u} & \textbf{Descripci\'on} \\\midrule\endhead
\texttt{claw\_\allowbreak{}open} & \texttt{90} & Garra & TODO(calibrar) \\
\texttt{claw\_\allowbreak{}cube} & \texttt{35} & Garra & TODO(calibrar) \\
\texttt{claw\_\allowbreak{}ball} & \texttt{45} & Garra & TODO(calibrar) \\
\texttt{com\_\allowbreak{}fwd} & \texttt{-1.0} &  & TODO(medir) centro de masa del chasis, adelante del centro de giro \\
\texttt{com\_\allowbreak{}h} & \texttt{6.5} &  & TODO(medir) altura del centro de masa del chasis sobre el piso \\
\texttt{arm\_\allowbreak{}mass} & \texttt{0.15} &  & TODO(medir) fracción del peso que está en brazo + garra \\
\texttt{arm\_\allowbreak{}w} & \texttt{4.9} &  & MEDIDO (F, vista superior): ancho del marco del brazo \\
\texttt{arm\_\allowbreak{}lat} & \texttt{-0.75} &  & MEDIDO (E = 5.97 de la oruga izquierda al marco): centro del marco 0.75 a la DERECHA del centro (+ = izquierda); (15.35-4.9)/2 - 5.97 \\
\texttt{claw\_\allowbreak{}open\_\allowbreak{}w} & \texttt{9.0} &  & TODO(medir) separación interior de los dedos con la garra abierta (el prisma mide \textasciitilde{}7 cm de base en el plano; pregunta abierta) \\
\bottomrule\end{longtable}
\subsection*{Línea}
\begin{longtable}{@{}p{0.2\linewidth}p{0.17\linewidth}lp{0.47\linewidth}@{}}\toprule
\textbf{Clave} & \textbf{Valor} & \textbf{Men\'u} & \textbf{Descripci\'on} \\\midrule\endhead
\texttt{line\_\allowbreak{}mode} & \texttt{1} & Linea & 1 = analógico calibrado (recomendado), 0 = get\_offset\_track \\
\texttt{dark\_\allowbreak{}th} & \texttt{0.45} & Linea & 0..1 desde qué oscuridad una sonda "ve" línea \\
\texttt{cross\_\allowbreak{}win} & \texttt{4.0} & Linea & cm de avance en los que cuentan las dos sondas exteriores \\
\texttt{cross\_\allowbreak{}th} & \texttt{0.8} & Linea & 0..1 oscuridad de una exterior para contar un cruce (> verde) \\
\texttt{kp} & \texttt{0.45} & Linea &  \\
\texttt{kd} & \texttt{1.5} & Linea &  \\
\texttt{ki} & \texttt{0.0} &  &  \\
\texttt{lost\_\allowbreak{}s} & \texttt{0.6} &  & s sin línea antes de darla por perdida \\
\texttt{green} & \texttt{0} & Linea & 1 = girar con marcas verdes (opcional; la ruta no depende de ellas) \\
\bottomrule\end{longtable}
\subsection*{Velocidades (\% potencia)}
\begin{longtable}{@{}p{0.2\linewidth}p{0.17\linewidth}lp{0.47\linewidth}@{}}\toprule
\textbf{Clave} & \textbf{Valor} & \textbf{Men\'u} & \textbf{Descripci\'on} \\\midrule\endhead
\texttt{v\_\allowbreak{}line} & \texttt{45} & Linea &  \\
\texttt{v\_\allowbreak{}slow} & \texttt{25} & Movimiento &  \\
\texttt{v\_\allowbreak{}straight} & \texttt{45} &  &  \\
\texttt{v\_\allowbreak{}turn\_\allowbreak{}max} & \texttt{45} &  &  \\
\texttt{v\_\allowbreak{}turn\_\allowbreak{}min} & \texttt{18} & Movimiento &  \\
\texttt{v\_\allowbreak{}ramp} & \texttt{70} & Movimiento &  \\
\texttt{v\_\allowbreak{}stairs} & \texttt{100} & Movimiento &  \\
\texttt{v\_\allowbreak{}push} & \texttt{30} &  &  \\
\bottomrule\end{longtable}
\subsection*{Giros y rampas}
\begin{longtable}{@{}p{0.2\linewidth}p{0.17\linewidth}lp{0.47\linewidth}@{}}\toprule
\textbf{Clave} & \textbf{Valor} & \textbf{Men\'u} & \textbf{Descripci\'on} \\\midrule\endhead
\texttt{turn\_\allowbreak{}kp} & \texttt{1.0} & Movimiento &  \\
\texttt{turn\_\allowbreak{}tol} & \texttt{2.5} &  &  \\
\texttt{head\_\allowbreak{}kp} & \texttt{1.2} & Movimiento &  \\
\texttt{pitch\_\allowbreak{}ramp} & \texttt{10} &  & grados: ya va subiendo \\
\texttt{pitch\_\allowbreak{}flat} & \texttt{4} &  & grados: ya está plano arriba \\
\texttt{ramp\_\allowbreak{}margin} & \texttt{3.0} &  & cm extra para quedar "completamente dentro" \\
\texttt{tip\_\allowbreak{}pitch} & \texttt{8} &  & grados: desde aquí la potencia solo sube de a poco \\
\texttt{tip\_\allowbreak{}slew} & \texttt{120} & Movimiento & \%/s: cuánto puede subir la potencia por segundo inclinado \\
\texttt{slope\_\allowbreak{}min\_\allowbreak{}v} & \texttt{4.0} &  & cm/s: subiendo más lento que esto (o resbalando) -> asistencia \\
\texttt{slope\_\allowbreak{}boost} & \texttt{60} &  & \%/s: cuánto sube la potencia extra mientras haga falta (máx. +60) \\
\texttt{h\_\allowbreak{}top} & \texttt{23.0} &  & servo a 90\textdegree{}: garra arriba, encima de la nariz (rodear la torre) \\
\texttt{h\_\allowbreak{}climb} & \texttt{10.5} & Brazo & garra baja y al frente al subir (adelanta y baja el centro de masa). No a 8.5 (servo 0\textdegree{}): el brazo baja al detectar la rampa (\textasciitilde{}10\textdegree{} de pitch, cola aún en el piso) y el cubo pasaba a 1.8 cm de la rampa (0.4 cm con una rampa 20 \% más empinada). Con 10.5: 4.0 / 2.8 cm. \\
\bottomrule\end{longtable}
\subsection*{Medidas de la PISTA que se pueden ajustar desde el menú (grupo "Pista").}
Sus valores se cambian en software/core/pista.py (ahí están explicadas, por nivel).

\begin{longtable}{@{}p{0.2\linewidth}p{0.17\linewidth}lp{0.47\linewidth}@{}}\toprule
\textbf{Clave} & \textbf{Valor} & \textbf{Men\'u} & \textbf{Descripci\'on} \\\midrule\endhead
\texttt{ped1\_\allowbreak{}from\_\allowbreak{}cross} & \texttt{pista.nivel\_\allowbreak{}1\_\allowbreak{}cruce\_\allowbreak{}a\_\allowbreak{}pedestal\_\allowbreak{}cm} & Pista &  \\
\texttt{ball\_\allowbreak{}ped\_\allowbreak{}dist} & \texttt{pista.nivel\_\allowbreak{}3\_\allowbreak{}pelota\_\allowbreak{}linea\_\allowbreak{}a\_\allowbreak{}pedestal\_\allowbreak{}cm} & Pista &  \\
\texttt{tower\_\allowbreak{}from\_\allowbreak{}edge} & \texttt{pista.nivel\_\allowbreak{}4\_\allowbreak{}zona\_\allowbreak{}verde\_\allowbreak{}a\_\allowbreak{}torre\_\allowbreak{}cm} & Pista &  \\
\texttt{tower\_\allowbreak{}size} & \texttt{pista.nivel\_\allowbreak{}4\_\allowbreak{}torre\_\allowbreak{}lado\_\allowbreak{}cm} &  &  \\
\texttt{side\_\allowbreak{}margin} & \texttt{pista.nivel\_\allowbreak{}4\_\allowbreak{}holgura\_\allowbreak{}torre\_\allowbreak{}cm} & Pista &  \\
\texttt{bar\_\allowbreak{}side} & \texttt{pista.nivel\_\allowbreak{}4\_\allowbreak{}barra\_\allowbreak{}lado} & Pista &  \\
\texttt{bar\_\allowbreak{}reach} & \texttt{pista.nivel\_\allowbreak{}4\_\allowbreak{}barra\_\allowbreak{}avance\_\allowbreak{}cm} & Pista &  \\
\bottomrule\end{longtable}
\subsection*{Colores}
\begin{longtable}{@{}p{0.2\linewidth}p{0.17\linewidth}lp{0.47\linewidth}@{}}\toprule
\textbf{Clave} & \textbf{Valor} & \textbf{Men\'u} & \textbf{Descripci\'on} \\\midrule\endhead
\texttt{color\_\allowbreak{}mode} & \texttt{0} &  & 0 = nombre del sensor (get\_color\_sta), 1 = RGB más cercano \\
\bottomrule\end{longtable}
\section{Decisiones de dise\~no}
\medskip\noindent\textbf{SW-DEC-001 -- Dos programas, cada uno en un solo archivo}\par
\noindent\textbf{Decisi\'on:} El autónomo (rover\_x\_auto.py) y el manual (rover\_x\_manual.py) son programas separados. El código se escribe en módulos (software/) y tools/scripts/bundle\_for\_mblock.py los junta en un solo archivo por programa (mblock/).\par
\noindent\textbf{Justificaci\'on:} mBlock sube un solo archivo de Python al CyberPi. Los módulos hacen el código legible y se prueban por separado en la PC; el archivo empaquetado solo importa módulos de Makeblock y de MicroPython (lo verifica una prueba).\par
\medskip\noindent\textbf{SW-DEC-002 -- Capa de hardware única (hal.py)}\par
\noindent\textbf{Decisi\'on:} Solo hal.py llama a la API del mBot2 (cyberpi, mbot2, mbuild, gamepad), con convenciones fijas de signos.\par
\noindent\textbf{Justificaci\'on:} Un signo mal montado (motor, giroscopio, sensor de línea) se corrige en un solo lugar (config.py) y en la PC se reemplaza el hardware por módulos falsos: el mismo programa corre en las pruebas y en el simulador.\par
\medskip\noindent\textbf{SW-DEC-003 -- La ruta como datos}\par
\noindent\textbf{Decisi\'on:} Las misiones son listas de pasos (tuplas) en missions.py; autonomous.run\_step() sabe ejecutar cada tipo de paso.\par
\noindent\textbf{Justificaci\'on:} Cambiar la ruta no requiere tocar la lógica de control. Los argumentos pueden ser nombres de parámetros de cfg.P, que se leen al ejecutar: lo ajustado en el menú del CyberPi se aplica sin recargar el programa.\par
\medskip\noindent\textbf{SW-DEC-004 -- Medidas de la pista en un solo archivo}\par
\noindent\textbf{Decisi\'on:} Todas las distancias, puntos del mapa y rumbos que dependen de la pista están en pista.py, por nivel y con nombres nivel\_N\_...\par
\noindent\textbf{Justificaci\'on:} Son lo primero que hay que medir y ajustar en la pista real. Cada medida está marcada MEDIR o MEDIDO para saber qué falta.\par
\medskip\noindent\textbf{SW-DEC-005 -- Autónomo por tramos}\par
\noindent\textbf{Decisi\'on:} En la pantalla de inicio se elige el nivel con el joystick y B corre solo el tramo de ese nivel; el robot se detiene al terminarlo (o al presionar A).\par
\noindent\textbf{Justificaci\'on:} Así lo usa el equipo en competencia: tras una falla se pide reinicio en la zona del nivel y se repite solo ese tramo. Las misiones de manipulación ya logradas se saltan.\par
\medskip\noindent\textbf{SW-DEC-006 -- Sin ultrasónico: distancias desde referencias del piso}\par
\noindent\textbf{Decisi\'on:} Los acercamientos a pedestales y la torre se miden con los encoders desde un cruce de línea o el borde de una zona de color.\par
\noindent\textbf{Justificaci\'on:} El diseño del robot no lleva ultrasónico. Los cruces y bordes son referencias fijas que reinician el error acumulado.\par
\medskip\noindent\textbf{SW-DEC-007 -- Odometría robusta al patinaje}\par
\noindent\textbf{Decisi\'on:} odometry.py integra encoders y giroscopio; cuando el giroscopio y los encoders no coinciden, mide el avance con la oruga que no patinó.\par
\noindent\textbf{Justificaci\'on:} En la zona de obstáculos los recortes de cartón hacen patinar una oruga: el robot avanza 8 a 28 cm menos de lo que cuentan los encoders (simulador).\par
\medskip\noindent\textbf{SW-DEC-008 -- Cruces con memoria}\par
\noindent\textbf{Decisi\'on:} Un cruce se arma con lo que vean las sondas exteriores dentro de una ventana de avance (cross\_win cm), no con las dos a la vez.\par
\noindent\textbf{Justificaci\'on:} Con una línea de 10 a 30 mm y el robot oscilando, las dos sondas exteriores casi nunca ven negro al mismo tiempo.\par
\medskip\noindent\textbf{SW-DEC-009 -- Control contra volcadura}\par
\noindent\textbf{Decisi\'on:} Inclinado, hal.drive solo deja subir la potencia de a poco (tip\_slew) y suma asistencia si el robot no avanza; el brazo baja al frente al subir.\par
\noindent\textbf{Justificaci\'on:} Un arranque brusco en rampa o escalera pasa el centro de masa detrás de las orugas y el robot se voltea hacia atrás.\par
\section{Referencia de los m\'odulos}
Cada funci\'on con su descripci\'on Doxygen, par\'ametros, lo que devuelve, las funciones del proyecto que llama y su c\'odigo.
\subsection*{Programas para el robot}
\subsection{\texttt{software/robot/rover\_\allowbreak{}x\_\allowbreak{}auto.py}}\label{mod:rover-x-auto}
\textbf{Programa AUTÓNOMO del robot (se carga en el CyberPi como mblock/rover\_x\_auto.py).}

Botones del CyberPi:

{\small\begin{alltt}
joystick   arriba = nivel 1, abajo = 2, izquierda = 3, derecha = 4 (cfg.LEVEL\_KEYS)
B          empezar: corre el tramo de ese nivel y se detiene al terminarlo (cfg.START\_BUTTON)
A          parar el robot durante el tramo (cfg.ABORT\_BUTTON)
centro     en la pantalla de inicio: menú de ajustes (cfg.MENU\_BUTTON)
\end{alltt}}
Al terminar o parar vuelve a la pantalla de inicio con el nivel que sigue ya elegido.

El control manual es otro programa: rover\_x\_manual.py. Archivo listo para mBlock: mblock/rover\_x\_auto.py (se regenera con: python3 tools/scripts/bundle\_for\_mblock.py --all).

\noindent Archivo fuente: \texttt{software/robot/rover\_\allowbreak{}x\_\allowbreak{}auto.py} (52 l\'ineas).

\subsubsection{\texttt{main()}}\label{fn:rover_x_auto:main}
\noindent Programa autónomo: pantalla de inicio -> tramo del nivel elegido -> otra vez la pantalla.

Carga los parámetros guardados, inicializa el hardware y pone la garra en posición de viaje. En un ciclo: elige el nivel (choose\_level), corre el tramo (autonomous.run) y, al terminar o al parar con A, vuelve a la pantalla de inicio con el siguiente nivel ya elegido. El centro del joystick abre el menú de ajustes.

\noindent\textbf{Llama a:} \hyperref[fn:settings:remember-defaults]{\texttt{settings.remember\_\allowbreak{}defaults}}, \hyperref[fn:settings:load]{\texttt{settings.load}}, \hyperref[fn:hal:init]{\texttt{hal.init}}, \hyperref[fn:hal:show]{\texttt{hal.show}}, \hyperref[fn:gripper:home]{\texttt{gripper.home}}, \hyperref[fn:autonomous:choose-level]{\texttt{autonomous.choose\_\allowbreak{}level}}, \hyperref[fn:hal:clear-abort]{\texttt{hal.clear\_\allowbreak{}abort}}, \hyperref[fn:autonomous:run]{\texttt{autonomous.run}}, \hyperref[fn:hal:stop]{\texttt{hal.stop}}, \hyperref[fn:hal:abort-requested]{\texttt{hal.abort\_\allowbreak{}requested}}, \hyperref[fn:hal:wait]{\texttt{hal.wait}}, \hyperref[fn:settings:menu]{\texttt{settings.menu}}\par
\lstinputlisting[firstline=26,lastline=49,firstnumber=26]{../../software/robot/rover_x_auto.py}
\subsection{\texttt{software/robot/rover\_\allowbreak{}x\_\allowbreak{}manual.py}}\label{mod:rover-x-manual}
\textbf{Programa MANUAL del robot (se carga en el CyberPi como mblock/rover\_x\_manual.py).}

Botones del CyberPi:

{\small\begin{alltt}
B   empezar el control manual (cfg.START\_BUTTON)
A   parar: detiene el robot y vuelve a la espera (cfg.ABORT\_BUTTON)
\end{alltt}}
Los mandos del gamepad están en docs/control\_manual.md.

El autónomo es otro programa: rover\_x\_auto.py. Archivo listo para mBlock: mblock/rover\_x\_manual.py (se regenera con: python3 tools/scripts/bundle\_for\_mblock.py --all).

\noindent Archivo fuente: \texttt{software/robot/rover\_\allowbreak{}x\_\allowbreak{}manual.py} (50 l\'ineas).

\subsubsection{\texttt{wait\_\allowbreak{}start()}}\label{fn:rover_x_manual:wait-start}
\noindent Pantalla de espera: muestra los botones y espera B.

Espera también a que se suelte B, para que no cuente como otra pulsación. Muestra si el gamepad está conectado.

\noindent\textbf{Llama a:} \hyperref[fn:hal:show]{\texttt{hal.show}}, \hyperref[fn:hal:button]{\texttt{hal.button}}, \hyperref[fn:hal:wait]{\texttt{hal.wait}}, \hyperref[fn:hal:pad-ok]{\texttt{hal.pad\_\allowbreak{}ok}}\par
\lstinputlisting[firstline=22,lastline=30,firstnumber=22]{../../software/robot/rover_x_manual.py}
\subsubsection{\texttt{main()}}\label{fn:rover_x_manual:main}
\noindent Programa manual: espera B -> control con el gamepad -> A para -> otra vez la espera.

Carga los parámetros guardados, inicializa el hardware y pone la garra en posición de viaje antes del primer arranque.

\noindent\textbf{Llama a:} \hyperref[fn:settings:remember-defaults]{\texttt{settings.remember\_\allowbreak{}defaults}}, \hyperref[fn:settings:load]{\texttt{settings.load}}, \hyperref[fn:hal:init]{\texttt{hal.init}}, \hyperref[fn:gripper:home]{\texttt{gripper.home}}, \hyperref[fn:rover_x_manual:wait-start]{\texttt{wait\_\allowbreak{}start}}, \hyperref[fn:manual:loop]{\texttt{manual.loop}}, \hyperref[fn:hal:stop]{\texttt{hal.stop}}, \hyperref[fn:hal:button]{\texttt{hal.button}}, \hyperref[fn:hal:clear-abort]{\texttt{hal.clear\_\allowbreak{}abort}}, \hyperref[fn:hal:wait]{\texttt{hal.wait}}\par
\lstinputlisting[firstline=36,lastline=47,firstnumber=36]{../../software/robot/rover_x_manual.py}
\subsection*{Autónomo}
\subsection{\texttt{software/autonomous/autonomous.py}}\label{mod:autonomous}
\textbf{Ejecutor del área automática: corre las misiones de missions.py.}

Ejecuta las misiones desde el nivel de inicio elegido (un TRAMO), muestra el avance en la pantalla del CyberPi y se detiene al terminar el tramo, si se acaba el tiempo, si falla una misión crítica o si se presiona A (cfg.ABORT\_BUTTON: parar para pedir reinicio). Cada misión es una lista de PASOS (tuplas); run\_step() ejecuta uno. La lista de pasos y su formato están en missions.py. También tiene la pantalla de inicio (choose\_level) y las misiones de tarjetas de colores (scan\_cards, drop\_red) y meta (stop\_in).

\noindent Archivo fuente: \texttt{software/autonomous/autonomous.py} (734 l\'ineas).

\subsubsection{\texttt{reset\_\allowbreak{}state()}}\label{fn:autonomous:reset-state}
\noindent Borra todo lo que una corrida anterior pudo dejar.

Tarjetas, resultados, último acercamiento, referencia de rumbo y estado del seguidor. Al relanzar el autónomo sin apagar el robot, una referencia vieja movía el primer giro (simulador).

\noindent\textbf{Llama a:} \hyperref[fn:linefollow:reset]{\texttt{linefollow.reset}}\par
\lstinputlisting[firstline=38,lastline=48,firstnumber=38]{../../software/autonomous/autonomous.py}
\subsubsection{\texttt{val(x)}}\label{fn:autonomous:val}
\noindent Valor de un argumento de paso.

En las misiones un valor puede ser un número o el NOMBRE de un parámetro de cfg.P; el nombre se lee al ejecutar, así lo que se cambie en el menú se aplica sin recargar el programa.

\par\noindent{\small\begin{tabular}{@{}p{0.2\linewidth}p{0.74\linewidth}@{}}
\textbf{Par\'ametro} & \textbf{Descripci\'on} \\
\texttt{x} & Número o nombre de parámetro (str). \\
\end{tabular}}\par\smallskip
\noindent\textbf{Devuelve:} El número, o cfg.P[x].\par
\lstinputlisting[firstline=56,lastline=59,firstnumber=56]{../../software/autonomous/autonomous.py}
\subsubsection{\texttt{make\_\allowbreak{}cond(spec)}}\label{fn:autonomous:make-cond}
\noindent Convierte la condición de un paso "line" en una función de parada.

("junction", n, tipo), ("color", nombre), ("color\_cm", color, cm), ("cm", d), ("leave", color), ("obstacle", cm), ("end",), ("pitch",), ("head", rumbo, tolerancia).

\par\noindent{\small\begin{tabular}{@{}p{0.2\linewidth}p{0.74\linewidth}@{}}
\textbf{Par\'ametro} & \textbf{Descripci\'on} \\
\texttt{spec} & Tupla de la condición. \\
\end{tabular}}\par\smallskip
\noindent\textbf{Devuelve:} Función cond(bits, junction) para linefollow.follow\_until().\par
\noindent\textbf{Llama a:} \hyperref[fn:linefollow:at-junction]{\texttt{linefollow.at\_\allowbreak{}junction}}, \hyperref[fn:linefollow:at-color]{\texttt{linefollow.at\_\allowbreak{}color}}, \hyperref[fn:linefollow:color-then-cm]{\texttt{linefollow.color\_\allowbreak{}then\_\allowbreak{}cm}}, \hyperref[fn:linefollow:after-cm]{\texttt{linefollow.after\_\allowbreak{}cm}}, \hyperref[fn:linefollow:leave-color]{\texttt{linefollow.leave\_\allowbreak{}color}}, \hyperref[fn:linefollow:obstacle-within]{\texttt{linefollow.obstacle\_\allowbreak{}within}}, \hyperref[fn:linefollow:line-end]{\texttt{linefollow.line\_\allowbreak{}end}}, \hyperref[fn:linefollow:pitch-above]{\texttt{linefollow.pitch\_\allowbreak{}above}}, \hyperref[fn:autonomous:val]{\texttt{val}}, \hyperref[fn:motion:wrap180]{\texttt{motion.wrap180}}, \hyperref[fn:hal:heading]{\texttt{hal.heading}}\par
\lstinputlisting[firstline=68,lastline=92,firstnumber=68]{../../software/autonomous/autonomous.py}
\subsubsection{\texttt{red\_\allowbreak{}index(found)}}\label{fn:autonomous:red-index}
\noindent Posición de la tarjeta roja en la lista leída.

\par\noindent{\small\begin{tabular}{@{}p{0.2\linewidth}p{0.74\linewidth}@{}}
\textbf{Par\'ametro} & \textbf{Descripci\'on} \\
\texttt{found} & Lista [(color, odometría)]. \\
\end{tabular}}\par\smallskip
\noindent\textbf{Devuelve:} Índice 0..5, o -1 si no hay roja.\par
\lstinputlisting[firstline=99,lastline=103,firstnumber=99]{../../software/autonomous/autonomous.py}
\subsubsection{\texttt{\_\allowbreak{}note\_\allowbreak{}card(c, odo)}}\label{fn:autonomous:-note-card}
\noindent Registra una tarjeta si no es una ya vista.

Mismo color a menos de 6 cm = la misma (la tarjeta mide 4.7 cm y al retroceder y volver a pasar se lee en sus dos orillas). La muestra en los LEDs y la pantalla.

\par\noindent{\small\begin{tabular}{@{}p{0.2\linewidth}p{0.74\linewidth}@{}}
\textbf{Par\'ametro} & \textbf{Descripci\'on} \\
\texttt{c} & Color leído. \\
\texttt{odo} & Odometría absoluta (cm) al leerla. \\
\end{tabular}}\par\smallskip
\noindent\textbf{Devuelve:} True si era nueva.\par
\noindent\textbf{Llama a:} \hyperref[fn:colors:show]{\texttt{colors.show}}\par
\lstinputlisting[firstline=112,lastline=118,firstnumber=112]{../../software/autonomous/autonomous.py}
\subsubsection{\texttt{\_\allowbreak{}scan\_\allowbreak{}back(back\_\allowbreak{}cm)}}\label{fn:autonomous:-scan-back}
\noindent Retrocede recto leyendo tarjetas.

Al volver a la línea después de rodear la torre, la sonda ya pasó la primera tarjeta.

\par\noindent{\small\begin{tabular}{@{}p{0.2\linewidth}p{0.74\linewidth}@{}}
\textbf{Par\'ametro} & \textbf{Descripci\'on} \\
\texttt{back\_\allowbreak{}cm} & cm a retroceder. \\
\end{tabular}}\par\smallskip
\noindent\textbf{Llama a:} \hyperref[fn:hal:odo-abs]{\texttt{hal.odo\_\allowbreak{}abs}}, \hyperref[fn:hal:heading]{\texttt{hal.heading}}, \hyperref[fn:hal:now]{\texttt{hal.now}}, \hyperref[fn:hal:stop]{\texttt{hal.stop}}, \hyperref[fn:hal:abort-requested]{\texttt{hal.abort\_\allowbreak{}requested}}, \hyperref[fn:colors:read]{\texttt{colors.read}}, \hyperref[fn:hal:clamp]{\texttt{hal.clamp}}, \hyperref[fn:hal:drive]{\texttt{hal.drive}}, \hyperref[fn:hal:wait]{\texttt{hal.wait}}, \hyperref[fn:colors:is-card-color]{\texttt{colors.is\_\allowbreak{}card\_\allowbreak{}color}}, \hyperref[fn:autonomous:-note-card]{\texttt{\_\allowbreak{}note\_\allowbreak{}card}}\par
\lstinputlisting[firstline=124,lastline=143,firstnumber=124]{../../software/autonomous/autonomous.py}
\subsubsection{\texttt{scan\_\allowbreak{}cards(back\_\allowbreak{}cm)}}\label{fn:autonomous:scan-cards}
\noindent Misión 11: sigue la línea despacio y registra cada tarjeta de color.

Para al ver la X azul de la meta (azul por más de GOAL\_RUN\_CM seguidos; una tarjeta mide \textasciitilde{}4.7 cm) o al recorrer CARD\_SCAN\_MAX\_CM. El azul de la meta no cuenta como tarjeta.

\par\noindent{\small\begin{tabular}{@{}p{0.2\linewidth}p{0.74\linewidth}@{}}
\textbf{Par\'ametro} & \textbf{Descripci\'on} \\
\texttt{back\_\allowbreak{}cm} & Antes, retrocede eso leyendo (tarjetas que quedaron detrás de la sonda). \\
\end{tabular}}\par\smallskip
\noindent\textbf{Devuelve:} True si leyó al menos una tarjeta.\par
\noindent\textbf{Llama a:} \hyperref[fn:linefollow:reset]{\texttt{linefollow.reset}}, \hyperref[fn:hal:odo-abs]{\texttt{hal.odo\_\allowbreak{}abs}}, \hyperref[fn:hal:now]{\texttt{hal.now}}, \hyperref[fn:hal:stop]{\texttt{hal.stop}}, \hyperref[fn:autonomous:-scan-back]{\texttt{\_\allowbreak{}scan\_\allowbreak{}back}}, \hyperref[fn:motion:sees-color]{\texttt{motion.sees\_\allowbreak{}color}}, \hyperref[fn:hal:abort-requested]{\texttt{hal.abort\_\allowbreak{}requested}}, \hyperref[fn:colors:read]{\texttt{colors.read}}, \hyperref[fn:linefollow:step]{\texttt{linefollow.step}}, \hyperref[fn:hal:wait]{\texttt{hal.wait}}, \hyperref[fn:colors:is-card-color]{\texttt{colors.is\_\allowbreak{}card\_\allowbreak{}color}}, \hyperref[fn:autonomous:-note-card]{\texttt{\_\allowbreak{}note\_\allowbreak{}card}}\par
\lstinputlisting[firstline=151,lastline=193,firstnumber=151]{../../software/autonomous/autonomous.py}
\subsubsection{\texttt{drop\_\allowbreak{}red()}}\label{fn:autonomous:drop-red}
\noindent Misión 12 (omitida por ahora): deja la pelota en la casilla de la tarjeta roja.

Se coloca frente a la tarjeta roja, gira hacia la barra (bar\_side), baja la garra a h\_bar, avanza bar\_reach, suelta y regresa a la línea.

\noindent\textbf{Devuelve:} True si la dejó.\par
\noindent\textbf{Llama a:} \hyperref[fn:autonomous:red-index]{\texttt{red\_\allowbreak{}index}}, \hyperref[fn:motion:straight]{\texttt{motion.straight}}, \hyperref[fn:gripper:arm-to]{\texttt{gripper.arm\_\allowbreak{}to}}, \hyperref[fn:gripper:open-claw]{\texttt{gripper.open\_\allowbreak{}claw}}, \hyperref[fn:motion:turn]{\texttt{motion.turn}}, \hyperref[fn:motion:spin-to-line]{\texttt{motion.spin\_\allowbreak{}to\_\allowbreak{}line}}, \hyperref[fn:geometry:card-to-center]{\texttt{geometry.card\_\allowbreak{}to\_\allowbreak{}center}}, \hyperref[fn:hal:odo-abs]{\texttt{hal.odo\_\allowbreak{}abs}}\par
\lstinputlisting[firstline=200,lastline=214,firstnumber=200]{../../software/autonomous/autonomous.py}
\subsubsection{\texttt{stop\_\allowbreak{}in(color)}}\label{fn:autonomous:stop-in}
\noindent Misión 13: para completamente dentro de una zona (meta).

Avanza zone\_extra() contado desde que las sondas ENTRARON a la zona; si la lectura de tarjetas ya entró, no vuelve a contar desde aquí (antes se pasaba y salía por el otro lado).

\par\noindent{\small\begin{tabular}{@{}p{0.2\linewidth}p{0.74\linewidth}@{}}
\textbf{Par\'ametro} & \textbf{Descripci\'on} \\
\texttt{color} & Color de la zona. \\
\end{tabular}}\par\smallskip
\noindent\textbf{Devuelve:} True si quedó dentro.\par
\noindent\textbf{Llama a:} \hyperref[fn:hal:stop]{\texttt{hal.stop}}, \hyperref[fn:geometry:zone-extra]{\texttt{geometry.zone\_\allowbreak{}extra}}, \hyperref[fn:motion:straight]{\texttt{motion.straight}}, \hyperref[fn:motion:sees-color]{\texttt{motion.sees\_\allowbreak{}color}}, \hyperref[fn:linefollow:follow-until]{\texttt{linefollow.follow\_\allowbreak{}until}}, \hyperref[fn:hal:odo-abs]{\texttt{hal.odo\_\allowbreak{}abs}}\par
\lstinputlisting[firstline=222,lastline=239,firstnumber=222]{../../software/autonomous/autonomous.py}
\subsubsection{\texttt{\_\allowbreak{}advance(cm)}}\label{fn:autonomous:-advance}
\noindent Avance de un acercamiento; lo guarda para que "retreat" regrese lo mismo.

\par\noindent{\small\begin{tabular}{@{}p{0.2\linewidth}p{0.74\linewidth}@{}}
\textbf{Par\'ametro} & \textbf{Descripci\'on} \\
\texttt{cm} & cm a avanzar. \\
\end{tabular}}\par\smallskip
\noindent\textbf{Devuelve:} True si recorrió la distancia.\par
\noindent\textbf{Llama a:} \hyperref[fn:motion:straight]{\texttt{motion.straight}}\par
\lstinputlisting[firstline=245,lastline=248,firstnumber=245]{../../software/autonomous/autonomous.py}
\subsubsection{\texttt{approach\_\allowbreak{}pedestal(h, from\_\allowbreak{}cross)}}\label{fn:autonomous:approach-pedestal}
\noindent Deja la garra sobre el centro de un pedestal medido desde el último cruce.

\par\noindent{\small\begin{tabular}{@{}p{0.2\linewidth}p{0.74\linewidth}@{}}
\textbf{Par\'ametro} & \textbf{Descripci\'on} \\
\texttt{h} & Altura de la garra (número o nombre de parámetro). \\
\texttt{from\_\allowbreak{}cross} & cm del cruce al centro del pedestal (número o nombre de parámetro). \\
\end{tabular}}\par\smallskip
\noindent\textbf{Devuelve:} True si avanzó.\par
\noindent\textbf{Llama a:} \hyperref[fn:autonomous:-advance]{\texttt{\_\allowbreak{}advance}}, \hyperref[fn:geometry:pedestal-advance]{\texttt{geometry.pedestal\_\allowbreak{}advance}}, \hyperref[fn:hal:odo-abs]{\texttt{hal.odo\_\allowbreak{}abs}}, \hyperref[fn:autonomous:val]{\texttt{val}}\par
\lstinputlisting[firstline=255,lastline=259,firstnumber=255]{../../software/autonomous/autonomous.py}
\subsubsection{\texttt{approach\_\allowbreak{}object(h, dist)}}\label{fn:autonomous:approach-object}
\noindent Deja la garra sobre un objeto al frente.

\par\noindent{\small\begin{tabular}{@{}p{0.2\linewidth}p{0.74\linewidth}@{}}
\textbf{Par\'ametro} & \textbf{Descripci\'on} \\
\texttt{h} & Altura de la garra. \\
\texttt{dist} & cm del centro de giro al centro del objeto. \\
\end{tabular}}\par\smallskip
\noindent\textbf{Devuelve:} True si avanzó.\par
\noindent\textbf{Llama a:} \hyperref[fn:autonomous:-advance]{\texttt{\_\allowbreak{}advance}}, \hyperref[fn:geometry:reach-advance]{\texttt{geometry.reach\_\allowbreak{}advance}}, \hyperref[fn:autonomous:val]{\texttt{val}}\par
\lstinputlisting[firstline=266,lastline=267,firstnumber=266]{../../software/autonomous/autonomous.py}
\subsubsection{\texttt{to\_\allowbreak{}tower()}}\label{fn:autonomous:to-tower}
\noindent Sigue la línea hasta quedar a side\_margin de la torre (desde el borde de la zona verde).

\noindent\textbf{Devuelve:} True si llegó.\par
\noindent\textbf{Llama a:} \hyperref[fn:geometry:tower-advance]{\texttt{geometry.tower\_\allowbreak{}advance}}, \hyperref[fn:linefollow:follow-until]{\texttt{linefollow.follow\_\allowbreak{}until}}, \hyperref[fn:motion:straight]{\texttt{motion.straight}}, \hyperref[fn:linefollow:after-cm]{\texttt{linefollow.after\_\allowbreak{}cm}}, \hyperref[fn:hal:odo-abs]{\texttt{hal.odo\_\allowbreak{}abs}}\par
\lstinputlisting[firstline=272,lastline=281,firstnumber=272]{../../software/autonomous/autonomous.py}
\subsubsection{\texttt{map\_\allowbreak{}target(deg)}}\label{fn:autonomous:map-target}
\noindent Rumbo de mapa -> valor de hal.heading() al que hay que llegar por el camino corto.

hal.heading() acumula vueltas; sin normalizar, una recta podía "sostener" un rumbo a 360\textdegree{} del real y torcerse (simulador).

\par\noindent{\small\begin{tabular}{@{}p{0.2\linewidth}p{0.74\linewidth}@{}}
\textbf{Par\'ametro} & \textbf{Descripci\'on} \\
\texttt{deg} & Rumbo de mapa (0 este, 90 sur, 180 oeste, -90 norte). \\
\end{tabular}}\par\smallskip
\noindent\textbf{Devuelve:} Grados en la escala de hal.heading().\par
\noindent\textbf{Llama a:} \hyperref[fn:hal:heading]{\texttt{hal.heading}}, \hyperref[fn:motion:wrap180]{\texttt{motion.wrap180}}\par
\lstinputlisting[firstline=293,lastline=295,firstnumber=293]{../../software/autonomous/autonomous.py}
\subsubsection{\texttt{run\_\allowbreak{}step(step)}}\label{fn:autonomous:run-step}
\noindent Ejecuta UN paso de una misión.

El primer elemento de la tupla dice qué hacer: "line", "jturn", "straight", "center", "turn", "ref", "head", "drive", "pose", "square", "track\_line", "diag", "track", "back\_to", "ramp\_up", "fix\_line", "bias", "fix", "pivot", "goto", "edge", "cross", "seek", "spin", "approach\_ped", "approach\_obj", "retreat", "to\_tower", "climb", "line\_ref", "if\_placed", "if\_skipped", "if\_ran", "if\_start", "crest", "stairs", "claw", "arm", "hold", "detour", "cards", "drop\_red", "stop\_in", "wait". El formato de cada uno está en missions.py.

\par\noindent{\small\begin{tabular}{@{}p{0.2\linewidth}p{0.74\linewidth}@{}}
\textbf{Par\'ametro} & \textbf{Descripci\'on} \\
\texttt{step} & Tupla del paso. \\
\end{tabular}}\par\smallskip
\noindent\textbf{Devuelve:} True si el paso terminó bien.\par
\noindent\textbf{Llama a:} \hyperref[fn:linefollow:follow-until]{\texttt{linefollow.follow\_\allowbreak{}until}}, \hyperref[fn:linefollow:turn-at-junction]{\texttt{linefollow.turn\_\allowbreak{}at\_\allowbreak{}junction}}, \hyperref[fn:motion:straight]{\texttt{motion.straight}}, \hyperref[fn:motion:turn]{\texttt{motion.turn}}, \hyperref[fn:motion:turn-to]{\texttt{motion.turn\_\allowbreak{}to}}, \hyperref[fn:autonomous:map-target]{\texttt{map\_\allowbreak{}target}}, \hyperref[fn:odometry:pose]{\texttt{odometry.pose}}, \hyperref[fn:odometry:set-pose]{\texttt{odometry.set\_\allowbreak{}pose}}, \hyperref[fn:odometry:track]{\texttt{odometry.track}}, \hyperref[fn:hal:stop]{\texttt{hal.stop}}, \hyperref[fn:hal:wait]{\texttt{hal.wait}}, \hyperref[fn:motion:straight-until]{\texttt{motion.straight\_\allowbreak{}until}}, \hyperref[fn:autonomous:val]{\texttt{val}}, \hyperref[fn:motion:pivot]{\texttt{motion.pivot}}, \hyperref[fn:odometry:goto]{\texttt{odometry.goto}}, \hyperref[fn:motion:drive-to-edge]{\texttt{motion.drive\_\allowbreak{}to\_\allowbreak{}edge}}, \hyperref[fn:motion:drive-to-cross]{\texttt{motion.drive\_\allowbreak{}to\_\allowbreak{}cross}}, \hyperref[fn:motion:seek-line]{\texttt{motion.seek\_\allowbreak{}line}}, \hyperref[fn:motion:spin-to-line]{\texttt{motion.spin\_\allowbreak{}to\_\allowbreak{}line}}, \hyperref[fn:autonomous:approach-pedestal]{\texttt{approach\_\allowbreak{}pedestal}}, \hyperref[fn:autonomous:approach-object]{\texttt{approach\_\allowbreak{}object}}, \hyperref[fn:autonomous:to-tower]{\texttt{to\_\allowbreak{}tower}}, \hyperref[fn:motion:climb]{\texttt{motion.climb}}, \hyperref[fn:autonomous:make-cond]{\texttt{make\_\allowbreak{}cond}}, \hyperref[fn:motion:climb-crest]{\texttt{motion.climb\_\allowbreak{}crest}}, \hyperref[fn:motion:stairs]{\texttt{motion.stairs}}, \hyperref[fn:gripper:claw]{\texttt{gripper.claw}}, \hyperref[fn:gripper:arm-to]{\texttt{gripper.arm\_\allowbreak{}to}}, \hyperref[fn:gripper:hold]{\texttt{gripper.hold}}, \hyperref[fn:motion:detour]{\texttt{motion.detour}}, \hyperref[fn:autonomous:scan-cards]{\texttt{scan\_\allowbreak{}cards}}, \hyperref[fn:autonomous:drop-red]{\texttt{drop\_\allowbreak{}red}}, \hyperref[fn:autonomous:stop-in]{\texttt{stop\_\allowbreak{}in}}, \hyperref[fn:geometry:turn-advance]{\texttt{geometry.turn\_\allowbreak{}advance}}, \hyperref[fn:hal:heading]{\texttt{hal.heading}}, \hyperref[fn:hal:line-bits]{\texttt{hal.line\_\allowbreak{}bits}}, \hyperref[fn:hal:line-read]{\texttt{hal.line\_\allowbreak{}read}}, \hyperref[fn:motion:wrap180]{\texttt{motion.wrap180}}, \hyperref[fn:motion:sees-color]{\texttt{motion.sees\_\allowbreak{}color}}, \hyperref[fn:hal:pitch]{\texttt{hal.pitch}}\par
\lstinputlisting[firstline=308,lastline=623,firstnumber=308]{../../software/autonomous/autonomous.py}
\subsubsection{\texttt{run\_\allowbreak{}mission(m)}}\label{fn:autonomous:run-mission}
\noindent Ejecuta todos los pasos de una misión.

\par\noindent{\small\begin{tabular}{@{}p{0.2\linewidth}p{0.74\linewidth}@{}}
\textbf{Par\'ametro} & \textbf{Descripci\'on} \\
\texttt{m} & Dict de la misión (de missions.MISSIONS). \\
\end{tabular}}\par\smallskip
\noindent\textbf{Devuelve:} True si todos los pasos terminaron bien (False al primer paso que falla o si se pidió parar).\par
\noindent\textbf{Llama a:} \hyperref[fn:hal:show]{\texttt{hal.show}}, \hyperref[fn:hal:abort-requested]{\texttt{hal.abort\_\allowbreak{}requested}}, \hyperref[fn:autonomous:run-step]{\texttt{run\_\allowbreak{}step}}\par
\lstinputlisting[firstline=629,lastline=636,firstnumber=629]{../../software/autonomous/autonomous.py}
\subsubsection{\texttt{run(start\_\allowbreak{}level, budget\_\allowbreak{}s, t\_\allowbreak{}start, skip)}}\label{fn:autonomous:run}
\noindent Corre el TRAMO que empieza en start\_level.

Con cfg.SEGMENTS el robot se detiene al terminar el tramo (misión con "end\_level") o al fallar una misión crítica; next\_level queda con el nivel que conviene elegir después.

\par\noindent{\small\begin{tabular}{@{}p{0.2\linewidth}p{0.74\linewidth}@{}}
\textbf{Par\'ametro} & \textbf{Descripci\'on} \\
\texttt{start\_\allowbreak{}level} & Nivel elegido en la pantalla de inicio (1..4). \\
\texttt{budget\_\allowbreak{}s} & Tiempo máximo (por defecto cfg.AUTO\_BUDGET\_S). \\
\texttt{t\_\allowbreak{}start} & Momento de inicio (por defecto ahora). \\
\texttt{skip} & Misiones ya logradas en este encendido: solo se saltan las de manipulación (las de navegación, críticas, se repiten porque llevan al robot a donde empieza la siguiente). \\
\end{tabular}}\par\smallskip
\noindent\textbf{Devuelve:} Puntos estimados del tramo.\par
\noindent\textbf{Llama a:} \hyperref[fn:autonomous:reset-state]{\texttt{reset\_\allowbreak{}state}}, \hyperref[fn:hal:stop]{\texttt{hal.stop}}, \hyperref[fn:hal:now]{\texttt{hal.now}}, \hyperref[fn:autonomous:run-mission]{\texttt{run\_\allowbreak{}mission}}, \hyperref[fn:hal:abort-requested]{\texttt{hal.abort\_\allowbreak{}requested}}, \hyperref[fn:hal:leds]{\texttt{hal.leds}}, \hyperref[fn:hal:show]{\texttt{hal.show}}\par
\lstinputlisting[firstline=648,lastline=698,firstnumber=648]{../../software/autonomous/autonomous.py}
\subsubsection{\texttt{detect\_\allowbreak{}level()}}\label{fn:autonomous:detect-level}
\noindent Nivel según el color de la zona de inicio bajo el robot.

\noindent\textbf{Devuelve:} 1..4 (rojo, amarillo, azul, verde); 1 si no se reconoce.\par
\noindent\textbf{Llama a:} \hyperref[fn:colors:zone-level]{\texttt{colors.zone\_\allowbreak{}level}}, \hyperref[fn:colors:read-stable]{\texttt{colors.read\_\allowbreak{}stable}}\par
\lstinputlisting[firstline=704,lastline=709,firstnumber=704]{../../software/autonomous/autonomous.py}
\subsubsection{\texttt{choose\_\allowbreak{}level(level)}}\label{fn:autonomous:choose-level}
\noindent Pantalla de inicio del autónomo.

Joystick del CyberPi: elige el nivel (cfg.LEVEL\_KEYS: arriba 1, abajo 2, izquierda 3, derecha 4). B (cfg.START\_BUTTON): arrancar el tramo. Centro (cfg.MENU\_BUTTON): menú de ajustes.

\par\noindent{\small\begin{tabular}{@{}p{0.2\linewidth}p{0.74\linewidth}@{}}
\textbf{Par\'ametro} & \textbf{Descripci\'on} \\
\texttt{level} & Nivel sugerido (None = leerlo del color de la zona). \\
\end{tabular}}\par\smallskip
\noindent\textbf{Devuelve:} Nivel elegido 1..4, o 0 = ir al menú de ajustes.\par
\noindent\textbf{Llama a:} \hyperref[fn:hal:wait]{\texttt{hal.wait}}, \hyperref[fn:autonomous:detect-level]{\texttt{detect\_\allowbreak{}level}}, \hyperref[fn:hal:show]{\texttt{hal.show}}, \hyperref[fn:hal:button]{\texttt{hal.button}}\par
\lstinputlisting[firstline=717,lastline=734,firstnumber=717]{../../software/autonomous/autonomous.py}
\subsection{\texttt{software/autonomous/missions.py}}\label{mod:missions}
\textbf{Misiones del área automática, escritas como listas de pasos (la RUTA).}

BORRADOR DE RUTA: se armó con el plano de la arena del reglamento, reconstruido en simulation/arena.py, y se prueba con el simulador (python3 simulation/run\_auto.py). Los conteos de cruces, giros y distancias se ajustan en la pista. La lógica de cada paso está en autonomous.py (run\_step).

Los valores pueden ser números o NOMBRES de parámetros de cfg.P ("v\_line", "h\_pedestal", "claw\_cube"...). Los nombres se leen al ejecutar, así que lo que se cambie en el menú del CyberPi se aplica sin volver a cargar el programa.

Pasos:

{\small\begin{alltt}
("line", cond, vel, timeout[, verde])  seguir la línea hasta cond:
     ("junction", n, tipo)  n cruces; tipo "any" / "X" / "L" / "R"
     ("color", nombre)  ("leave", color) salir de una zona  ("cm", d)  ("end",)
     ("color\_cm", color, d)  d cm después de ver el color (terminar una rampa sin parar en ella)
     ("pitch",) empieza una rampa   ("obstacle", cm) solo si hay ultrasónico
("jturn", grados)          girar EN el último cruce (centro de giro sobre él) y alinear
("straight", cm[, vel])    recto con giroscopio (negativo = atrás)
("center",)                avanzar lo que hay de las sondas al centro de giro (para girar
                           sobre la marca o el cruce que acaban de ver las sondas)
("turn", grados)           giro relativo, positivo = derecha
("ref", rumbo)             "ahora miro hacia este rumbo de mapa" (0 este, 90 sur, 180 oeste)
("line\_ref", rumbo, cond, vel, timeout)  seguir la línea promediando el rumbo -> referencia
("head", rumbo)            girar hasta ese rumbo de mapa
("drive", rumbo, cm[, vel])  girar a ese rumbo y avanzar sosteniéndolo con el IMU
("pose", x, y, rumbo)      posición conocida en el mapa (cm; x este, y sur) para la odometría
("goto", x, y[, vel])      ir a ese punto con odometría encoders + giroscopio (patinaje)
("square", rumbo, s, "x"|"y", valor)  apoyarse en un muro (por tiempo) para fijar esa coordenada
("track\_line", rumbo, "x"|"y", valor, max\_cm)  seguir esa recta del mapa hasta ver una línea
("spin", dir)              girar hasta encontrar la línea (1 der, -1 izq)
("cross", vel, max\_cm[, rumbo])  recto hasta que las dos exteriores ven una línea que cruza
("seek"[, vel, max\_cm, rumbo])  recto con giroscopio hasta que las sondas centrales vean línea
("approach\_ped", altura, dist)  dejar la garra (a esa altura) sobre el centro de un pedestal
                           que está a `dist` cm del último cruce (medida de la pista)
("approach\_obj", altura, dist)  igual, para un objeto a `dist` cm del centro de giro
("retreat",)               regresar lo mismo que avanzó el último acercamiento
("to\_tower",)              seguir la línea hasta quedar junto a la torre (desde la zona verde)
\end{alltt}}
El robot NO lleva ultrasónico: los acercamientos se hacen con el encoder, desde un cruce o el borde de una zona, usando medidas de la pista (software/core/pista.py).

{\small\begin{alltt}
("climb",) / ("stairs"[, extra\_cm])   subir rampa / escalera esperando a quedar plano
("if\_skipped", id, [pasos]) / ("if\_ran", id, [pasos])  según si la misión id se saltó (ya
                           estaba hecha antes de un reinicio) o se corrió en esta pasada
("if\_placed", [pasos])       pasos si el robot recién fue colocado a mano (arranque o reinicio)
("if\_start", nivel, [pasos])  pasos que solo se hacen si el autónomo arrancó en ese nivel
("crest", cm)              subir hasta pasar la cima y avanzar cm (poco espacio arriba)
("claw", apertura) / ("arm", altura)
("hold",)                  sostener 5 s (validación de cubo y pelota)
("detour", lado)           rodear la torre (-1 izquierda, 1 derecha)
("cards",) / ("drop\_red",) / ("stop\_in", color)
\end{alltt}}
critical=True: si falla, el robot quedó fuera de ruta y el tramo se detiene (el equipo pide reinicio en la zona del nivel). critical=False: se sigue con la siguiente misión. end\_level=True: termina el tramo (el robot ya está completamente dentro del siguiente nivel y, en el tramo 1, el cubo ya está en el pedestal): con cfg.SEGMENTS el robot se detiene ahí; el equipo pide reinicio en la zona del siguiente nivel, lo elige con el joystick y arranca.

\noindent Archivo fuente: \texttt{software/autonomous/missions.py} (301 l\'ineas).

\subsubsection{\texttt{junctions(route)}}\label{fn:missions:junctions}
\noindent Tabla (cruces, tipo, giro) -> pasos "line" + "jturn".

Ya no la usa el nivel 1; queda para rutas por línea.

\par\noindent{\small\begin{tabular}{@{}p{0.2\linewidth}p{0.74\linewidth}@{}}
\textbf{Par\'ametro} & \textbf{Descripci\'on} \\
\texttt{route} & Lista de (n cruces, tipo "any"/"X"/"L"/"R", grados a girar). \\
\end{tabular}}\par\smallskip
\noindent\textbf{Devuelve:} Lista de pasos.\par
\lstinputlisting[firstline=179,lastline=184,firstnumber=179]{../../software/autonomous/missions.py}
\subsubsection{\texttt{total\_\allowbreak{}points()}}\label{fn:missions:total-points}
\noindent Puntos posibles de todas las misiones de MISSIONS.

\noindent\textbf{Devuelve:} Suma de "pts".\par
\lstinputlisting[firstline=297,lastline=301,firstnumber=297]{../../software/autonomous/missions.py}
\subsection*{Control manual}
\subsection{\texttt{software/manual/manual.py}}\label{mod:manual}
\textbf{Control manual con el Bluetooth Controller de Makeblock (misiones de alianza).}

Misiones manuales (reglamento): cubos M-A-K-E-X en la barra del nivel 2 (en orden y de frente), activar las rampas (el primer acceso es por la ESCALERA) y 4 cilindros (5 x 15 cm) parados en el área negra.

Mapeo (ver docs/control\_manual.md):

{\small\begin{alltt}
Joystick izq. \(\updownarrow\) (Ly)  avanzar / retroceder      Joystick der. \(\leftrightarrow\) (Rx)  girar
Cruceta               movimientos finos lentos (si los joysticks están en el centro)
L1 / L2               brazo arriba / abajo (mantener)
R1 / R2               garra abrir / cerrar (mantener; entre claw\_open y claw\_cube - 20\textdegree{})
N1                    modo rápido / preciso (LEDs azul / naranja)
N2                    brazo a la altura de la barra (cubos MakeX)
N3 (mantener)         escalera: recto con el giroscopio, potencia de escalera
N4                    brazo arriba (garra sobre la nariz): viajar, girar cerca de cosas
Joystick der. \(\updownarrow\) (Ry)  brazo fino (además de L1/L2)
\end{alltt}}
Ayudas:

\begin{itemize}\setlength\itemsep{0pt}
\item Al avanzar sin girar, el giroscopio SOSTIENE EL RUMBO (rampas, recortes, la oruga izquierda que avanza 2.3 \% menos). Al soltar el giro se toma el rumbo nuevo.
\item hal.drive limita la aceleración mientras va inclinado (no se vuelca en rampas/escalera).
\item El brazo se mueve por ÁNGULO DEL SERVO (0..90\textdegree{}, tabla medida en config.py): a 0\textdegree{} la garra queda baja y al frente (\textasciitilde{}8.5 cm), a 90\textdegree{} arriba sobre la nariz.
\end{itemize}
\noindent Archivo fuente: \texttt{software/manual/manual.py} (206 l\'ineas).

\subsubsection{\texttt{shape(v)}}\label{fn:manual:shape}
\noindent Zona muerta + curva expo del joystick: más precisión cerca del centro.

\par\noindent{\small\begin{tabular}{@{}p{0.2\linewidth}p{0.74\linewidth}@{}}
\textbf{Par\'ametro} & \textbf{Descripci\'on} \\
\texttt{v} & Valor del eje (-100..100). \\
\end{tabular}}\par\smallskip
\noindent\textbf{Devuelve:} Valor con zona muerta (PAD\_DEADZONE) y curva (PAD\_EXPO), -100..100.\par
\noindent\textbf{Llama a:} \hyperref[fn:hal:clamp]{\texttt{hal.clamp}}\par
\lstinputlisting[firstline=41,lastline=46,firstnumber=41]{../../software/manual/manual.py}
\subsubsection{\texttt{mix(forward, turn, limit)}}\label{fn:manual:mix}
\noindent Mezcla arcade -> tanque, escalando para no saturar.

\par\noindent{\small\begin{tabular}{@{}p{0.2\linewidth}p{0.74\linewidth}@{}}
\textbf{Par\'ametro} & \textbf{Descripci\'on} \\
\texttt{forward} & Avance (-100..100). \\
\texttt{turn} & Giro (-100..100; + = derecha). \\
\texttt{limit} & Potencia máxima (PAD\_FAST o PAD\_SLOW). \\
\end{tabular}}\par\smallskip
\noindent\textbf{Devuelve:} (izquierda, derecha) en \% de potencia.\par
\lstinputlisting[firstline=54,lastline=58,firstnumber=54]{../../software/manual/manual.py}
\subsubsection{\texttt{\_\allowbreak{}edge(state, key)}}\label{fn:manual:-edge}
\noindent Detecta el momento en que se presiona una tecla del gamepad.

\par\noindent{\small\begin{tabular}{@{}p{0.2\linewidth}p{0.74\linewidth}@{}}
\textbf{Par\'ametro} & \textbf{Descripci\'on} \\
\texttt{state} & Dict donde se guarda el estado anterior de cada tecla. \\
\texttt{key} & Nombre de la tecla. \\
\end{tabular}}\par\smallskip
\noindent\textbf{Devuelve:} True solo en el ciclo en que se presiona.\par
\noindent\textbf{Llama a:} \hyperref[fn:hal:pad-key]{\texttt{hal.pad\_\allowbreak{}key}}\par
\lstinputlisting[firstline=65,lastline=69,firstnumber=65]{../../software/manual/manual.py}
\subsubsection{\texttt{\_\allowbreak{}heading\_\allowbreak{}hold(f, t)}}\label{fn:manual:-heading-hold}
\noindent Corrección de giro para ir recto con el giroscopio mientras no se pide girar.

Al empezar a avanzar sin girar guarda el rumbo; al pedir un giro lo suelta y al terminar de girar toma el rumbo nuevo.

\par\noindent{\small\begin{tabular}{@{}p{0.2\linewidth}p{0.74\linewidth}@{}}
\textbf{Par\'ametro} & \textbf{Descripci\'on} \\
\texttt{f} & Avance pedido. \\
\texttt{t} & Giro pedido. \\
\end{tabular}}\par\smallskip
\noindent\textbf{Devuelve:} Giro a aplicar (el pedido, o la corrección para sostener el rumbo, \(\pm\)25).\par
\noindent\textbf{Llama a:} \hyperref[fn:hal:clamp]{\texttt{hal.clamp}}, \hyperref[fn:hal:heading]{\texttt{hal.heading}}\par
\lstinputlisting[firstline=78,lastline=87,firstnumber=78]{../../software/manual/manual.py}
\subsubsection{\texttt{\_\allowbreak{}claw\_\allowbreak{}dir()}}\label{fn:manual:-claw-dir}
\noindent Sentido de apertura del servo de la garra.

\noindent\textbf{Devuelve:} +1 si abrir es subir el ángulo (claw\_open >= claw\_cube), -1 si no.\par
\lstinputlisting[firstline=95,lastline=96,firstnumber=95]{../../software/manual/manual.py}
\subsubsection{\texttt{\_\allowbreak{}claw\_\allowbreak{}limit(deg)}}\label{fn:manual:-claw-limit}
\noindent Limita el ángulo de la garra a su recorrido útil.

De claw\_open (abierta) a claw\_cube + CLAW\_SQUEEZE (apretando): abrir o cerrar de más fuerza el servo contra el tope.

\par\noindent{\small\begin{tabular}{@{}p{0.2\linewidth}p{0.74\linewidth}@{}}
\textbf{Par\'ametro} & \textbf{Descripci\'on} \\
\texttt{deg} & Ángulo pedido. \\
\end{tabular}}\par\smallskip
\noindent\textbf{Devuelve:} Ángulo limitado.\par
\noindent\textbf{Llama a:} \hyperref[fn:hal:clamp]{\texttt{hal.clamp}}, \hyperref[fn:manual:-claw-dir]{\texttt{\_\allowbreak{}claw\_\allowbreak{}dir}}\par
\lstinputlisting[firstline=104,lastline=106,firstnumber=104]{../../software/manual/manual.py}
\subsubsection{\texttt{\_\allowbreak{}arm(deg)}}\label{fn:manual:-arm}
\noindent Fija el ángulo del servo del brazo (0..90\textdegree{}).

\par\noindent{\small\begin{tabular}{@{}p{0.2\linewidth}p{0.74\linewidth}@{}}
\textbf{Par\'ametro} & \textbf{Descripci\'on} \\
\texttt{deg} & Ángulo pedido. \\
\end{tabular}}\par\smallskip
\noindent\textbf{Llama a:} \hyperref[fn:hal:clamp]{\texttt{hal.clamp}}\par
\lstinputlisting[firstline=111,lastline=113,firstnumber=111]{../../software/manual/manual.py}
\subsubsection{\texttt{step(state)}}\label{fn:manual:step}
\noindent Un ciclo del control manual (cada DT segundos).

Lee el gamepad: N1 modo rápido/preciso; L1/L2 y Ry brazo; N2 brazo a la barra; N4 brazo arriba; R1/R2 garra; N3 (mantener) escalera recto con el giroscopio; Ly/Rx manejo con sostén de rumbo; cruceta movimientos finos. Mapeo completo en docs/control\_manual.md.

\par\noindent{\small\begin{tabular}{@{}p{0.2\linewidth}p{0.74\linewidth}@{}}
\textbf{Par\'ametro} & \textbf{Descripci\'on} \\
\texttt{state} & Dict de estado de las teclas (para \_edge). \\
\end{tabular}}\par\smallskip
\noindent\textbf{Llama a:} \hyperref[fn:manual:-edge]{\texttt{\_\allowbreak{}edge}}, \hyperref[fn:hal:pad-key]{\texttt{hal.pad\_\allowbreak{}key}}, \hyperref[fn:manual:shape]{\texttt{shape}}, \hyperref[fn:hal:servo]{\texttt{hal.servo}}, \hyperref[fn:manual:-heading-hold]{\texttt{\_\allowbreak{}heading\_\allowbreak{}hold}}, \hyperref[fn:manual:mix]{\texttt{mix}}, \hyperref[fn:hal:drive]{\texttt{hal.drive}}, \hyperref[fn:manual:-arm]{\texttt{\_\allowbreak{}arm}}, \hyperref[fn:hal:pad-axis]{\texttt{hal.pad\_\allowbreak{}axis}}, \hyperref[fn:manual:-claw-limit]{\texttt{\_\allowbreak{}claw\_\allowbreak{}limit}}, \hyperref[fn:hal:clamp]{\texttt{hal.clamp}}, \hyperref[fn:hal:leds]{\texttt{hal.leds}}, \hyperref[fn:geometry:arm-servo]{\texttt{geometry.arm\_\allowbreak{}servo}}, \hyperref[fn:hal:heading]{\texttt{hal.heading}}, \hyperref[fn:manual:-claw-dir]{\texttt{\_\allowbreak{}claw\_\allowbreak{}dir}}\par
\lstinputlisting[firstline=121,lastline=175,firstnumber=121]{../../software/manual/manual.py}
\subsubsection{\texttt{status()}}\label{fn:manual:status}
\noindent Texto de estado para la pantalla.

\noindent\textbf{Devuelve:} Modo (RAPIDO/PRECISO), ángulo del brazo y altura de la garra.\par
\noindent\textbf{Llama a:} \hyperref[fn:geometry:arm-at]{\texttt{geometry.arm\_\allowbreak{}at}}\par
\lstinputlisting[firstline=180,lastline=182,firstnumber=180]{../../software/manual/manual.py}
\subsubsection{\texttt{loop(duration\_\allowbreak{}s)}}\label{fn:manual:loop}
\noindent Bucle de control manual.

Sale al pasar duration\_s o con A (parar) en el CyberPi; al salir detiene el robot.

\par\noindent{\small\begin{tabular}{@{}p{0.2\linewidth}p{0.74\linewidth}@{}}
\textbf{Par\'ametro} & \textbf{Descripci\'on} \\
\texttt{duration\_\allowbreak{}s} & Segundos máximos (None = sin límite). \\
\end{tabular}}\par\smallskip
\noindent\textbf{Llama a:} \hyperref[fn:hal:clear-abort]{\texttt{hal.clear\_\allowbreak{}abort}}, \hyperref[fn:hal:show]{\texttt{hal.show}}, \hyperref[fn:hal:now]{\texttt{hal.now}}, \hyperref[fn:hal:stop]{\texttt{hal.stop}}, \hyperref[fn:hal:pad-ok]{\texttt{hal.pad\_\allowbreak{}ok}}, \hyperref[fn:hal:abort-requested]{\texttt{hal.abort\_\allowbreak{}requested}}, \hyperref[fn:manual:step]{\texttt{step}}, \hyperref[fn:hal:wait]{\texttt{hal.wait}}, \hyperref[fn:manual:status]{\texttt{status}}\par
\lstinputlisting[firstline=188,lastline=206,firstnumber=188]{../../software/manual/manual.py}
\subsection*{Control}
\subsection{\texttt{software/controls/motion.py}}\label{mod:motion}
\textbf{Movimientos con giroscopio y encoders (sin seguir la línea).}

Todas las funciones de movimiento devuelven True si terminaron bien y False si se acabó el tiempo o se pidió parar (botón A del CyberPi, cfg.ABORT\_BUTTON). Convención de giros: positivo = a la DERECHA (igual que hal.heading()).

\noindent Archivo fuente: \texttt{software/controls/motion.py} (413 l\'ineas).

\subsubsection{\texttt{sign(v)}}\label{fn:motion:sign}
\noindent Signo de un número.

\par\noindent{\small\begin{tabular}{@{}p{0.2\linewidth}p{0.74\linewidth}@{}}
\textbf{Par\'ametro} & \textbf{Descripci\'on} \\
\texttt{v} & Número. \\
\end{tabular}}\par\smallskip
\noindent\textbf{Devuelve:} -1 si es negativo, 1 si no.\par
\lstinputlisting[firstline=16,lastline=19,firstnumber=16]{../../software/controls/motion.py}
\subsubsection{\texttt{turn(deg, timeout)}}\label{fn:motion:turn}
\noindent Giro relativo en el lugar.

Control proporcional (turn\_kp) con potencia entre v\_turn\_min y v\_turn\_max; termina cuando el error queda dentro de turn\_tol durante TURN\_SETTLE\_S.

\par\noindent{\small\begin{tabular}{@{}p{0.2\linewidth}p{0.74\linewidth}@{}}
\textbf{Par\'ametro} & \textbf{Descripci\'on} \\
\texttt{deg} & Grados a girar; positivo = derecha. \\
\texttt{timeout} & s máximos (por defecto cfg.TURN\_TIMEOUT\_S). \\
\end{tabular}}\par\smallskip
\noindent\textbf{Devuelve:} True si llegó al ángulo.\par
\noindent\textbf{Llama a:} \hyperref[fn:hal:now]{\texttt{hal.now}}, \hyperref[fn:hal:stop]{\texttt{hal.stop}}, \hyperref[fn:hal:heading]{\texttt{hal.heading}}, \hyperref[fn:hal:abort-requested]{\texttt{hal.abort\_\allowbreak{}requested}}, \hyperref[fn:hal:wait]{\texttt{hal.wait}}, \hyperref[fn:hal:drive]{\texttt{hal.drive}}, \hyperref[fn:motion:sign]{\texttt{sign}}, \hyperref[fn:hal:clamp]{\texttt{hal.clamp}}\par
\lstinputlisting[firstline=28,lastline=51,firstnumber=28]{../../software/controls/motion.py}
\subsubsection{\texttt{pivot(deg, speed, timeout)}}\label{fn:motion:pivot}
\noindent Giro relativo sobre UNA oruga (la otra detenida).

Avanza la oruga de afuera: al girar a la derecha la derecha queda quieta. Si la oruga que se mueve patina (recortes del pasillo) el robot solo gira más lento: el centro del giro es la oruga quieta y la posición final depende solo del ángulo (en un giro en el lugar una oruga que patina corre el robot hasta 6 cm sin que la odometría lo vea).

\par\noindent{\small\begin{tabular}{@{}p{0.2\linewidth}p{0.74\linewidth}@{}}
\textbf{Par\'ametro} & \textbf{Descripci\'on} \\
\texttt{deg} & Grados a girar; positivo = derecha. \\
\texttt{speed} & \% de potencia (por defecto v\_turn\_max); negativo = pivote hacia atrás. \\
\texttt{timeout} & s máximos (por defecto el doble de TURN\_TIMEOUT\_S). \\
\end{tabular}}\par\smallskip
\noindent\textbf{Devuelve:} True si llegó al ángulo.\par
\noindent\textbf{Llama a:} \hyperref[fn:hal:now]{\texttt{hal.now}}, \hyperref[fn:hal:stop]{\texttt{hal.stop}}, \hyperref[fn:hal:heading]{\texttt{hal.heading}}, \hyperref[fn:hal:abort-requested]{\texttt{hal.abort\_\allowbreak{}requested}}, \hyperref[fn:hal:wait]{\texttt{hal.wait}}, \hyperref[fn:hal:clamp]{\texttt{hal.clamp}}, \hyperref[fn:hal:drive]{\texttt{hal.drive}}\par
\lstinputlisting[firstline=63,lastline=93,firstnumber=63]{../../software/controls/motion.py}
\subsubsection{\texttt{wrap180(deg)}}\label{fn:motion:wrap180}
\noindent Lleva un ángulo a -180..180.

\par\noindent{\small\begin{tabular}{@{}p{0.2\linewidth}p{0.74\linewidth}@{}}
\textbf{Par\'ametro} & \textbf{Descripci\'on} \\
\texttt{deg} & Grados. \\
\end{tabular}}\par\smallskip
\noindent\textbf{Devuelve:} Grados equivalentes en -180..180.\par
\lstinputlisting[firstline=99,lastline=104,firstnumber=99]{../../software/controls/motion.py}
\subsubsection{\texttt{turn\_\allowbreak{}to(heading, timeout)}}\label{fn:motion:turn-to}
\noindent Gira hasta un rumbo absoluto por el camino más corto.

\par\noindent{\small\begin{tabular}{@{}p{0.2\linewidth}p{0.74\linewidth}@{}}
\textbf{Par\'ametro} & \textbf{Descripci\'on} \\
\texttt{heading} & Rumbo objetivo (grados de hal.heading()). \\
\texttt{timeout} & s máximos. \\
\end{tabular}}\par\smallskip
\noindent\textbf{Devuelve:} True si llegó.\par
\noindent\textbf{Llama a:} \hyperref[fn:motion:turn]{\texttt{turn}}, \hyperref[fn:motion:wrap180]{\texttt{wrap180}}, \hyperref[fn:hal:heading]{\texttt{hal.heading}}\par
\lstinputlisting[firstline=111,lastline=112,firstnumber=111]{../../software/controls/motion.py}
\subsubsection{\texttt{\_\allowbreak{}hold(speed, target)}}\label{fn:motion:-hold}
\noindent Un ciclo de avance sosteniendo un rumbo con el giroscopio.

El giro sale de la DIFERENCIA entre orugas (izquierda > derecha = a la derecha), igual hacia adelante que hacia atrás. Antes se invertía la corrección al retroceder y el rumbo se iba (en un retroceso largo sobre un recorte el robot daba vueltas).

\par\noindent{\small\begin{tabular}{@{}p{0.2\linewidth}p{0.74\linewidth}@{}}
\textbf{Par\'ametro} & \textbf{Descripci\'on} \\
\texttt{speed} & \% de potencia (negativo = atrás). \\
\texttt{target} & Rumbo a sostener (grados de hal.heading()). \\
\end{tabular}}\par\smallskip
\noindent\textbf{Llama a:} \hyperref[fn:hal:drive]{\texttt{hal.drive}}, \hyperref[fn:hal:heading]{\texttt{hal.heading}}\par
\lstinputlisting[firstline=121,lastline=126,firstnumber=121]{../../software/controls/motion.py}
\subsubsection{\texttt{straight\_\allowbreak{}until(cond, speed, timeout, target)}}\label{fn:motion:straight-until}
\noindent Avanza recto (manteniendo el rumbo) hasta que se cumpla una condición.

\par\noindent{\small\begin{tabular}{@{}p{0.2\linewidth}p{0.74\linewidth}@{}}
\textbf{Par\'ametro} & \textbf{Descripci\'on} \\
\texttt{cond} & Función sin argumentos; el avance termina cuando devuelve True. \\
\texttt{speed} & \% de potencia (por defecto v\_straight; negativo = atrás). \\
\texttt{timeout} & s máximos. \\
\texttt{target} & Rumbo a sostener (por defecto el actual). \\
\end{tabular}}\par\smallskip
\noindent\textbf{Devuelve:} True si terminó por cond().\par
\noindent\textbf{Llama a:} \hyperref[fn:hal:now]{\texttt{hal.now}}, \hyperref[fn:hal:stop]{\texttt{hal.stop}}, \hyperref[fn:hal:heading]{\texttt{hal.heading}}, \hyperref[fn:hal:abort-requested]{\texttt{hal.abort\_\allowbreak{}requested}}, \hyperref[fn:motion:-hold]{\texttt{\_\allowbreak{}hold}}, \hyperref[fn:hal:wait]{\texttt{hal.wait}}\par
\lstinputlisting[firstline=135,lastline=150,firstnumber=135]{../../software/controls/motion.py}
\subsubsection{\texttt{straight(cm, speed, timeout, target)}}\label{fn:motion:straight}
\noindent Avanza o retrocede una distancia medida con encoders.

Sin target se sostiene el rumbo con el que empezó; con target, el error del giro anterior y los golpes del piso se corrigen durante el avance.

\par\noindent{\small\begin{tabular}{@{}p{0.2\linewidth}p{0.74\linewidth}@{}}
\textbf{Par\'ametro} & \textbf{Descripci\'on} \\
\texttt{cm} & Distancia: positiva = adelante, negativa = atrás. \\
\texttt{speed} & \% de potencia (por defecto v\_straight; se usa su valor absoluto). \\
\texttt{timeout} & s máximos (por defecto 1 + |cm| / 8). \\
\texttt{target} & Rumbo (hal.heading) a sostener con el giroscopio. \\
\end{tabular}}\par\smallskip
\noindent\textbf{Devuelve:} True si recorrió la distancia.\par
\noindent\textbf{Llama a:} \hyperref[fn:hal:reset-odometer]{\texttt{hal.reset\_\allowbreak{}odometer}}, \hyperref[fn:motion:straight-until]{\texttt{straight\_\allowbreak{}until}}, \hyperref[fn:motion:sign]{\texttt{sign}}, \hyperref[fn:hal:odometer-cm]{\texttt{hal.odometer\_\allowbreak{}cm}}\par
\lstinputlisting[firstline=161,lastline=169,firstnumber=161]{../../software/controls/motion.py}
\subsubsection{\texttt{climb(speed, timeout, extra\_\allowbreak{}cm)}}\label{fn:motion:climb}
\noindent Sube una rampa o escalera y entra completamente al nivel.

Espera a inclinarse (pitch\_ramp), luego a quedar plano (pitch\_flat durante FLAT\_CONFIRM\_S) y avanza extra\_cm.

\par\noindent{\small\begin{tabular}{@{}p{0.2\linewidth}p{0.74\linewidth}@{}}
\textbf{Par\'ametro} & \textbf{Descripci\'on} \\
\texttt{speed} & \% de potencia (por defecto v\_ramp). \\
\texttt{timeout} & s máximos (por defecto cfg.RAMP\_TIMEOUT\_S). \\
\texttt{extra\_\allowbreak{}cm} & Avance final (por defecto geometry.ramp\_extra()). \\
\end{tabular}}\par\smallskip
\noindent\textbf{Devuelve:} True si subió y entró.\par
\noindent\textbf{Llama a:} \hyperref[fn:hal:heading]{\texttt{hal.heading}}, \hyperref[fn:hal:now]{\texttt{hal.now}}, \hyperref[fn:hal:stop]{\texttt{hal.stop}}, \hyperref[fn:geometry:ramp-extra]{\texttt{geometry.ramp\_\allowbreak{}extra}}, \hyperref[fn:hal:abort-requested]{\texttt{hal.abort\_\allowbreak{}requested}}, \hyperref[fn:hal:pitch]{\texttt{hal.pitch}}, \hyperref[fn:motion:-hold]{\texttt{\_\allowbreak{}hold}}, \hyperref[fn:hal:wait]{\texttt{hal.wait}}, \hyperref[fn:motion:straight]{\texttt{straight}}\par
\lstinputlisting[firstline=179,lastline=209,firstnumber=179]{../../software/controls/motion.py}
\subsubsection{\texttt{stairs(extra\_\allowbreak{}cm)}}\label{fn:motion:stairs}
\noindent Sube la escalera (climb con v\_stairs y STAIRS\_TIMEOUT\_S).

\par\noindent{\small\begin{tabular}{@{}p{0.2\linewidth}p{0.74\linewidth}@{}}
\textbf{Par\'ametro} & \textbf{Descripci\'on} \\
\texttt{extra\_\allowbreak{}cm} & Avance final después de quedar plano. \\
\end{tabular}}\par\smallskip
\noindent\textbf{Devuelve:} True si subió.\par
\noindent\textbf{Llama a:} \hyperref[fn:motion:climb]{\texttt{climb}}\par
\lstinputlisting[firstline=215,lastline=216,firstnumber=215]{../../software/controls/motion.py}
\subsubsection{\texttt{sees\_\allowbreak{}color(name)}}\label{fn:motion:sees-color}
\noindent ¿Alguna de las cuatro sondas ve ese color?

La línea pasa por encima de las zonas, así que las centrales leen negro; las exteriores van a los lados (simulador).

\par\noindent{\small\begin{tabular}{@{}p{0.2\linewidth}p{0.74\linewidth}@{}}
\textbf{Par\'ametro} & \textbf{Descripci\'on} \\
\texttt{name} & Nombre del color. \\
\end{tabular}}\par\smallskip
\noindent\textbf{Devuelve:} True si alguna sonda lo ve.\par
\noindent\textbf{Llama a:} \hyperref[fn:colors:read]{\texttt{colors.read}}\par
\lstinputlisting[firstline=224,lastline=228,firstnumber=224]{../../software/controls/motion.py}
\subsubsection{\texttt{climb\_\allowbreak{}crest(after\_\allowbreak{}cm, speed, timeout, target, color, arm\_\allowbreak{}h)}}\label{fn:motion:climb-crest}
\noindent Sube hasta pasar la cima y avanza after\_cm sin detenerse.

Para cuando arriba no hay espacio para esperar a quedar plano (nivel 4: la torre está a \textasciitilde{}20 cm de la escalera). La cima se detecta cuando el pitch bajó CREST\_DROP grados desde su máximo o, con color, cuando las sondas ven el color de la zona de arriba.

\par\noindent{\small\begin{tabular}{@{}p{0.2\linewidth}p{0.74\linewidth}@{}}
\textbf{Par\'ametro} & \textbf{Descripci\'on} \\
\texttt{after\_\allowbreak{}cm} & cm a avanzar después de la cima. \\
\texttt{speed} & \% de potencia (por defecto v\_stairs). \\
\texttt{timeout} & s máximos. \\
\texttt{target} & Rumbo a sostener. \\
\texttt{color} & Color de la zona de arriba (referencia exacta), o None. \\
\texttt{arm\_\allowbreak{}h} & Altura a la que sube el brazo al llegar (sin parar), o None. \\
\end{tabular}}\par\smallskip
\noindent\textbf{Devuelve:} True si pasó la cima y avanzó.\par
\noindent\textbf{Llama a:} \hyperref[fn:hal:now]{\texttt{hal.now}}, \hyperref[fn:hal:stop]{\texttt{hal.stop}}, \hyperref[fn:hal:heading]{\texttt{hal.heading}}, \hyperref[fn:hal:abort-requested]{\texttt{hal.abort\_\allowbreak{}requested}}, \hyperref[fn:hal:pitch]{\texttt{hal.pitch}}, \hyperref[fn:motion:-hold]{\texttt{\_\allowbreak{}hold}}, \hyperref[fn:hal:wait]{\texttt{hal.wait}}, \hyperref[fn:motion:sees-color]{\texttt{sees\_\allowbreak{}color}}, \hyperref[fn:motion:straight]{\texttt{straight}}, \hyperref[fn:hal:servo]{\texttt{hal.servo}}, \hyperref[fn:geometry:arm-servo]{\texttt{geometry.arm\_\allowbreak{}servo}}\par
\lstinputlisting[firstline=242,lastline=271,firstnumber=242]{../../software/controls/motion.py}
\subsubsection{\texttt{spin\_\allowbreak{}to\_\allowbreak{}line(direction, timeout, leave)}}\label{fn:motion:spin-to-line}
\noindent Gira en el lugar hasta que las sondas centrales vean línea.

\par\noindent{\small\begin{tabular}{@{}p{0.2\linewidth}p{0.74\linewidth}@{}}
\textbf{Par\'ametro} & \textbf{Descripci\'on} \\
\texttt{direction} & 1 = derecha, -1 = izquierda. \\
\texttt{timeout} & s máximos. \\
\texttt{leave} & True: si ya está sobre una línea, primero sale de ella (para buscar OTRA línea). False: si ya ve la línea, termina (tras un giro en un cruce las sondas ya quedan sobre la nueva). \\
\end{tabular}}\par\smallskip
\noindent\textbf{Devuelve:} True si encontró la línea.\par
\noindent\textbf{Llama a:} \hyperref[fn:hal:now]{\texttt{hal.now}}, \hyperref[fn:hal:stop]{\texttt{hal.stop}}, \hyperref[fn:hal:line-bits]{\texttt{hal.line\_\allowbreak{}bits}}, \hyperref[fn:hal:drive]{\texttt{hal.drive}}, \hyperref[fn:hal:wait]{\texttt{hal.wait}}, \hyperref[fn:hal:abort-requested]{\texttt{hal.abort\_\allowbreak{}requested}}\par
\lstinputlisting[firstline=280,lastline=299,firstnumber=280]{../../software/controls/motion.py}
\subsubsection{\texttt{on\_\allowbreak{}line()}}\label{fn:motion:on-line}
\noindent ¿Las sondas centrales (L1 o R1) ven línea?

\noindent\textbf{Devuelve:} True si alguna ve negro.\par
\noindent\textbf{Llama a:} \hyperref[fn:hal:line-bits]{\texttt{hal.line\_\allowbreak{}bits}}\par
\lstinputlisting[firstline=304,lastline=306,firstnumber=304]{../../software/controls/motion.py}
\subsubsection{\texttt{seek\_\allowbreak{}line(speed, max\_\allowbreak{}cm, target)}}\label{fn:motion:seek-line}
\noindent Avanza recto (con giroscopio) hasta que las sondas centrales vean una línea.

\par\noindent{\small\begin{tabular}{@{}p{0.2\linewidth}p{0.74\linewidth}@{}}
\textbf{Par\'ametro} & \textbf{Descripci\'on} \\
\texttt{speed} & \% de potencia (por defecto v\_slow). \\
\texttt{max\_\allowbreak{}cm} & cm máximos. \\
\texttt{target} & Rumbo a sostener. \\
\end{tabular}}\par\smallskip
\noindent\textbf{Devuelve:} True si encontró la línea.\par
\noindent\textbf{Llama a:} \hyperref[fn:hal:reset-odometer]{\texttt{hal.reset\_\allowbreak{}odometer}}, \hyperref[fn:motion:straight-until]{\texttt{straight\_\allowbreak{}until}}, \hyperref[fn:motion:on-line]{\texttt{on\_\allowbreak{}line}}, \hyperref[fn:hal:odometer-cm]{\texttt{hal.odometer\_\allowbreak{}cm}}\par
\lstinputlisting[firstline=314,lastline=328,firstnumber=314]{../../software/controls/motion.py}
\subsubsection{\texttt{drive\_\allowbreak{}to\_\allowbreak{}edge(color, speed, max\_\allowbreak{}cm, target)}}\label{fn:motion:drive-to-edge}
\noindent Avanza recto hasta SALIR de una zona de color.

Referencia fija para empezar una maniobra sin importar dónde quedó o se colocó el robot dentro de la zona. Pide COLOR\_CONFIRM lecturas seguidas fuera de la zona.

\par\noindent{\small\begin{tabular}{@{}p{0.2\linewidth}p{0.74\linewidth}@{}}
\textbf{Par\'ametro} & \textbf{Descripci\'on} \\
\texttt{color} & Color de la zona. \\
\texttt{speed} & \% de potencia (por defecto v\_slow). \\
\texttt{max\_\allowbreak{}cm} & cm máximos. \\
\texttt{target} & Rumbo a sostener. \\
\end{tabular}}\par\smallskip
\noindent\textbf{Devuelve:} True si salió de la zona.\par
\noindent\textbf{Llama a:} \hyperref[fn:hal:reset-odometer]{\texttt{hal.reset\_\allowbreak{}odometer}}, \hyperref[fn:motion:straight-until]{\texttt{straight\_\allowbreak{}until}}, \hyperref[fn:motion:sees-color]{\texttt{sees\_\allowbreak{}color}}, \hyperref[fn:hal:odometer-cm]{\texttt{hal.odometer\_\allowbreak{}cm}}\par
\lstinputlisting[firstline=339,lastline=357,firstnumber=339]{../../software/controls/motion.py}
\subsubsection{\texttt{drive\_\allowbreak{}to\_\allowbreak{}cross(speed, max\_\allowbreak{}cm, target)}}\label{fn:motion:drive-to-cross}
\noindent Avanza recto hasta una línea que cruza el camino.

Las DOS sondas exteriores deben ver negro fuerte (cross\_th), cada una con memoria (pueden verla en distinto momento si la línea no es perpendicular), a menos de cross\_win cm una de otra. No es seguir la línea: sirve de referencia para corregir la distancia. Las marcas verdes y el tope no cuentan.

\par\noindent{\small\begin{tabular}{@{}p{0.2\linewidth}p{0.74\linewidth}@{}}
\textbf{Par\'ametro} & \textbf{Descripci\'on} \\
\texttt{speed} & \% de potencia (por defecto v\_slow). \\
\texttt{max\_\allowbreak{}cm} & cm máximos. \\
\texttt{target} & Rumbo a sostener. \\
\end{tabular}}\par\smallskip
\noindent\textbf{Devuelve:} True si encontró el cruce.\par
\noindent\textbf{Llama a:} \hyperref[fn:hal:reset-odometer]{\texttt{hal.reset\_\allowbreak{}odometer}}, \hyperref[fn:motion:straight-until]{\texttt{straight\_\allowbreak{}until}}, \hyperref[fn:hal:line-read]{\texttt{hal.line\_\allowbreak{}read}}, \hyperref[fn:hal:odometer-cm]{\texttt{hal.odometer\_\allowbreak{}cm}}\par
\lstinputlisting[firstline=369,lastline=389,firstnumber=369]{../../software/controls/motion.py}
\subsubsection{\texttt{detour(side)}}\label{fn:motion:detour}
\noindent Rodea un obstáculo fijo (torre) y vuelve a la línea.

Distancias desde geometry.py (detour\_out, detour\_pass). Al final gira hacia el mismo lado del rodeo para seguir en el rumbo original.

\par\noindent{\small\begin{tabular}{@{}p{0.2\linewidth}p{0.74\linewidth}@{}}
\textbf{Par\'ametro} & \textbf{Descripci\'on} \\
\texttt{side} & -1 = por la izquierda, 1 = por la derecha. \\
\end{tabular}}\par\smallskip
\noindent\textbf{Devuelve:} True si volvió a la línea.\par
\noindent\textbf{Llama a:} \hyperref[fn:hal:heading]{\texttt{hal.heading}}, \hyperref[fn:motion:turn]{\texttt{turn}}, \hyperref[fn:motion:straight]{\texttt{straight}}, \hyperref[fn:motion:spin-to-line]{\texttt{spin\_\allowbreak{}to\_\allowbreak{}line}}, \hyperref[fn:motion:turn-to]{\texttt{turn\_\allowbreak{}to}}, \hyperref[fn:motion:straight-until]{\texttt{straight\_\allowbreak{}until}}, \hyperref[fn:geometry:turn-advance]{\texttt{geometry.turn\_\allowbreak{}advance}}, \hyperref[fn:geometry:detour-out]{\texttt{geometry.detour\_\allowbreak{}out}}, \hyperref[fn:geometry:detour-pass]{\texttt{geometry.detour\_\allowbreak{}pass}}\par
\lstinputlisting[firstline=397,lastline=413,firstnumber=397]{../../software/controls/motion.py}
\subsection{\texttt{software/controls/linefollow.py}}\label{mod:linefollow}
\textbf{Seguidor de línea PID con detección de cruces robusta para línea delgada.}

Cruces "con memoria": con una línea de 10-30 mm y el robot oscilando, las dos sondas exteriores casi nunca ven negro AL MISMO TIEMPO. Por eso un cruce se arma con lo que vean L2 y R2 dentro de una ventana de avance (P["cross\_win"] cm, medida con encoders): "X" = las dos exteriores vieron negro, "L" = solo la izquierda, "R" = solo la derecha. Los giros se escriben en la ruta (missions.py); las marcas verdes son opcionales (P["green"]). Las funciones at\_junction, at\_color, color\_then\_cm, after\_cm, leave\_color, obstacle\_within, line\_end y pitch\_above crean CONDICIONES DE PARADA cond(bits, junction) para follow\_until().

\noindent Archivo fuente: \texttt{software/controls/linefollow.py} (345 l\'ineas).

\subsubsection{\texttt{reset()}}\label{fn:linefollow:reset}
\noindent Reinicia el PID, la vuelta pendiente de las marcas verdes y el rastreador de cruces.

\noindent\textbf{Llama a:} \hyperref[fn:linefollow:reset-junctions]{\texttt{reset\_\allowbreak{}junctions}}\par
\lstinputlisting[firstline=34,lastline=39,firstnumber=34]{../../software/controls/linefollow.py}
\subsubsection{\texttt{reset\_\allowbreak{}junctions()}}\label{fn:linefollow:reset-junctions}
\noindent Reinicia el rastreador de cruces (queda armado).

\lstinputlisting[firstline=43,lastline=49,firstnumber=43]{../../software/controls/linefollow.py}
\subsubsection{\texttt{correction(err, last\_\allowbreak{}err, integral)}}\label{fn:linefollow:correction}
\noindent Corrección PID (función pura, probada en PC).

\par\noindent{\small\begin{tabular}{@{}p{0.2\linewidth}p{0.74\linewidth}@{}}
\textbf{Par\'ametro} & \textbf{Descripci\'on} \\
\texttt{err} & Error actual (offset de la línea, -100..100). \\
\texttt{last\_\allowbreak{}err} & Error del ciclo anterior. \\
\texttt{integral} & Suma de errores. \\
\end{tabular}}\par\smallskip
\noindent\textbf{Devuelve:} Corrección en \% de potencia (kp\(\cdot\)err + ki\(\cdot\)integral + kd\(\cdot\)(err - last\_err)).\par
\lstinputlisting[firstline=57,lastline=58,firstnumber=57]{../../software/controls/linefollow.py}
\subsubsection{\texttt{drive\_\allowbreak{}with(err, speed)}}\label{fn:linefollow:drive-with}
\noindent Aplica la corrección PID a las orugas.

Línea a la izquierda (err < 0) -> frena la oruga izquierda. Con bias\_cm el robot va corrido de la línea esos cm.

\par\noindent{\small\begin{tabular}{@{}p{0.2\linewidth}p{0.74\linewidth}@{}}
\textbf{Par\'ametro} & \textbf{Descripci\'on} \\
\texttt{err} & Offset de la línea (-100..100). \\
\texttt{speed} & \% de potencia base. \\
\end{tabular}}\par\smallskip
\noindent\textbf{Llama a:} \hyperref[fn:hal:clamp]{\texttt{hal.clamp}}, \hyperref[fn:linefollow:correction]{\texttt{correction}}, \hyperref[fn:hal:drive]{\texttt{hal.drive}}\par
\lstinputlisting[firstline=66,lastline=75,firstnumber=66]{../../software/controls/linefollow.py}
\subsubsection{\texttt{step(speed)}}\label{fn:linefollow:step}
\noindent Un ciclo de control del seguidor.

\par\noindent{\small\begin{tabular}{@{}p{0.2\linewidth}p{0.74\linewidth}@{}}
\textbf{Par\'ametro} & \textbf{Descripci\'on} \\
\texttt{speed} & \% de potencia base. \\
\end{tabular}}\par\smallskip
\noindent\textbf{Devuelve:} Bits leídos (dict sonda -> ve línea).\par
\noindent\textbf{Llama a:} \hyperref[fn:hal:line-read]{\texttt{hal.line\_\allowbreak{}read}}, \hyperref[fn:linefollow:drive-with]{\texttt{drive\_\allowbreak{}with}}\par
\lstinputlisting[firstline=81,lastline=84,firstnumber=81]{../../software/controls/linefollow.py}
\subsubsection{\texttt{sees\_\allowbreak{}line(bits)}}\label{fn:linefollow:sees-line}
\noindent ¿Alguna sonda ve línea?

\par\noindent{\small\begin{tabular}{@{}p{0.2\linewidth}p{0.74\linewidth}@{}}
\textbf{Par\'ametro} & \textbf{Descripci\'on} \\
\texttt{bits} & Dict de hal.line\_read(). \\
\end{tabular}}\par\smallskip
\noindent\textbf{Devuelve:} True si alguna de las cuatro ve negro.\par
\lstinputlisting[firstline=90,lastline=91,firstnumber=90]{../../software/controls/linefollow.py}
\subsubsection{\texttt{track\_\allowbreak{}junction(bits, odo)}}\label{fn:linefollow:track-junction}
\noindent Actualiza el rastreador de cruces con una lectura.

Tras cerrar un cruce queda desarmado hasta recorrer 1.5 cm con las exteriores en blanco (no cuenta dos veces el mismo cruce). Guarda el cruce en last\_junction.

\par\noindent{\small\begin{tabular}{@{}p{0.2\linewidth}p{0.74\linewidth}@{}}
\textbf{Par\'ametro} & \textbf{Descripci\'on} \\
\texttt{bits} & Dict de hal.line\_read(). \\
\texttt{odo} & Odometría absoluta (cm). \\
\end{tabular}}\par\smallskip
\noindent\textbf{Devuelve:} "X", "L" o "R" en el ciclo en que se cierra un cruce; si no, None.\par
\lstinputlisting[firstline=100,lastline=130,firstnumber=100]{../../software/controls/linefollow.py}
\subsubsection{\texttt{\_\allowbreak{}check\_\allowbreak{}green()}}\label{fn:linefollow:-check-green}
\noindent Lee las marcas verdes con las sondas exteriores y deja pendiente la vuelta.

Verde a los dos lados = media vuelta (2); a la izquierda = -1; a la derecha = 1.

\noindent\textbf{Llama a:} \hyperref[fn:colors:read]{\texttt{colors.read}}\par
\lstinputlisting[firstline=135,lastline=144,firstnumber=135]{../../software/controls/linefollow.py}
\subsubsection{\texttt{turn\_\allowbreak{}at\_\allowbreak{}junction(deg)}}\label{fn:linefollow:turn-at-junction}
\noindent Gira EN el último cruce (paso "jturn") y se alinea con la línea nueva.

Avanza hasta que el CENTRO DE GIRO quede sobre el cruce, gira el 85 \% con el giroscopio y termina buscando la línea (spin\_to\_line). El otro brazo del cruce sigue bajo las sondas: el siguiente seguimiento empieza desarmado.

\par\noindent{\small\begin{tabular}{@{}p{0.2\linewidth}p{0.74\linewidth}@{}}
\textbf{Par\'ametro} & \textbf{Descripci\'on} \\
\texttt{deg} & Grados; positivo = derecha, 0 = solo centrarse. \\
\end{tabular}}\par\smallskip
\noindent\textbf{Devuelve:} True si encontró la línea.\par
\noindent\textbf{Llama a:} \hyperref[fn:motion:turn]{\texttt{motion.turn}}, \hyperref[fn:motion:spin-to-line]{\texttt{motion.spin\_\allowbreak{}to\_\allowbreak{}line}}, \hyperref[fn:linefollow:reset]{\texttt{reset}}, \hyperref[fn:geometry:turn-advance]{\texttt{geometry.turn\_\allowbreak{}advance}}, \hyperref[fn:motion:straight]{\texttt{motion.straight}}, \hyperref[fn:hal:odo-abs]{\texttt{hal.odo\_\allowbreak{}abs}}\par
\lstinputlisting[firstline=153,lastline=168,firstnumber=153]{../../software/controls/linefollow.py}
\subsubsection{\texttt{follow\_\allowbreak{}until(cond, speed, timeout, green)}}\label{fn:linefollow:follow-until}
\noindent Sigue la línea hasta que se cumpla una condición.

\par\noindent{\small\begin{tabular}{@{}p{0.2\linewidth}p{0.74\linewidth}@{}}
\textbf{Par\'ametro} & \textbf{Descripci\'on} \\
\texttt{cond} & Función cond(bits, junction); junction es "X"/"L"/"R" en el ciclo en que se cierra un cruce, si no None. \\
\texttt{speed} & \% de potencia (por defecto v\_line). \\
\texttt{timeout} & s máximos. \\
\texttt{green} & True = girar con las marcas verdes (si P["green"] == 1). \\
\end{tabular}}\par\smallskip
\noindent\textbf{Devuelve:} True si se cumplió cond; False si se acabó el tiempo, se pidió parar o la línea se perdió más de P["lost\_s"] segundos.\par
\noindent\textbf{Llama a:} \hyperref[fn:linefollow:reset]{\texttt{reset}}, \hyperref[fn:hal:now]{\texttt{hal.now}}, \hyperref[fn:hal:stop]{\texttt{hal.stop}}, \hyperref[fn:hal:abort-requested]{\texttt{hal.abort\_\allowbreak{}requested}}, \hyperref[fn:hal:line-read]{\texttt{hal.line\_\allowbreak{}read}}, \hyperref[fn:linefollow:track-junction]{\texttt{track\_\allowbreak{}junction}}, \hyperref[fn:linefollow:sees-line]{\texttt{sees\_\allowbreak{}line}}, \hyperref[fn:linefollow:drive-with]{\texttt{drive\_\allowbreak{}with}}, \hyperref[fn:hal:wait]{\texttt{hal.wait}}, \hyperref[fn:hal:odo-abs]{\texttt{hal.odo\_\allowbreak{}abs}}, \hyperref[fn:linefollow:-check-green]{\texttt{\_\allowbreak{}check\_\allowbreak{}green}}, \hyperref[fn:linefollow:turn-at-junction]{\texttt{turn\_\allowbreak{}at\_\allowbreak{}junction}}\par
\lstinputlisting[firstline=179,lastline=214,firstnumber=179]{../../software/controls/linefollow.py}
\subsubsection{\texttt{at\_\allowbreak{}junction(count, kind)}}\label{fn:linefollow:at-junction}
\noindent Condición: contar cruces.

\par\noindent{\small\begin{tabular}{@{}p{0.2\linewidth}p{0.74\linewidth}@{}}
\textbf{Par\'ametro} & \textbf{Descripci\'on} \\
\texttt{count} & Cruces a contar. \\
\texttt{kind} & "any", "X", "L" o "R" ("L" y "R" también cuentan las "X"). \\
\end{tabular}}\par\smallskip
\noindent\textbf{Devuelve:} Función cond(bits, junction) para follow\_until().\par
\lstinputlisting[firstline=222,lastline=229,firstnumber=222]{../../software/controls/linefollow.py}
\subsubsection{\texttt{at\_\allowbreak{}color(name, probes)}}\label{fn:linefollow:at-color}
\noindent Condición: alguna de las sondas indicadas ve un color.

\par\noindent{\small\begin{tabular}{@{}p{0.2\linewidth}p{0.74\linewidth}@{}}
\textbf{Par\'ametro} & \textbf{Descripci\'on} \\
\texttt{name} & Nombre del color. \\
\texttt{probes} & Sondas a revisar (por defecto las centrales). \\
\end{tabular}}\par\smallskip
\noindent\textbf{Devuelve:} Función cond(bits, junction) para follow\_until().\par
\noindent\textbf{Llama a:} \hyperref[fn:colors:read]{\texttt{colors.read}}\par
\lstinputlisting[firstline=236,lastline=242,firstnumber=236]{../../software/controls/linefollow.py}
\subsubsection{\texttt{color\_\allowbreak{}then\_\allowbreak{}cm(name, cm, probes)}}\label{fn:linefollow:color-then-cm}
\noindent Condición: cm recorridos después de que alguna sonda vio un color.

Sirve para terminar arriba de una rampa SIN detenerse sobre la pendiente (parado ahí el robot se resbalaba hacia atrás al arrancar; simulador con una rampa de 40\textdegree{}).

\par\noindent{\small\begin{tabular}{@{}p{0.2\linewidth}p{0.74\linewidth}@{}}
\textbf{Par\'ametro} & \textbf{Descripci\'on} \\
\texttt{name} & Nombre del color. \\
\texttt{cm} & cm a recorrer después de verlo. \\
\texttt{probes} & Sondas a revisar (por defecto las cuatro). \\
\end{tabular}}\par\smallskip
\noindent\textbf{Devuelve:} Función cond(bits, junction) para follow\_until().\par
\noindent\textbf{Llama a:} \hyperref[fn:hal:odo-abs]{\texttt{hal.odo\_\allowbreak{}abs}}, \hyperref[fn:colors:read]{\texttt{colors.read}}\par
\lstinputlisting[firstline=252,lastline=265,firstnumber=252]{../../software/controls/linefollow.py}
\subsubsection{\texttt{after\_\allowbreak{}cm(cm)}}\label{fn:linefollow:after-cm}
\noindent Condición: recorrer una distancia.

\par\noindent{\small\begin{tabular}{@{}p{0.2\linewidth}p{0.74\linewidth}@{}}
\textbf{Par\'ametro} & \textbf{Descripci\'on} \\
\texttt{cm} & cm desde que se crea la condición. \\
\end{tabular}}\par\smallskip
\noindent\textbf{Devuelve:} Función cond(bits, junction) para follow\_until().\par
\noindent\textbf{Llama a:} \hyperref[fn:hal:odo-abs]{\texttt{hal.odo\_\allowbreak{}abs}}\par
\lstinputlisting[firstline=271,lastline=276,firstnumber=271]{../../software/controls/linefollow.py}
\subsubsection{\texttt{leave\_\allowbreak{}color(name, probes)}}\label{fn:linefollow:leave-color}
\noindent Condición: salir de una zona de color.

Verdadera cuando las sondas, después de haber visto la zona, la dejan de ver (COLOR\_CONFIRM lecturas seguidas). Guarda la odometría en last\_edge: es la referencia para medir desde el borde de la zona.

\par\noindent{\small\begin{tabular}{@{}p{0.2\linewidth}p{0.74\linewidth}@{}}
\textbf{Par\'ametro} & \textbf{Descripci\'on} \\
\texttt{name} & Nombre del color. \\
\texttt{probes} & Sondas a revisar (por defecto las cuatro: la línea cruza la zona y las centrales leen negro). \\
\end{tabular}}\par\smallskip
\noindent\textbf{Devuelve:} Función cond(bits, junction) para follow\_until().\par
\noindent\textbf{Llama a:} \hyperref[fn:hal:odo-abs]{\texttt{hal.odo\_\allowbreak{}abs}}, \hyperref[fn:colors:read]{\texttt{colors.read}}\par
\lstinputlisting[firstline=287,lastline=307,firstnumber=287]{../../software/controls/linefollow.py}
\subsubsection{\texttt{obstacle\_\allowbreak{}within(cm)}}\label{fn:linefollow:obstacle-within}
\noindent Condición: obstáculo a menos de cm (solo con ultrasónico, cfg.HAS\_US).

\par\noindent{\small\begin{tabular}{@{}p{0.2\linewidth}p{0.74\linewidth}@{}}
\textbf{Par\'ametro} & \textbf{Descripci\'on} \\
\texttt{cm} & Distancia en cm. \\
\end{tabular}}\par\smallskip
\noindent\textbf{Devuelve:} Función cond(bits, junction) para follow\_until().\par
\noindent\textbf{Llama a:} \hyperref[fn:hal:distance-cm]{\texttt{hal.distance\_\allowbreak{}cm}}\par
\lstinputlisting[firstline=313,lastline=317,firstnumber=313]{../../software/controls/linefollow.py}
\subsubsection{\texttt{line\_\allowbreak{}end(cm)}}\label{fn:linefollow:line-end}
\noindent Condición: fin de la línea.

Ninguna sonda ve línea durante cm de avance (en una curva cerrada la línea se pierde un instante y no es el final; simulador).

\par\noindent{\small\begin{tabular}{@{}p{0.2\linewidth}p{0.74\linewidth}@{}}
\textbf{Par\'ametro} & \textbf{Descripci\'on} \\
\texttt{cm} & cm sin línea. \\
\end{tabular}}\par\smallskip
\noindent\textbf{Devuelve:} Función cond(bits, junction) para follow\_until().\par
\noindent\textbf{Llama a:} \hyperref[fn:linefollow:sees-line]{\texttt{sees\_\allowbreak{}line}}, \hyperref[fn:hal:odo-abs]{\texttt{hal.odo\_\allowbreak{}abs}}\par
\lstinputlisting[firstline=325,lastline=336,firstnumber=325]{../../software/controls/linefollow.py}
\subsubsection{\texttt{pitch\_\allowbreak{}above(deg)}}\label{fn:linefollow:pitch-above}
\noindent Condición: empieza una rampa (pitch mayor o igual a deg).

\par\noindent{\small\begin{tabular}{@{}p{0.2\linewidth}p{0.74\linewidth}@{}}
\textbf{Par\'ametro} & \textbf{Descripci\'on} \\
\texttt{deg} & Grados de inclinación. \\
\end{tabular}}\par\smallskip
\noindent\textbf{Devuelve:} Función cond(bits, junction) para follow\_until().\par
\noindent\textbf{Llama a:} \hyperref[fn:hal:pitch]{\texttt{hal.pitch}}\par
\lstinputlisting[firstline=342,lastline=345,firstnumber=342]{../../software/controls/linefollow.py}
\subsection{\texttt{software/controls/odometry.py}}\label{mod:odometry}
\textbf{Odometría de posición (x, y, rumbo) con encoders + giroscopio, robusta al patinaje.}

Marco del mapa (el de las vistas del reglamento): x hacia el este desde la orilla izquierda, y hacia el SUR desde la orilla de arriba, en cm; rumbo 0 = este, 90 = sur (horario positivo, igual que hal.heading()). Patinaje (recortes de cartón en la zona de obstáculos): la oruga que pisa uno gira sin avanzar y su encoder sigue contando. El giroscopio dice cuánto giró DE VERDAD el robot y los encoders cuánto "debió" girar; la diferencia dice qué oruga patinó, y el avance se mide con la otra:

\begin{itemize}\setlength\itemsep{0pt}
\item (adelante) giró más a la izquierda de lo esperado -> patinó la izquierda -> se mide con la derecha
\item (adelante) giró más a la derecha de lo esperado -> patinó la derecha -> se mide con la izquierda
\item en reversa los signos se invierten (la que patina avanza menos en su sentido)
\end{itemize}
B = distancia entre orugas * P["skid"] (las orugas derrapan al girar; CALIBRAR). update() se registra en hal.\_on\_wait: se integra en cada ciclo de control.

\noindent Archivo fuente: \texttt{software/controls/odometry.py} (202 l\'ineas).

\subsubsection{\texttt{set\_\allowbreak{}pose(x, y, heading)}}\label{fn:odometry:set-pose}
\noindent Fija la posición conocida del robot (paso "pose" de las misiones).

\par\noindent{\small\begin{tabular}{@{}p{0.2\linewidth}p{0.74\linewidth}@{}}
\textbf{Par\'ametro} & \textbf{Descripci\'on} \\
\texttt{x} & cm hacia el este. \\
\texttt{y} & cm hacia el sur. \\
\texttt{heading} & Rumbo de mapa en grados (0 este, 90 sur). \\
\end{tabular}}\par\smallskip
\noindent\textbf{Llama a:} \hyperref[fn:hal:tracks-cm]{\texttt{hal.tracks\_\allowbreak{}cm}}, \hyperref[fn:hal:heading]{\texttt{hal.heading}}\par
\lstinputlisting[firstline=32,lastline=39,firstnumber=32]{../../software/controls/odometry.py}
\subsubsection{\texttt{heading()}}\label{fn:odometry:heading}
\noindent Rumbo de mapa actual.

\noindent\textbf{Devuelve:} Grados (0 este, 90 sur, 180 oeste, -90 norte).\par
\noindent\textbf{Llama a:} \hyperref[fn:hal:heading]{\texttt{hal.heading}}\par
\lstinputlisting[firstline=44,lastline=45,firstnumber=44]{../../software/controls/odometry.py}
\subsubsection{\texttt{pose()}}\label{fn:odometry:pose}
\noindent Posición estimada.

\noindent\textbf{Devuelve:} (x, y, rumbo) en cm y grados de mapa.\par
\noindent\textbf{Llama a:} \hyperref[fn:odometry:heading]{\texttt{heading}}\par
\lstinputlisting[firstline=50,lastline=51,firstnumber=50]{../../software/controls/odometry.py}
\subsubsection{\texttt{update()}}\label{fn:odometry:update}
\noindent Integra el movimiento desde la última llamada.

Elige con qué oruga medir el avance según el patinaje: en un pivote, la oruga quieta (no puede patinar); en un giro en el lugar o sin patinaje, el promedio; si no, la que no patinó (ver la explicación del archivo). Acumula en slip\_cm lo que contaron de más.

\noindent\textbf{Llama a:} \hyperref[fn:hal:tracks-cm]{\texttt{hal.tracks\_\allowbreak{}cm}}, \hyperref[fn:hal:heading]{\texttt{hal.heading}}, \hyperref[fn:geometry:track-w]{\texttt{geometry.track\_\allowbreak{}w}}\par
\lstinputlisting[firstline=58,lastline=89,firstnumber=58]{../../software/controls/odometry.py}
\subsubsection{\texttt{wrap(d)}}\label{fn:odometry:wrap}
\noindent Lleva un ángulo a -180..180.

\par\noindent{\small\begin{tabular}{@{}p{0.2\linewidth}p{0.74\linewidth}@{}}
\textbf{Par\'ametro} & \textbf{Descripci\'on} \\
\texttt{d} & Grados. \\
\end{tabular}}\par\smallskip
\noindent\textbf{Devuelve:} Grados equivalentes en -180..180.\par
\lstinputlisting[firstline=100,lastline=105,firstnumber=100]{../../software/controls/odometry.py}
\subsubsection{\texttt{goto(x, y, speed, tol, timeout)}}\label{fn:odometry:goto}
\noindent Va a un punto del mapa.

Gira hacia él y avanza corrigiendo el rumbo hacia el punto con la posición estimada (si un recorte lo desvía, el controlador lo regresa). Termina al llegar o al rebasarlo; frena en los últimos 6 cm.

\par\noindent{\small\begin{tabular}{@{}p{0.2\linewidth}p{0.74\linewidth}@{}}
\textbf{Par\'ametro} & \textbf{Descripci\'on} \\
\texttt{x} & cm hacia el este. \\
\texttt{y} & cm hacia el sur. \\
\texttt{speed} & \% de potencia (por defecto v\_slow). \\
\texttt{tol} & cm de tolerancia para dar por llegado. \\
\texttt{timeout} & s máximos (por defecto 2 + distancia / 5). \\
\end{tabular}}\par\smallskip
\noindent\textbf{Devuelve:} True si llegó.\par
\noindent\textbf{Llama a:} \hyperref[fn:odometry:update]{\texttt{update}}, \hyperref[fn:hal:now]{\texttt{hal.now}}, \hyperref[fn:hal:stop]{\texttt{hal.stop}}, \hyperref[fn:motion:turn-to]{\texttt{motion.turn\_\allowbreak{}to}}, \hyperref[fn:hal:abort-requested]{\texttt{hal.abort\_\allowbreak{}requested}}, \hyperref[fn:odometry:wrap]{\texttt{wrap}}, \hyperref[fn:hal:drive]{\texttt{hal.drive}}, \hyperref[fn:hal:wait]{\texttt{hal.wait}}, \hyperref[fn:hal:heading]{\texttt{hal.heading}}, \hyperref[fn:odometry:heading]{\texttt{heading}}\par
\lstinputlisting[firstline=118,lastline=154,firstnumber=118]{../../software/controls/odometry.py}
\subsubsection{\texttt{track(x0, y0, hdg, speed, timeout, max\_\allowbreak{}cm, stop, max\_\allowbreak{}deg, halt)}}\label{fn:odometry:track}
\noindent Sigue una recta del mapa corrigiendo rumbo y desvío lateral.

Si el patinaje lo empuja de lado, regresa a la recta (ángulo máximo de corrección max\_deg). Termina cuando stop() es verdadero, al recorrer max\_cm o al acabarse el tiempo.

\par\noindent{\small\begin{tabular}{@{}p{0.2\linewidth}p{0.74\linewidth}@{}}
\textbf{Par\'ametro} & \textbf{Descripci\'on} \\
\texttt{x0} & x de un punto de la recta (cm). \\
\texttt{y0} & y de un punto de la recta (cm). \\
\texttt{hdg} & Rumbo de la recta (grados de mapa) = hacia donde mira el robot. \\
\texttt{speed} & \% de potencia; negativo = hacia atrás por la misma recta. \\
\texttt{timeout} & s máximos. \\
\texttt{max\_\allowbreak{}cm} & cm máximos a recorrer sobre la recta. \\
\texttt{stop} & Función sin argumentos que dice cuándo terminar (p. ej. "ve la línea"). \\
\texttt{max\_\allowbreak{}deg} & Ángulo máximo contra la recta al corregir (por defecto track\_max). \\
\texttt{halt} & False = al cumplirse stop() no frena (encadenar otro tramo sin parar en una rampa). \\
\end{tabular}}\par\smallskip
\noindent\textbf{Devuelve:} True si terminó por stop().\par
\noindent\textbf{Llama a:} \hyperref[fn:odometry:update]{\texttt{update}}, \hyperref[fn:hal:now]{\texttt{hal.now}}, \hyperref[fn:hal:stop]{\texttt{hal.stop}}, \hyperref[fn:hal:abort-requested]{\texttt{hal.abort\_\allowbreak{}requested}}, \hyperref[fn:odometry:wrap]{\texttt{wrap}}, \hyperref[fn:hal:drive]{\texttt{hal.drive}}, \hyperref[fn:hal:wait]{\texttt{hal.wait}}, \hyperref[fn:hal:clamp]{\texttt{hal.clamp}}, \hyperref[fn:odometry:heading]{\texttt{heading}}\par
\lstinputlisting[firstline=170,lastline=202,firstnumber=170]{../../software/controls/odometry.py}
\subsection{\texttt{software/controls/gripper.py}}\label{mod:gripper}
\textbf{Garra (servo de cierre) y brazo (servo de altura).}

El brazo se manda por ALTURA en cm (geometry.arm\_servo convierte a ángulo del servo), así que al cambiar el brazo solo se mide de nuevo la tabla arm\_h* / arm\_x* (config.py).

\noindent Archivo fuente: \texttt{software/controls/gripper.py} (57 l\'ineas).

\subsubsection{\texttt{\_\allowbreak{}move(port, angle)}}\label{fn:gripper:-move}
\noindent Mueve un servo y espera a que llegue (cfg.SERVO\_STEP\_S).

\par\noindent{\small\begin{tabular}{@{}p{0.2\linewidth}p{0.74\linewidth}@{}}
\textbf{Par\'ametro} & \textbf{Descripci\'on} \\
\texttt{port} & Puerto del servo. \\
\texttt{angle} & Ángulo en grados. \\
\end{tabular}}\par\smallskip
\noindent\textbf{Llama a:} \hyperref[fn:hal:servo]{\texttt{hal.servo}}, \hyperref[fn:hal:wait]{\texttt{hal.wait}}\par
\lstinputlisting[firstline=14,lastline=16,firstnumber=14]{../../software/controls/gripper.py}
\subsubsection{\texttt{claw(angle)}}\label{fn:gripper:claw}
\noindent Mueve la garra a una apertura.

\par\noindent{\small\begin{tabular}{@{}p{0.2\linewidth}p{0.74\linewidth}@{}}
\textbf{Par\'ametro} & \textbf{Descripci\'on} \\
\texttt{angle} & Ángulo del servo de la garra (claw\_open = abierta, claw\_cube = cerrada en el cubo). \\
\end{tabular}}\par\smallskip
\noindent\textbf{Llama a:} \hyperref[fn:gripper:-move]{\texttt{\_\allowbreak{}move}}\par
\lstinputlisting[firstline=21,lastline=22,firstnumber=21]{../../software/controls/gripper.py}
\subsubsection{\texttt{open\_\allowbreak{}claw()}}\label{fn:gripper:open-claw}
\noindent Abre la garra (cfg.P["claw\_open"]).

\noindent\textbf{Llama a:} \hyperref[fn:gripper:claw]{\texttt{claw}}\par
\lstinputlisting[firstline=26,lastline=27,firstnumber=26]{../../software/controls/gripper.py}
\subsubsection{\texttt{arm\_\allowbreak{}to(h)}}\label{fn:gripper:arm-to}
\noindent Sube o baja el brazo hasta que la garra quede a una altura.

\par\noindent{\small\begin{tabular}{@{}p{0.2\linewidth}p{0.74\linewidth}@{}}
\textbf{Par\'ametro} & \textbf{Descripci\'on} \\
\texttt{h} & Altura del centro de la garra sobre el piso (cm). \\
\end{tabular}}\par\smallskip
\noindent\textbf{Llama a:} \hyperref[fn:gripper:-move]{\texttt{\_\allowbreak{}move}}, \hyperref[fn:geometry:arm-servo]{\texttt{geometry.arm\_\allowbreak{}servo}}\par
\lstinputlisting[firstline=32,lastline=33,firstnumber=32]{../../software/controls/gripper.py}
\subsubsection{\texttt{home()}}\label{fn:gripper:home}
\noindent Posición de viaje: brazo arriba (h\_top) y garra abierta.

Garra sobre la nariz: lo más compacto para colocar el robot en las zonas de inicio, que son más chicas que él.

\noindent\textbf{Llama a:} \hyperref[fn:gripper:arm-to]{\texttt{arm\_\allowbreak{}to}}, \hyperref[fn:gripper:open-claw]{\texttt{open\_\allowbreak{}claw}}\par
\lstinputlisting[firstline=39,lastline=41,firstnumber=39]{../../software/controls/gripper.py}
\subsubsection{\texttt{hold(seconds)}}\label{fn:gripper:hold}
\noindent Sostiene el objeto quieto el tiempo que pide el reglamento (>= 5 s).

\par\noindent{\small\begin{tabular}{@{}p{0.2\linewidth}p{0.74\linewidth}@{}}
\textbf{Par\'ametro} & \textbf{Descripci\'on} \\
\texttt{seconds} & Segundos (por defecto cfg.HOLD\_VALID\_S). \\
\end{tabular}}\par\smallskip
\noindent\textbf{Devuelve:} True si terminó; False si se presionó A (parar).\par
\noindent\textbf{Llama a:} \hyperref[fn:hal:stop]{\texttt{hal.stop}}, \hyperref[fn:hal:now]{\texttt{hal.now}}, \hyperref[fn:hal:abort-requested]{\texttt{hal.abort\_\allowbreak{}requested}}, \hyperref[fn:hal:show]{\texttt{hal.show}}, \hyperref[fn:hal:wait]{\texttt{hal.wait}}\par
\lstinputlisting[firstline=47,lastline=57,firstnumber=47]{../../software/controls/gripper.py}
\subsection*{Base (core)}
\subsection{\texttt{software/core/hal.py}}\label{mod:hal}
\textbf{Capa de hardware (HAL): ÚNICO lugar que llama a la API del mBot2.}

Usa cyberpi, mbot2, mbuild y gamepad (opcional). El resto del programa usa estas funciones con convenciones fijas:

\begin{itemize}\setlength\itemsep{0pt}
\item drive(izq, der): potencia -100..100, positivo = adelante en ambas orugas.
\item heading(): grados acumulados, positivo = giro a la DERECHA.
\item pitch(): grados, positivo = nariz ARRIBA.
\item line\_offset(): -100..100, negativo = la línea está a la IZQUIERDA del centro.
\item darkness(): por sonda, 0 = blanco calibrado, 1 = negro calibrado.
\end{itemize}
Referencia de la API: docs/api\_mbot2.md. En la PC (pruebas y simulador) se usan módulos falsos con la misma interfaz (tests/fakes).

\noindent Archivo fuente: \texttt{software/core/hal.py} (426 l\'ineas).

\subsubsection{\texttt{clamp(v, lo, hi)}}\label{fn:hal:clamp}
\noindent Limita un valor a un rango.

\par\noindent{\small\begin{tabular}{@{}p{0.2\linewidth}p{0.74\linewidth}@{}}
\textbf{Par\'ametro} & \textbf{Descripci\'on} \\
\texttt{v} & Valor a limitar. \\
\texttt{lo} & Mínimo. \\
\texttt{hi} & Máximo. \\
\end{tabular}}\par\smallskip
\noindent\textbf{Devuelve:} v recortado a [lo, hi].\par
\lstinputlisting[firstline=40,lastline=45,firstnumber=40]{../../software/core/hal.py}
\subsubsection{\texttt{init()}}\label{fn:hal:init}
\noindent Inicializa el hardware al arrancar un programa.

Detiene los motores, pone en cero el rumbo y el odómetro, enciende la luz blanca del Quad RGB (la usan las sondas) y apaga los LEDs del CyberPi. Borra el aviso de abortar.

\noindent\textbf{Llama a:} \hyperref[fn:hal:stop]{\texttt{stop}}, \hyperref[fn:hal:reset-heading]{\texttt{reset\_\allowbreak{}heading}}, \hyperref[fn:hal:reset-odometer]{\texttt{reset\_\allowbreak{}odometer}}\par
\lstinputlisting[firstline=51,lastline=58,firstnumber=51]{../../software/core/hal.py}
\subsubsection{\texttt{now()}}\label{fn:hal:now}
\noindent Tiempo desde que se encendió el CyberPi.

\noindent\textbf{Devuelve:} Segundos (float).\par
\lstinputlisting[firstline=64,lastline=65,firstnumber=64]{../../software/core/hal.py}
\subsubsection{\texttt{wait(s)}}\label{fn:hal:wait}
\noindent Espera sin dejar de actualizar la odometría.

Antes de dormir llama a las funciones registradas en \_on\_wait (p. ej. odometry.update), así la posición se sigue integrando durante las esperas.

\par\noindent{\small\begin{tabular}{@{}p{0.2\linewidth}p{0.74\linewidth}@{}}
\textbf{Par\'ametro} & \textbf{Descripci\'on} \\
\texttt{s} & Segundos a esperar. \\
\end{tabular}}\par\smallskip
\lstinputlisting[firstline=75,lastline=78,firstnumber=75]{../../software/core/hal.py}
\subsubsection{\texttt{\_\allowbreak{}tip\_\allowbreak{}guard(left, right)}}\label{fn:hal:-tip-guard}
\noindent Control contra volcadura en rampas y escalera.

Inclinado, un arranque brusco empuja el centro de masa detrás de la punta de las orugas y el robot se voltea hacia atrás: con pitch >= tip\_pitch la potencia COMÚN (el avance) solo puede subir tip\_slew \%/s; bajar siempre está permitido. La diferencia entre orugas (la dirección) queda libre para no frenar las correcciones de rumbo. También hace la asistencia de subida: si con potencia hacia adelante el robot casi no avanza (menos de slope\_min\_v cm/s) o resbala, suma potencia de a poco (hasta +60). No recorta la potencia en pendientes fuertes: en una rampa de 40\textdegree{} eso dejaba al robot sin fuerza a media subida (simulador). El otro control es el brazo (gripper.climb\_pose).

\par\noindent{\small\begin{tabular}{@{}p{0.2\linewidth}p{0.74\linewidth}@{}}
\textbf{Par\'ametro} & \textbf{Descripci\'on} \\
\texttt{left} & Potencia pedida a la oruga izquierda (-100..100). \\
\texttt{right} & Potencia pedida a la oruga derecha (-100..100). \\
\end{tabular}}\par\smallskip
\noindent\textbf{Devuelve:} (izquierda, derecha) con el límite aplicado.\par
\noindent\textbf{Llama a:} \hyperref[fn:hal:now]{\texttt{now}}, \hyperref[fn:hal:odo-abs]{\texttt{odo\_\allowbreak{}abs}}, \hyperref[fn:hal:pitch]{\texttt{pitch}}, \hyperref[fn:hal:clamp]{\texttt{clamp}}\par
\lstinputlisting[firstline=100,lastline=131,firstnumber=100]{../../software/core/hal.py}
\subsubsection{\texttt{drive(left, right)}}\label{fn:hal:drive}
\noindent Mueve las orugas con potencia (convención del programa: + = adelante).

Pasa por \_tip\_guard y, si una oruga pasaría de 100 \%, baja las dos lo mismo para conservar la diferencia (la corrección de rumbo): recortar solo una dejaba casi sin corrección a potencia máxima y en la escalera una oruga se salía de su tramo (simulador). Aplica los signos de montaje LEFT\_SIGN / RIGHT\_SIGN.

\par\noindent{\small\begin{tabular}{@{}p{0.2\linewidth}p{0.74\linewidth}@{}}
\textbf{Par\'ametro} & \textbf{Descripci\'on} \\
\texttt{left} & Potencia de la oruga izquierda, -100..100. \\
\texttt{right} & Potencia de la oruga derecha, -100..100. \\
\end{tabular}}\par\smallskip
\noindent\textbf{Llama a:} \hyperref[fn:hal:-tip-guard]{\texttt{\_\allowbreak{}tip\_\allowbreak{}guard}}, \hyperref[fn:hal:clamp]{\texttt{clamp}}\par
\lstinputlisting[firstline=141,lastline=159,firstnumber=141]{../../software/core/hal.py}
\subsubsection{\texttt{stop()}}\label{fn:hal:stop}
\noindent Detiene las dos orugas (potencia 0).

\lstinputlisting[firstline=163,lastline=166,firstnumber=163]{../../software/core/hal.py}
\subsubsection{\texttt{\_\allowbreak{}encoder\_\allowbreak{}deg()}}\label{fn:hal:-encoder-deg}
\noindent Avance del robot en grados de encoder, con signo.

Promedio de las dos orugas con su signo: un giro en el lugar no suma (antes se usaba abs() y los giros cambiaban la odometría, lo que movía las distancias medidas desde un cruce).

\noindent\textbf{Devuelve:} Grados (adelante +, atrás -).\par
\lstinputlisting[firstline=173,lastline=176,firstnumber=173]{../../software/core/hal.py}
\subsubsection{\texttt{tracks\_\allowbreak{}cm()}}\label{fn:hal:tracks-cm}
\noindent Avance de cada oruga en cm desde el arranque, con signo.

\noindent\textbf{Devuelve:} (izquierda, derecha) en cm; usa los cm por grado de cada oruga (lr\_ratio).\par
\noindent\textbf{Llama a:} \hyperref[fn:geometry:cm-per-deg-lr]{\texttt{geometry.cm\_\allowbreak{}per\_\allowbreak{}deg\_\allowbreak{}lr}}\par
\lstinputlisting[firstline=181,lastline=184,firstnumber=181]{../../software/core/hal.py}
\subsubsection{\texttt{reset\_\allowbreak{}odometer()}}\label{fn:hal:reset-odometer}
\noindent Pone en cero el odómetro relativo (odometer\_cm).

\noindent\textbf{Llama a:} \hyperref[fn:hal:-encoder-deg]{\texttt{\_\allowbreak{}encoder\_\allowbreak{}deg}}\par
\lstinputlisting[firstline=188,lastline=190,firstnumber=188]{../../software/core/hal.py}
\subsubsection{\texttt{odo\_\allowbreak{}abs()}}\label{fn:hal:odo-abs}
\noindent Avance absoluto desde el arranque (no se reinicia).

\noindent\textbf{Devuelve:} cm con signo.\par
\noindent\textbf{Llama a:} \hyperref[fn:hal:-encoder-deg]{\texttt{\_\allowbreak{}encoder\_\allowbreak{}deg}}, \hyperref[fn:geometry:cm-per-deg]{\texttt{geometry.cm\_\allowbreak{}per\_\allowbreak{}deg}}\par
\lstinputlisting[firstline=195,lastline=196,firstnumber=195]{../../software/core/hal.py}
\subsubsection{\texttt{odometer\_\allowbreak{}cm()}}\label{fn:hal:odometer-cm}
\noindent Avance desde el último reset\_odometer().

\noindent\textbf{Devuelve:} cm con signo.\par
\noindent\textbf{Llama a:} \hyperref[fn:geometry:cm-per-deg]{\texttt{geometry.cm\_\allowbreak{}per\_\allowbreak{}deg}}, \hyperref[fn:hal:-encoder-deg]{\texttt{\_\allowbreak{}encoder\_\allowbreak{}deg}}\par
\lstinputlisting[firstline=201,lastline=202,firstnumber=201]{../../software/core/hal.py}
\subsubsection{\texttt{servo(port, angle)}}\label{fn:hal:servo}
\noindent Mueve un servo.

\par\noindent{\small\begin{tabular}{@{}p{0.2\linewidth}p{0.74\linewidth}@{}}
\textbf{Par\'ametro} & \textbf{Descripci\'on} \\
\texttt{port} & Puerto del servo ("s1" garra, "s2" brazo; ver cfg.SERVO\_CLAW / SERVO\_ARM). \\
\texttt{angle} & Ángulo en grados; se limita a 0..180. \\
\end{tabular}}\par\smallskip
\noindent\textbf{Llama a:} \hyperref[fn:hal:clamp]{\texttt{clamp}}\par
\lstinputlisting[firstline=208,lastline=209,firstnumber=208]{../../software/core/hal.py}
\subsubsection{\texttt{\_\allowbreak{}raw\_\allowbreak{}heading()}}\label{fn:hal:-raw-heading}
\noindent Rumbo del giroscopio sin el cero del programa.

Aplica HEADING\_SIGN y gyro\_scale (grados reales por grado leído; TODO calibrar: girar el robot A MANO 360\textdegree{} contra una escuadra y poner 360 / lectura).

\noindent\textbf{Devuelve:} Grados acumulados, + = derecha.\par
\lstinputlisting[firstline=217,lastline=220,firstnumber=217]{../../software/core/hal.py}
\subsubsection{\texttt{reset\_\allowbreak{}heading()}}\label{fn:hal:reset-heading}
\noindent Toma el rumbo actual como 0.

\noindent\textbf{Llama a:} \hyperref[fn:hal:-raw-heading]{\texttt{\_\allowbreak{}raw\_\allowbreak{}heading}}\par
\lstinputlisting[firstline=224,lastline=226,firstnumber=224]{../../software/core/hal.py}
\subsubsection{\texttt{heading()}}\label{fn:hal:heading}
\noindent Rumbo del robot.

\noindent\textbf{Devuelve:} Grados acumulados desde reset\_heading(); positivo = giro a la DERECHA.\par
\noindent\textbf{Llama a:} \hyperref[fn:hal:-raw-heading]{\texttt{\_\allowbreak{}raw\_\allowbreak{}heading}}\par
\lstinputlisting[firstline=231,lastline=232,firstnumber=231]{../../software/core/hal.py}
\subsubsection{\texttt{pitch()}}\label{fn:hal:pitch}
\noindent Inclinación del robot (cabeceo).

\noindent\textbf{Devuelve:} Grados; positivo = nariz ARRIBA (aplica PITCH\_SIGN).\par
\lstinputlisting[firstline=237,lastline=238,firstnumber=237]{../../software/core/hal.py}
\subsubsection{\texttt{probe\_\allowbreak{}light(probe)}}\label{fn:hal:probe-light}
\noindent Luz reflejada que lee una sonda del Quad RGB.

\par\noindent{\small\begin{tabular}{@{}p{0.2\linewidth}p{0.74\linewidth}@{}}
\textbf{Par\'ametro} & \textbf{Descripci\'on} \\
\texttt{probe} & Nombre de la sonda: "L2", "L1", "R1" o "R2". \\
\end{tabular}}\par\smallskip
\noindent\textbf{Devuelve:} Lectura cruda de luz.\par
\lstinputlisting[firstline=249,lastline=250,firstnumber=249]{../../software/core/hal.py}
\subsubsection{\texttt{darkness()}}\label{fn:hal:darkness}
\noindent Oscuridad de cada sonda según la calibración LINE\_CAL.

Con línea delgada (10-30 mm) la línea puede quedar ENTRE dos sondas; una sonda medio cubierta da \textasciitilde{}0.5 y el centro ponderado (offset\_from) sigue siendo correcto.

\noindent\textbf{Devuelve:} Dict sonda -> 0..1 (0 = blanco calibrado, 1 = negro calibrado).\par
\noindent\textbf{Llama a:} \hyperref[fn:hal:clamp]{\texttt{clamp}}, \hyperref[fn:hal:probe-light]{\texttt{probe\_\allowbreak{}light}}\par
\lstinputlisting[firstline=257,lastline=265,firstnumber=257]{../../software/core/hal.py}
\subsubsection{\texttt{offset\_\allowbreak{}from(dk)}}\label{fn:hal:offset-from}
\noindent Posición de la línea a partir de la oscuridad de las sondas (función pura).

Centro ponderado usando la posición real de cada sonda en cm (no están igual de separadas: 2.4 / 1.25 / 2.4). Antes resta el fondo (la sonda más clara): sobre una zona oscura (azul, verde) todas las sondas suben igual y la línea se perdía en el promedio.

\par\noindent{\small\begin{tabular}{@{}p{0.2\linewidth}p{0.74\linewidth}@{}}
\textbf{Par\'ametro} & \textbf{Descripci\'on} \\
\texttt{dk} & Dict sonda -> oscuridad 0..1 (de darkness()). \\
\end{tabular}}\par\smallskip
\noindent\textbf{Devuelve:} -100..100; negativo = línea a la izquierda; 0 si no hay línea.\par
\noindent\textbf{Llama a:} \hyperref[fn:hal:clamp]{\texttt{clamp}}\par
\lstinputlisting[firstline=274,lastline=291,firstnumber=274]{../../software/core/hal.py}
\subsubsection{\texttt{line\_\allowbreak{}offset()}}\label{fn:hal:line-offset}
\noindent Posición de la línea bajo el sensor.

line\_mode 1: analógico calibrado (offset\_from); 0: get\_offset\_track del sensor.

\noindent\textbf{Devuelve:} -100..100; negativo = la línea está a la IZQUIERDA.\par
\noindent\textbf{Llama a:} \hyperref[fn:hal:offset-from]{\texttt{offset\_\allowbreak{}from}}, \hyperref[fn:hal:darkness]{\texttt{darkness}}\par
\lstinputlisting[firstline=297,lastline=300,firstnumber=297]{../../software/core/hal.py}
\subsubsection{\texttt{line\_\allowbreak{}read()}}\label{fn:hal:line-read}
\noindent Posición de la línea y sondas sobre negro con UNA sola lectura.

Para el lazo de control. Además de L2..R2 devuelve "L2x"/"R2x": las exteriores con umbral cross\_th (negro "fuerte"), porque las marcas verdes junto a la línea también son oscuras y adelantaban el cruce (simulador).

\noindent\textbf{Devuelve:} (offset, bits): offset -100..100 y dict sonda -> True si ve línea.\par
\noindent\textbf{Llama a:} \hyperref[fn:hal:darkness]{\texttt{darkness}}, \hyperref[fn:hal:line-offset]{\texttt{line\_\allowbreak{}offset}}, \hyperref[fn:hal:line-bits]{\texttt{line\_\allowbreak{}bits}}, \hyperref[fn:hal:offset-from]{\texttt{offset\_\allowbreak{}from}}\par
\lstinputlisting[firstline=308,lastline=320,firstnumber=308]{../../software/core/hal.py}
\subsubsection{\texttt{line\_\allowbreak{}bits()}}\label{fn:hal:line-bits}
\noindent Qué sondas ven línea negra.

\noindent\textbf{Devuelve:} Dict sonda -> True si ve línea (umbral dark\_th, o bits de get\_line\_sta en line\_mode 0).\par
\noindent\textbf{Llama a:} \hyperref[fn:hal:darkness]{\texttt{darkness}}\par
\lstinputlisting[firstline=325,lastline=336,firstnumber=325]{../../software/core/hal.py}
\subsubsection{\texttt{probe\_\allowbreak{}rgb(probe)}}\label{fn:hal:probe-rgb}
\noindent Color RGB que lee una sonda.

\par\noindent{\small\begin{tabular}{@{}p{0.2\linewidth}p{0.74\linewidth}@{}}
\textbf{Par\'ametro} & \textbf{Descripci\'on} \\
\texttt{probe} & Nombre de la sonda. \\
\end{tabular}}\par\smallskip
\noindent\textbf{Devuelve:} (rojo, verde, azul).\par
\lstinputlisting[firstline=342,lastline=344,firstnumber=342]{../../software/core/hal.py}
\subsubsection{\texttt{probe\_\allowbreak{}color\_\allowbreak{}builtin(probe)}}\label{fn:hal:probe-color-builtin}
\noindent Nombre del color que reconoce el propio sensor.

\par\noindent{\small\begin{tabular}{@{}p{0.2\linewidth}p{0.74\linewidth}@{}}
\textbf{Par\'ametro} & \textbf{Descripci\'on} \\
\texttt{probe} & Nombre de la sonda. \\
\end{tabular}}\par\smallskip
\noindent\textbf{Devuelve:} Nombre del color según get\_color\_sta (color\_mode 0).\par
\lstinputlisting[firstline=350,lastline=351,firstnumber=350]{../../software/core/hal.py}
\subsubsection{\texttt{distance\_\allowbreak{}cm()}}\label{fn:hal:distance-cm}
\noindent Distancia del Ultrasónico 2.

\noindent\textbf{Devuelve:} cm, o -1 si el robot no lo lleva (cfg.HAS\_US = False, diseño actual).\par
\lstinputlisting[firstline=357,lastline=360,firstnumber=357]{../../software/core/hal.py}
\subsubsection{\texttt{button(name)}}\label{fn:hal:button}
\noindent Lee un botón del CyberPi.

\par\noindent{\small\begin{tabular}{@{}p{0.2\linewidth}p{0.74\linewidth}@{}}
\textbf{Par\'ametro} & \textbf{Descripci\'on} \\
\texttt{name} & "a", "b", "up", "down", "left", "right" o "middle". \\
\end{tabular}}\par\smallskip
\noindent\textbf{Devuelve:} True si está presionado.\par
\lstinputlisting[firstline=367,lastline=368,firstnumber=367]{../../software/core/hal.py}
\subsubsection{\texttt{abort\_\allowbreak{}requested()}}\label{fn:hal:abort-requested}
\noindent ¿Se pidió parar? (botón A, cfg.ABORT\_BUTTON).

Queda marcado hasta clear\_abort(): una pulsación corta basta para detener el autónomo o el control manual.

\noindent\textbf{Devuelve:} True si se presionó A desde el último clear\_abort().\par
\lstinputlisting[firstline=375,lastline=379,firstnumber=375]{../../software/core/hal.py}
\subsubsection{\texttt{clear\_\allowbreak{}abort()}}\label{fn:hal:clear-abort}
\noindent Borra el aviso de parar (después de atenderlo).

\lstinputlisting[firstline=383,lastline=385,firstnumber=383]{../../software/core/hal.py}
\subsubsection{\texttt{show(text, line, size)}}\label{fn:hal:show}
\noindent Escribe un renglón en la pantalla del CyberPi.

\par\noindent{\small\begin{tabular}{@{}p{0.2\linewidth}p{0.74\linewidth}@{}}
\textbf{Par\'ametro} & \textbf{Descripci\'on} \\
\texttt{text} & Texto a mostrar. \\
\texttt{line} & Renglón (0 arriba; 18 px por renglón). \\
\texttt{size} & Tamaño de letra. \\
\end{tabular}}\par\smallskip
\lstinputlisting[firstline=392,lastline=393,firstnumber=392]{../../software/core/hal.py}
\subsubsection{\texttt{leds(r, g, b)}}\label{fn:hal:leds}
\noindent Enciende los LEDs del CyberPi de un color.

\par\noindent{\small\begin{tabular}{@{}p{0.2\linewidth}p{0.74\linewidth}@{}}
\textbf{Par\'ametro} & \textbf{Descripci\'on} \\
\texttt{r} & Rojo 0..255. \\
\texttt{g} & Verde 0..255. \\
\texttt{b} & Azul 0..255. \\
\end{tabular}}\par\smallskip
\lstinputlisting[firstline=400,lastline=401,firstnumber=400]{../../software/core/hal.py}
\subsubsection{\texttt{pad\_\allowbreak{}ok()}}\label{fn:hal:pad-ok}
\noindent ¿Hay módulo de gamepad (Bluetooth Controller)?

\noindent\textbf{Devuelve:} True si se pudo importar gamepad.\par
\lstinputlisting[firstline=407,lastline=408,firstnumber=407]{../../software/core/hal.py}
\subsubsection{\texttt{pad\_\allowbreak{}key(name)}}\label{fn:hal:pad-key}
\noindent Lee un botón del gamepad.

\par\noindent{\small\begin{tabular}{@{}p{0.2\linewidth}p{0.74\linewidth}@{}}
\textbf{Par\'ametro} & \textbf{Descripci\'on} \\
\texttt{name} & Nombre del botón (p. ej. "N1", "R1", "Up"; ver docs/control\_manual.md). \\
\end{tabular}}\par\smallskip
\noindent\textbf{Devuelve:} True si está presionado (False si no hay gamepad).\par
\lstinputlisting[firstline=414,lastline=417,firstnumber=414]{../../software/core/hal.py}
\subsubsection{\texttt{pad\_\allowbreak{}axis(name)}}\label{fn:hal:pad-axis}
\noindent Lee un eje de los joysticks del gamepad.

\par\noindent{\small\begin{tabular}{@{}p{0.2\linewidth}p{0.74\linewidth}@{}}
\textbf{Par\'ametro} & \textbf{Descripci\'on} \\
\texttt{name} & "Lx", "Ly", "Rx" o "Ry". \\
\end{tabular}}\par\smallskip
\noindent\textbf{Devuelve:} Valor del eje (0 si no hay gamepad).\par
\lstinputlisting[firstline=423,lastline=426,firstnumber=423]{../../software/core/hal.py}
\subsection{\texttt{software/core/config.py}}\label{mod:config}
\textbf{Configuración del robot (mBot2 + orugas + garra de dos servos): puertos, signos y parámetros ajustables cfg.P.}

Dos tipos de valores:

\begin{itemize}\setlength\itemsep{0pt}
\item Constantes fijas (puertos, nombres de sondas): variables normales en MAYÚSCULAS.
\item Parámetros ajustables en pista: diccionario P. Se cambian desde el menú del CyberPi (software/core/settings.py), se guardan en el robot y se leen como cfg.P["clave"].
\end{itemize}
Todas las distancias van en cm y los ángulos en grados. Lo marcado MEDIR sale del robot real (ver docs/geometria.md); lo marcado CALIBRAR sale de tools/robot/api\_check.py. Las MEDIDAS DE LA PISTA (por nivel) están en software/core/pista.py.

\noindent Archivo fuente: \texttt{software/core/config.py} (256 l\'ineas).

\subsection{\texttt{software/core/pista.py}}\label{mod:pista}
\textbf{MEDIDAS DE LA PISTA que usa el programa AUTÓNOMO, por nivel.}

Aquí están TODAS las medidas de la pista de las que depende la ruta, separadas por nivel. Cada nivel tiene dos partes:

\begin{itemize}\setlength\itemsep{0pt}
\item "Mediciones que necesita el robot": se miden con flexómetro en la pista real.
\item "Ajustes de la maniobra": retrocesos, holguras y rumbos; se afinan probando en la pista.
\end{itemize}
Los valores de ahora salen del plano del reglamento (simulation/arena.py) y se probaron en el simulador. Los marcados  \# MEDIR  hay que confirmarlos en la pista; al medirlos, cambiar "MEDIR" por "MEDIDO".

Después de cambiar algo:

{\small\begin{alltt}
1) python3 -m pytest              (que todo siga pasando)
2) python3 simulation/run\_auto.py (ver la ruta en el simulador)
3) python3 tools/scripts/bundle\_for\_mblock.py --all   y volver a cargar mblock/rover\_x\_auto.py
\end{alltt}}
Las medidas marcadas [MENU] también se pueden cambiar desde el menú de ajustes del CyberPi (grupo "Pista"). OJO: si ya se guardaron desde el menú, en el robot manda lo guardado y no lo que diga este archivo (para volver a este archivo: menú -> "Restaurar valores").

UNIDADES Y MAPA

\begin{itemize}\setlength\itemsep{0pt}
\item Distancias en cm, ángulos en grados.
\item Posiciones (x, y) en el mapa de la vista superior del reglamento: x hacia la DERECHA desde la orilla izquierda de la arena, y hacia ABAJO desde la orilla de arriba.
\item Rumbos: 0 = derecha (este), 90 = abajo (sur), 180 = izquierda (oeste), -90 = arriba (norte).
\item "Centro de giro" = punto medio de las orugas; las sondas de línea van 5.2 cm delante de él.
\end{itemize}
Medidas del robot que importan para las holguras: 21.35 cm de ancho x 20.3 de largo; girando en el lugar las esquinas barren 16 cm (nariz) y 13.6 cm (cola).

\noindent Archivo fuente: \texttt{software/core/pista.py} (165 l\'ineas).

\subsection{\texttt{software/core/geometry.py}}\label{mod:geometry}
\textbf{Cuentas que dependen de la geometría del robot (cfg.P).}

Así, cambiar la geometría (otro chasis, otra garra, el sensor más adelante) solo requiere medir de nuevo y ajustar cfg.P; las misiones no cambian. Esquema en docs/geometria.md. Los parámetros "\_fwd" se miden desde el punto de referencia (sobre la caja del mBot2); aquí se convierten al CENTRO DE GIRO (punto medio de las orugas). Todo en cm.

\noindent Archivo fuente: \texttt{software/core/geometry.py} (234 l\'ineas).

\subsubsection{\texttt{p(key)}}\label{fn:geometry:p}
\noindent Atajo para leer un parámetro ajustable.

\par\noindent{\small\begin{tabular}{@{}p{0.2\linewidth}p{0.74\linewidth}@{}}
\textbf{Par\'ametro} & \textbf{Descripci\'on} \\
\texttt{key} & Clave de cfg.P. \\
\end{tabular}}\par\smallskip
\noindent\textbf{Devuelve:} cfg.P[key].\par
\lstinputlisting[firstline=14,lastline=15,firstnumber=14]{../../software/core/geometry.py}
\subsubsection{\texttt{fwd(key)}}\label{fn:geometry:fwd}
\noindent Convierte una distancia medida desde la referencia al centro de giro.

\par\noindent{\small\begin{tabular}{@{}p{0.2\linewidth}p{0.74\linewidth}@{}}
\textbf{Par\'ametro} & \textbf{Descripci\'on} \\
\texttt{key} & Clave de cfg.P con una distancia "\_fwd" (desde el punto de referencia). \\
\end{tabular}}\par\smallskip
\noindent\textbf{Devuelve:} cm hacia adelante desde el CENTRO DE GIRO.\par
\noindent\textbf{Llama a:} \hyperref[fn:geometry:p]{\texttt{p}}\par
\lstinputlisting[firstline=21,lastline=22,firstnumber=21]{../../software/core/geometry.py}
\subsubsection{\texttt{track\_\allowbreak{}w()}}\label{fn:geometry:track-w}
\noindent Distancia entre los centros de las dos orugas.

\noindent\textbf{Devuelve:} cm (robot\_w - belt\_w).\par
\noindent\textbf{Llama a:} \hyperref[fn:geometry:p]{\texttt{p}}\par
\lstinputlisting[firstline=27,lastline=28,firstnumber=27]{../../software/core/geometry.py}
\subsubsection{\texttt{contact\_\allowbreak{}half()}}\label{fn:geometry:contact-half}
\noindent Del centro de giro a cada extremo del tramo apoyado.

\noindent\textbf{Devuelve:} cm (eje de la rueda blanca / eje de la 1a negra = track\_c / 2).\par
\noindent\textbf{Llama a:} \hyperref[fn:geometry:p]{\texttt{p}}\par
\lstinputlisting[firstline=33,lastline=34,firstnumber=33]{../../software/core/geometry.py}
\subsubsection{\texttt{nose\_\allowbreak{}fwd()}}\label{fn:geometry:nose-fwd}
\noindent Del centro de giro al eje de la rueda de la nariz, en horizontal.

\noindent\textbf{Devuelve:} cm.\par
\noindent\textbf{Llama a:} \hyperref[fn:geometry:contact-half]{\texttt{contact\_\allowbreak{}half}}, \hyperref[fn:geometry:p]{\texttt{p}}\par
\lstinputlisting[firstline=39,lastline=40,firstnumber=39]{../../software/core/geometry.py}
\subsubsection{\texttt{nose\_\allowbreak{}h()}}\label{fn:geometry:nose-h}
\noindent Altura del eje de la nariz sobre los ejes apoyados.

\noindent\textbf{Devuelve:} cm.\par
\noindent\textbf{Llama a:} \hyperref[fn:geometry:p]{\texttt{p}}\par
\lstinputlisting[firstline=45,lastline=46,firstnumber=45]{../../software/core/geometry.py}
\subsubsection{\texttt{rear\_\allowbreak{}len()}}\label{fn:geometry:rear-len}
\noindent Del centro de giro a la punta trasera de las orugas.

\noindent\textbf{Devuelve:} cm.\par
\noindent\textbf{Llama a:} \hyperref[fn:geometry:p]{\texttt{p}}, \hyperref[fn:geometry:contact-half]{\texttt{contact\_\allowbreak{}half}}\par
\lstinputlisting[firstline=51,lastline=52,firstnumber=51]{../../software/core/geometry.py}
\subsubsection{\texttt{track\_\allowbreak{}front()}}\label{fn:geometry:track-front}
\noindent Del centro de giro a la punta delantera de las orugas (la nariz).

\noindent\textbf{Devuelve:} cm.\par
\noindent\textbf{Llama a:} \hyperref[fn:geometry:p]{\texttt{p}}, \hyperref[fn:geometry:nose-fwd]{\texttt{nose\_\allowbreak{}fwd}}\par
\lstinputlisting[firstline=57,lastline=58,firstnumber=57]{../../software/core/geometry.py}
\subsubsection{\texttt{track\_\allowbreak{}len()}}\label{fn:geometry:track-len}
\noindent Largo de la oruga de punta a punta, visto desde arriba.

\noindent\textbf{Devuelve:} cm.\par
\noindent\textbf{Llama a:} \hyperref[fn:geometry:rear-len]{\texttt{rear\_\allowbreak{}len}}, \hyperref[fn:geometry:track-front]{\texttt{track\_\allowbreak{}front}}\par
\lstinputlisting[firstline=63,lastline=64,firstnumber=63]{../../software/core/geometry.py}
\subsubsection{\texttt{front\_\allowbreak{}len()}}\label{fn:geometry:front-len}
\noindent Del centro de giro a lo más adelantado del robot.

Lo que esté más adelante: la nariz de las orugas o la punta de la garra en posición de viaje (h\_carry).

\noindent\textbf{Devuelve:} cm.\par
\noindent\textbf{Llama a:} \hyperref[fn:geometry:track-front]{\texttt{track\_\allowbreak{}front}}, \hyperref[fn:geometry:claw-reach]{\texttt{claw\_\allowbreak{}reach}}, \hyperref[fn:geometry:p]{\texttt{p}}\par
\lstinputlisting[firstline=71,lastline=74,firstnumber=71]{../../software/core/geometry.py}
\subsubsection{\texttt{robot\_\allowbreak{}len()}}\label{fn:geometry:robot-len}
\noindent Largo total del robot (cola a lo más adelantado).

\noindent\textbf{Devuelve:} cm.\par
\noindent\textbf{Llama a:} \hyperref[fn:geometry:rear-len]{\texttt{rear\_\allowbreak{}len}}, \hyperref[fn:geometry:front-len]{\texttt{front\_\allowbreak{}len}}\par
\lstinputlisting[firstline=79,lastline=80,firstnumber=79]{../../software/core/geometry.py}
\subsubsection{\texttt{cm\_\allowbreak{}per\_\allowbreak{}deg()}}\label{fn:geometry:cm-per-deg}
\noindent Avance de la oruga por grado de encoder.

Usa el diámetro efectivo sprocket\_d y el factor de patinaje slip.

\noindent\textbf{Devuelve:} cm por grado (promedio de las dos orugas).\par
\noindent\textbf{Llama a:} \hyperref[fn:geometry:p]{\texttt{p}}\par
\lstinputlisting[firstline=86,lastline=87,firstnumber=86]{../../software/core/geometry.py}
\subsubsection{\texttt{cm\_\allowbreak{}per\_\allowbreak{}deg\_\allowbreak{}lr()}}\label{fn:geometry:cm-per-deg-lr}
\noindent cm por grado de cada oruga.

Las orugas no avanzan igual por grado (prueba de 1 m: 63.5 y 65 cm, lr\_ratio); su promedio es cm\_per\_deg().

\noindent\textbf{Devuelve:} (izquierda, derecha) en cm por grado.\par
\noindent\textbf{Llama a:} \hyperref[fn:geometry:cm-per-deg]{\texttt{cm\_\allowbreak{}per\_\allowbreak{}deg}}, \hyperref[fn:geometry:p]{\texttt{p}}\par
\lstinputlisting[firstline=94,lastline=97,firstnumber=94]{../../software/core/geometry.py}
\subsubsection{\texttt{\_\allowbreak{}arm\_\allowbreak{}rows()}}\label{fn:geometry:-arm-rows}
\noindent Tabla medida del brazo.

\noindent\textbf{Devuelve:} ((ángulo, altura, adelanto), ...) para el servo a 0, 45 y 90 grados.\par
\noindent\textbf{Llama a:} \hyperref[fn:geometry:p]{\texttt{p}}\par
\lstinputlisting[firstline=106,lastline=108,firstnumber=106]{../../software/core/geometry.py}
\subsubsection{\texttt{arm\_\allowbreak{}at(servo)}}\label{fn:geometry:arm-at}
\noindent Posición de la garra para un ángulo del servo del brazo.

Interpola la tabla medida (el brazo no es un paralelogramo: la garra gira al subir).

\par\noindent{\small\begin{tabular}{@{}p{0.2\linewidth}p{0.74\linewidth}@{}}
\textbf{Par\'ametro} & \textbf{Descripci\'on} \\
\texttt{servo} & Ángulo del servo del brazo en grados. \\
\end{tabular}}\par\smallskip
\noindent\textbf{Devuelve:} (altura del centro de la garra sobre el piso, cm delante de la nariz).\par
\noindent\textbf{Llama a:} \hyperref[fn:geometry:-arm-rows]{\texttt{\_\allowbreak{}arm\_\allowbreak{}rows}}\par
\lstinputlisting[firstline=115,lastline=125,firstnumber=115]{../../software/core/geometry.py}
\subsubsection{\texttt{arm\_\allowbreak{}servo(h)}}\label{fn:geometry:arm-servo}
\noindent Ángulo del servo del brazo para una altura de la garra.

\par\noindent{\small\begin{tabular}{@{}p{0.2\linewidth}p{0.74\linewidth}@{}}
\textbf{Par\'ametro} & \textbf{Descripci\'on} \\
\texttt{h} & Altura deseada del centro de la garra sobre el piso (cm). \\
\end{tabular}}\par\smallskip
\noindent\textbf{Devuelve:} Grados del servo; fuera del alcance, el tope (0 o 90).\par
\noindent\textbf{Llama a:} \hyperref[fn:geometry:-arm-rows]{\texttt{\_\allowbreak{}arm\_\allowbreak{}rows}}\par
\lstinputlisting[firstline=131,lastline=140,firstnumber=131]{../../software/core/geometry.py}
\subsubsection{\texttt{claw\_\allowbreak{}reach(h)}}\label{fn:geometry:claw-reach}
\noindent Distancia horizontal del centro de giro al centro de la garra.

\par\noindent{\small\begin{tabular}{@{}p{0.2\linewidth}p{0.74\linewidth}@{}}
\textbf{Par\'ametro} & \textbf{Descripci\'on} \\
\texttt{h} & Altura de la garra (cm). \\
\end{tabular}}\par\smallskip
\noindent\textbf{Devuelve:} cm.\par
\noindent\textbf{Llama a:} \hyperref[fn:geometry:track-front]{\texttt{track\_\allowbreak{}front}}, \hyperref[fn:geometry:arm-at]{\texttt{arm\_\allowbreak{}at}}, \hyperref[fn:geometry:arm-servo]{\texttt{arm\_\allowbreak{}servo}}\par
\lstinputlisting[firstline=146,lastline=147,firstnumber=146]{../../software/core/geometry.py}
\subsubsection{\texttt{turn\_\allowbreak{}advance()}}\label{fn:geometry:turn-advance}
\noindent Avance tras detectar un cruce con las sondas para girar SOBRE él.

\noindent\textbf{Devuelve:} cm (las sondas van sensor\_fwd delante del centro de giro).\par
\noindent\textbf{Llama a:} \hyperref[fn:geometry:fwd]{\texttt{fwd}}\par
\lstinputlisting[firstline=153,lastline=154,firstnumber=153]{../../software/core/geometry.py}
\subsubsection{\texttt{reach\_\allowbreak{}advance(target, h)}}\label{fn:geometry:reach-advance}
\noindent Avance para que la garra llegue a un objeto.

\par\noindent{\small\begin{tabular}{@{}p{0.2\linewidth}p{0.74\linewidth}@{}}
\textbf{Par\'ametro} & \textbf{Descripci\'on} \\
\texttt{target} & cm del centro de giro al centro del objeto, en la dirección en la que mira el robot. \\
\texttt{h} & Altura de la garra (cm). \\
\end{tabular}}\par\smallskip
\noindent\textbf{Devuelve:} cm a avanzar (negativo = retroceder).\par
\noindent\textbf{Llama a:} \hyperref[fn:geometry:claw-reach]{\texttt{claw\_\allowbreak{}reach}}\par
\lstinputlisting[firstline=161,lastline=162,firstnumber=161]{../../software/core/geometry.py}
\subsubsection{\texttt{pedestal\_\allowbreak{}advance(h, from\_\allowbreak{}cross, traveled)}}\label{fn:geometry:pedestal-advance}
\noindent Avance para dejar la garra sobre el centro de un pedestal (sin ultrasónico).

El pedestal está a from\_cross cm de un cruce; cuando las sondas detectaron el cruce, el centro de giro estaba sensor\_fwd atrás.

\par\noindent{\small\begin{tabular}{@{}p{0.2\linewidth}p{0.74\linewidth}@{}}
\textbf{Par\'ametro} & \textbf{Descripci\'on} \\
\texttt{h} & Altura de la garra (cm). \\
\texttt{from\_\allowbreak{}cross} & cm del cruce al centro del pedestal (medida de la pista, p. ej. pista.nivel\_1\_cruce\_a\_pedestal\_cm). \\
\texttt{traveled} & cm que ya avanzó el robot desde que las sondas vieron el cruce. \\
\end{tabular}}\par\smallskip
\noindent\textbf{Devuelve:} cm a avanzar.\par
\noindent\textbf{Llama a:} \hyperref[fn:geometry:fwd]{\texttt{fwd}}, \hyperref[fn:geometry:claw-reach]{\texttt{claw\_\allowbreak{}reach}}\par
\lstinputlisting[firstline=173,lastline=174,firstnumber=173]{../../software/core/geometry.py}
\subsubsection{\texttt{tower\_\allowbreak{}advance(traveled)}}\label{fn:geometry:tower-advance}
\noindent Avance siguiendo la línea hasta quedar a side\_margin de la torre.

\par\noindent{\small\begin{tabular}{@{}p{0.2\linewidth}p{0.74\linewidth}@{}}
\textbf{Par\'ametro} & \textbf{Descripci\'on} \\
\texttt{traveled} & cm que ya avanzó desde que las sondas salieron de la zona verde. \\
\end{tabular}}\par\smallskip
\noindent\textbf{Devuelve:} cm a avanzar.\par
\noindent\textbf{Llama a:} \hyperref[fn:geometry:fwd]{\texttt{fwd}}, \hyperref[fn:geometry:p]{\texttt{p}}, \hyperref[fn:geometry:front-len]{\texttt{front\_\allowbreak{}len}}\par
\lstinputlisting[firstline=180,lastline=181,firstnumber=180]{../../software/core/geometry.py}
\subsubsection{\texttt{ramp\_\allowbreak{}extra()}}\label{fn:geometry:ramp-extra}
\noindent Avance después de quedar plano arriba de una rampa para estar completamente dentro.

\noindent\textbf{Devuelve:} cm.\par
\noindent\textbf{Llama a:} \hyperref[fn:geometry:rear-len]{\texttt{rear\_\allowbreak{}len}}, \hyperref[fn:geometry:p]{\texttt{p}}\par
\lstinputlisting[firstline=186,lastline=187,firstnumber=186]{../../software/core/geometry.py}
\subsubsection{\texttt{zone\_\allowbreak{}extra()}}\label{fn:geometry:zone-extra}
\noindent Avance desde que las sondas entran a una zona hasta que todo el robot está dentro.

\noindent\textbf{Devuelve:} cm.\par
\noindent\textbf{Llama a:} \hyperref[fn:geometry:p]{\texttt{p}}, \hyperref[fn:geometry:fwd]{\texttt{fwd}}, \hyperref[fn:geometry:rear-len]{\texttt{rear\_\allowbreak{}len}}\par
\lstinputlisting[firstline=192,lastline=193,firstnumber=192]{../../software/core/geometry.py}
\subsubsection{\texttt{detour\_\allowbreak{}out()}}\label{fn:geometry:detour-out}
\noindent Desplazamiento lateral para rodear la torre sin tocarla.

\noindent\textbf{Devuelve:} cm.\par
\noindent\textbf{Llama a:} \hyperref[fn:geometry:p]{\texttt{p}}\par
\lstinputlisting[firstline=198,lastline=199,firstnumber=198]{../../software/core/geometry.py}
\subsubsection{\texttt{detour\_\allowbreak{}pass()}}\label{fn:geometry:detour-pass}
\noindent Avance paralelo para rebasar la torre por completo.

Partiendo de side\_margin antes de ella.

\noindent\textbf{Devuelve:} cm.\par
\noindent\textbf{Llama a:} \hyperref[fn:geometry:rear-len]{\texttt{rear\_\allowbreak{}len}}, \hyperref[fn:geometry:p]{\texttt{p}}, \hyperref[fn:geometry:front-len]{\texttt{front\_\allowbreak{}len}}\par
\lstinputlisting[firstline=205,lastline=206,firstnumber=205]{../../software/core/geometry.py}
\subsubsection{\texttt{card\_\allowbreak{}to\_\allowbreak{}center(card\_\allowbreak{}odo, now\_\allowbreak{}odo)}}\label{fn:geometry:card-to-center}
\noindent Cuánto moverse para dejar el centro de giro frente a una tarjeta.

\par\noindent{\small\begin{tabular}{@{}p{0.2\linewidth}p{0.74\linewidth}@{}}
\textbf{Par\'ametro} & \textbf{Descripci\'on} \\
\texttt{card\_\allowbreak{}odo} & Odometría (cm) cuando la sonda detectó la tarjeta. \\
\texttt{now\_\allowbreak{}odo} & Odometría actual (cm). \\
\end{tabular}}\par\smallskip
\noindent\textbf{Devuelve:} cm a moverse (negativo = atrás).\par
\noindent\textbf{Llama a:} \hyperref[fn:geometry:fwd]{\texttt{fwd}}\par
\lstinputlisting[firstline=213,lastline=214,firstnumber=213]{../../software/core/geometry.py}
\subsubsection{\texttt{probe\_\allowbreak{}x(probe)}}\label{fn:geometry:probe-x}
\noindent Posición lateral de una sonda.

\par\noindent{\small\begin{tabular}{@{}p{0.2\linewidth}p{0.74\linewidth}@{}}
\textbf{Par\'ametro} & \textbf{Descripci\'on} \\
\texttt{probe} & Nombre de la sonda ("L2", "L1", "R1", "R2"). \\
\end{tabular}}\par\smallskip
\noindent\textbf{Devuelve:} cm, negativo = izquierda.\par
\noindent\textbf{Llama a:} \hyperref[fn:geometry:probe-lat]{\texttt{probe\_\allowbreak{}lat}}, \hyperref[fn:geometry:p]{\texttt{p}}\par
\lstinputlisting[firstline=220,lastline=222,firstnumber=220]{../../software/core/geometry.py}
\subsubsection{\texttt{probe\_\allowbreak{}lat(i, pitch, outer)}}\label{fn:geometry:probe-lat}
\noindent Posición lateral de la sonda i (función pura).

Las sondas no están igual de separadas: las dos centrales a pitch entre sí y cada exterior a outer de su central.

\par\noindent{\small\begin{tabular}{@{}p{0.2\linewidth}p{0.74\linewidth}@{}}
\textbf{Par\'ametro} & \textbf{Descripci\'on} \\
\texttt{i} & Índice 0..3, de izquierda a derecha. \\
\texttt{pitch} & cm entre las dos sondas centrales. \\
\texttt{outer} & cm de cada central a la exterior. \\
\end{tabular}}\par\smallskip
\noindent\textbf{Devuelve:} cm laterales.\par
\lstinputlisting[firstline=232,lastline=234,firstnumber=232]{../../software/core/geometry.py}
\subsection{\texttt{software/core/colors.py}}\label{mod:colors}
\textbf{Reconocimiento de colores: zonas de inicio, marcas verdes y tarjetas.}

Dos modos (cfg.P["color\_mode"]): 0 = el nombre que da el propio sensor (get\_color\_sta); 1 = el color de referencia RGB más cercano (cfg.COLOR\_REF).

\noindent Archivo fuente: \texttt{software/core/colors.py} (105 l\'ineas).

\subsubsection{\texttt{nearest(rgb, refs)}}\label{fn:colors:nearest}
\noindent Color de referencia más cercano a una lectura RGB.

Distancia sobre cromaticidad + brillo. Normalizar por el brillo hace la clasificación menos sensible a cambios de luz (reglamento: el robot debe adaptarse a la iluminación).

\par\noindent{\small\begin{tabular}{@{}p{0.2\linewidth}p{0.74\linewidth}@{}}
\textbf{Par\'ametro} & \textbf{Descripci\'on} \\
\texttt{rgb} & (rojo, verde, azul) leídos por la sonda. \\
\texttt{refs} & Dict nombre -> (r, g, b) de referencia (cfg.COLOR\_REF). \\
\end{tabular}}\par\smallskip
\noindent\textbf{Devuelve:} Nombre del color más cercano.\par
\lstinputlisting[firstline=28,lastline=42,firstnumber=28]{../../software/core/colors.py}
\subsubsection{\texttt{normalize(name)}}\label{fn:colors:normalize}
\noindent Nombre del sensor -> nombre del programa.

\par\noindent{\small\begin{tabular}{@{}p{0.2\linewidth}p{0.74\linewidth}@{}}
\textbf{Par\'ametro} & \textbf{Descripci\'on} \\
\texttt{name} & Nombre que devuelve el sensor (puede venir en mayúsculas o ser None). \\
\end{tabular}}\par\smallskip
\noindent\textbf{Devuelve:} Nombre en minúsculas, o "unknown".\par
\lstinputlisting[firstline=48,lastline=51,firstnumber=48]{../../software/core/colors.py}
\subsubsection{\texttt{read(probe)}}\label{fn:colors:read}
\noindent Color que ve una sonda (según color\_mode).

\par\noindent{\small\begin{tabular}{@{}p{0.2\linewidth}p{0.74\linewidth}@{}}
\textbf{Par\'ametro} & \textbf{Descripci\'on} \\
\texttt{probe} & Nombre de la sonda. \\
\end{tabular}}\par\smallskip
\noindent\textbf{Devuelve:} Nombre del color ("red", "blue", ...).\par
\noindent\textbf{Llama a:} \hyperref[fn:colors:normalize]{\texttt{normalize}}, \hyperref[fn:colors:nearest]{\texttt{nearest}}, \hyperref[fn:hal:probe-color-builtin]{\texttt{hal.probe\_\allowbreak{}color\_\allowbreak{}builtin}}, \hyperref[fn:hal:probe-rgb]{\texttt{hal.probe\_\allowbreak{}rgb}}\par
\lstinputlisting[firstline=57,lastline=60,firstnumber=57]{../../software/core/colors.py}
\subsubsection{\texttt{read\_\allowbreak{}stable(probe, n)}}\label{fn:colors:read-stable}
\noindent Lee hasta tener n lecturas iguales seguidas.

\par\noindent{\small\begin{tabular}{@{}p{0.2\linewidth}p{0.74\linewidth}@{}}
\textbf{Par\'ametro} & \textbf{Descripci\'on} \\
\texttt{probe} & Nombre de la sonda. \\
\texttt{n} & Lecturas iguales necesarias (por defecto cfg.COLOR\_CONFIRM); máx. 4n intentos. \\
\end{tabular}}\par\smallskip
\noindent\textbf{Devuelve:} El color confirmado, o la última lectura si no se confirmó.\par
\noindent\textbf{Llama a:} \hyperref[fn:colors:read]{\texttt{read}}, \hyperref[fn:hal:wait]{\texttt{hal.wait}}\par
\lstinputlisting[firstline=67,lastline=82,firstnumber=67]{../../software/core/colors.py}
\subsubsection{\texttt{show(name)}}\label{fn:colors:show}
\noindent Muestra un color en los LEDs y la pantalla (misión 11).

\par\noindent{\small\begin{tabular}{@{}p{0.2\linewidth}p{0.74\linewidth}@{}}
\textbf{Par\'ametro} & \textbf{Descripci\'on} \\
\texttt{name} & Nombre del color. \\
\end{tabular}}\par\smallskip
\noindent\textbf{Llama a:} \hyperref[fn:hal:leds]{\texttt{hal.leds}}, \hyperref[fn:hal:show]{\texttt{hal.show}}\par
\lstinputlisting[firstline=87,lastline=90,firstnumber=87]{../../software/core/colors.py}
\subsubsection{\texttt{zone\_\allowbreak{}level(color)}}\label{fn:colors:zone-level}
\noindent Nivel de inicio según el color de la zona.

\par\noindent{\small\begin{tabular}{@{}p{0.2\linewidth}p{0.74\linewidth}@{}}
\textbf{Par\'ametro} & \textbf{Descripci\'on} \\
\texttt{color} & Color de la zona en la que está el robot. \\
\end{tabular}}\par\smallskip
\noindent\textbf{Devuelve:} 1..4 (cfg.ZONE\_LEVEL), 0 = no se sabe.\par
\lstinputlisting[firstline=96,lastline=97,firstnumber=96]{../../software/core/colors.py}
\subsubsection{\texttt{is\_\allowbreak{}card\_\allowbreak{}color(name)}}\label{fn:colors:is-card-color}
\noindent ¿Es color de tarjeta?

Una tarjeta es cualquier color que no sea el piso blanco ni la línea negra.

\par\noindent{\small\begin{tabular}{@{}p{0.2\linewidth}p{0.74\linewidth}@{}}
\textbf{Par\'ametro} & \textbf{Descripci\'on} \\
\texttt{name} & Nombre del color. \\
\end{tabular}}\par\smallskip
\noindent\textbf{Devuelve:} True si puede ser una tarjeta.\par
\lstinputlisting[firstline=104,lastline=105,firstnumber=104]{../../software/core/colors.py}
\subsection{\texttt{software/core/settings.py}}\label{mod:settings}
\textbf{Menú de ajustes en el CyberPi, calibración de línea y guardado de parámetros.}

Se entra desde la pantalla de inicio del autónomo con el CENTRO del joystick (cfg.MENU\_BUTTON). Dentro del menú:

\begin{itemize}\setlength\itemsep{0pt}
\item arriba/abajo = elegir parámetro o acción
\item izquierda/derecha = bajar / subir el valor
\item A = probar (mueve el brazo o la garra a ese valor, o ejecuta la acción)
\item B = guardar y salir
\end{itemize}
Los valores se guardan en el robot (cfg.PARAMS\_FILE) y se cargan al encender; lo guardado manda sobre config.py y pista.py. Si el CyberPi no permite escribir archivos, los cambios duran hasta apagarlo: anotarlos en config.py o pista.py.

\noindent Archivo fuente: \texttt{software/core/settings.py} (245 l\'ineas).

\subsubsection{\texttt{remember\_\allowbreak{}defaults()}}\label{fn:settings:remember-defaults}
\noindent Guarda los valores de fábrica de cfg.P (los de config.py y pista.py).

Se llama al arrancar, antes de load(); la acción "Restaurar valores" vuelve a ellos.

\lstinputlisting[firstline=26,lastline=28,firstnumber=26]{../../software/core/settings.py}
\subsubsection{\texttt{snapshot()}}\label{fn:settings:snapshot}
\noindent Copia de los parámetros y la calibración de línea actuales.

\noindent\textbf{Devuelve:} Dict \{"P": \{...\}, "LINE\_CAL": \{...\}\} listo para guardar como JSON.\par
\lstinputlisting[firstline=33,lastline=40,firstnumber=33]{../../software/core/settings.py}
\subsubsection{\texttt{apply(data)}}\label{fn:settings:apply}
\noindent Aplica un diccionario guardado.

Ignora claves desconocidas (archivos de versiones viejas del programa).

\par\noindent{\small\begin{tabular}{@{}p{0.2\linewidth}p{0.74\linewidth}@{}}
\textbf{Par\'ametro} & \textbf{Descripci\'on} \\
\texttt{data} & Dict como el de snapshot(). \\
\end{tabular}}\par\smallskip
\lstinputlisting[firstline=46,lastline=53,firstnumber=46]{../../software/core/settings.py}
\subsubsection{\texttt{save()}}\label{fn:settings:save}
\noindent Guarda los parámetros en el robot (cfg.PARAMS\_FILE).

\noindent\textbf{Devuelve:} True si se pudo escribir el archivo.\par
\noindent\textbf{Llama a:} \hyperref[fn:settings:snapshot]{\texttt{snapshot}}\par
\lstinputlisting[firstline=58,lastline=65,firstnumber=58]{../../software/core/settings.py}
\subsubsection{\texttt{load()}}\label{fn:settings:load}
\noindent Carga los parámetros guardados en el robot.

\noindent\textbf{Devuelve:} True si había archivo y se aplicó; False = se usan los de fábrica.\par
\noindent\textbf{Llama a:} \hyperref[fn:settings:apply]{\texttt{apply}}\par
\lstinputlisting[firstline=70,lastline=78,firstnumber=70]{../../software/core/settings.py}
\subsubsection{\texttt{cal\_\allowbreak{}from\_\allowbreak{}samples(samples)}}\label{fn:settings:cal-from-samples}
\noindent Calibración de cada sonda a partir de lecturas de luz (función pura).

\par\noindent{\small\begin{tabular}{@{}p{0.2\linewidth}p{0.74\linewidth}@{}}
\textbf{Par\'ametro} & \textbf{Descripci\'on} \\
\texttt{samples} & Dict sonda -> lista de lecturas de luz tomadas pasando sobre la línea. \\
\end{tabular}}\par\smallskip
\noindent\textbf{Devuelve:} Dict sonda -> [blanco, negro], o None si alguna sonda no vio contraste (no pasó sobre la línea).\par
\lstinputlisting[firstline=86,lastline=95,firstnumber=86]{../../software/core/settings.py}
\subsubsection{\texttt{calibrate\_\allowbreak{}line()}}\label{fn:settings:calibrate-line}
\noindent Calibra las sondas de línea (acción "Calibrar linea").

Con el robot centrado sobre la línea, gira \textasciitilde{}45\textdegree{} a cada lado leyendo las 4 sondas y guarda en cfg.LINE\_CAL el blanco y el negro de cada una.

\noindent\textbf{Devuelve:} True si todas las sondas vieron contraste.\par
\noindent\textbf{Llama a:} \hyperref[fn:hal:show]{\texttt{hal.show}}, \hyperref[fn:hal:heading]{\texttt{hal.heading}}, \hyperref[fn:hal:stop]{\texttt{hal.stop}}, \hyperref[fn:settings:cal-from-samples]{\texttt{cal\_\allowbreak{}from\_\allowbreak{}samples}}, \hyperref[fn:hal:now]{\texttt{hal.now}}, \hyperref[fn:hal:drive]{\texttt{hal.drive}}, \hyperref[fn:hal:wait]{\texttt{hal.wait}}, \hyperref[fn:hal:probe-light]{\texttt{hal.probe\_\allowbreak{}light}}\par
\lstinputlisting[firstline=102,lastline=127,firstnumber=102]{../../software/core/settings.py}
\subsubsection{\texttt{show\_\allowbreak{}sensors()}}\label{fn:settings:show-sensors}
\noindent Muestra oscuridad por sonda, offset, pitch y rumbo hasta presionar B.

\noindent\textbf{Llama a:} \hyperref[fn:hal:wait]{\texttt{hal.wait}}, \hyperref[fn:hal:button]{\texttt{hal.button}}, \hyperref[fn:hal:darkness]{\texttt{hal.darkness}}, \hyperref[fn:hal:show]{\texttt{hal.show}}, \hyperref[fn:hal:offset-from]{\texttt{hal.offset\_\allowbreak{}from}}, \hyperref[fn:hal:pitch]{\texttt{hal.pitch}}, \hyperref[fn:hal:heading]{\texttt{hal.heading}}\par
\lstinputlisting[firstline=131,lastline=139,firstnumber=131]{../../software/core/settings.py}
\subsubsection{\texttt{run\_\allowbreak{}action(name)}}\label{fn:settings:run-action}
\noindent Ejecuta una acción del menú.

\par\noindent{\small\begin{tabular}{@{}p{0.2\linewidth}p{0.74\linewidth}@{}}
\textbf{Par\'ametro} & \textbf{Descripci\'on} \\
\texttt{name} & Nombre de la acción (ver ACTIONS). \\
\end{tabular}}\par\smallskip
\noindent\textbf{Devuelve:} True si la acción terminó bien.\par
\noindent\textbf{Llama a:} \hyperref[fn:settings:calibrate-line]{\texttt{calibrate\_\allowbreak{}line}}, \hyperref[fn:settings:show-sensors]{\texttt{show\_\allowbreak{}sensors}}, \hyperref[fn:motion:turn]{\texttt{motion.turn}}, \hyperref[fn:motion:straight]{\texttt{motion.straight}}\par
\lstinputlisting[firstline=145,lastline=159,firstnumber=145]{../../software/core/settings.py}
\subsubsection{\texttt{items()}}\label{fn:settings:items}
\noindent Lista de renglones del menú: primero los parámetros (cfg.TUNE), luego las acciones.

\noindent\textbf{Devuelve:} Lista de ("param", spec) o ("action", nombre).\par
\lstinputlisting[firstline=165,lastline=171,firstnumber=165]{../../software/core/settings.py}
\subsubsection{\texttt{bump(spec, direction)}}\label{fn:settings:bump}
\noindent Cambia un parámetro un paso.

\par\noindent{\small\begin{tabular}{@{}p{0.2\linewidth}p{0.74\linewidth}@{}}
\textbf{Par\'ametro} & \textbf{Descripci\'on} \\
\texttt{spec} & (grupo, clave, paso, mínimo, máximo), de cfg.TUNE. \\
\texttt{direction} & 1 = subir, -1 = bajar. \\
\end{tabular}}\par\smallskip
\noindent\textbf{Devuelve:} El valor nuevo (recortado al rango).\par
\noindent\textbf{Llama a:} \hyperref[fn:hal:clamp]{\texttt{hal.clamp}}\par
\lstinputlisting[firstline=178,lastline=187,firstnumber=178]{../../software/core/settings.py}
\subsubsection{\texttt{try\_\allowbreak{}param(spec)}}\label{fn:settings:try-param}
\noindent Prueba física de un parámetro (botón A).

Alturas "h\_": el brazo va a esa altura; aperturas "claw\_": la garra va a esa apertura; tabla del brazo "arm\_": el brazo va a la altura de transporte.

\par\noindent{\small\begin{tabular}{@{}p{0.2\linewidth}p{0.74\linewidth}@{}}
\textbf{Par\'ametro} & \textbf{Descripci\'on} \\
\texttt{spec} & Renglón de cfg.TUNE. \\
\end{tabular}}\par\smallskip
\noindent\textbf{Llama a:} \hyperref[fn:gripper:arm-to]{\texttt{gripper.arm\_\allowbreak{}to}}, \hyperref[fn:gripper:claw]{\texttt{gripper.claw}}\par
\lstinputlisting[firstline=194,lastline=201,firstnumber=194]{../../software/core/settings.py}
\subsubsection{\texttt{menu()}}\label{fn:settings:menu}
\noindent Bucle del menú de ajustes; sale al presionar B (guarda antes).

\noindent\textbf{Llama a:} \hyperref[fn:settings:items]{\texttt{items}}, \hyperref[fn:hal:wait]{\texttt{hal.wait}}, \hyperref[fn:hal:show]{\texttt{hal.show}}, \hyperref[fn:hal:button]{\texttt{hal.button}}, \hyperref[fn:settings:bump]{\texttt{bump}}, \hyperref[fn:settings:try-param]{\texttt{try\_\allowbreak{}param}}, \hyperref[fn:settings:run-action]{\texttt{run\_\allowbreak{}action}}, \hyperref[fn:settings:save]{\texttt{save}}, \hyperref[fn:hal:clear-abort]{\texttt{hal.clear\_\allowbreak{}abort}}\par
\lstinputlisting[firstline=205,lastline=244,firstnumber=205]{../../software/core/settings.py}
\section{Programas de prueba y calibraci\'on}
Programas sueltos (cada uno ya es un solo archivo; se suben directo desde mBlock) para calibrar el movimiento del robot.
\subsection{\texttt{software/robot/calibrations.py}}\label{mod:calibrations}
\textbf{Programa de prueba de calibración (movimiento lineal, giros, círculo y giroscopio). Archivo único: se sube directo desde mBlock.}

círculo y giroscopio). Se sube desde mBlock (Python, modo "subir") como el programa del equipo.

BOTÓN A: secuencia completa (en un espacio libre de \textasciitilde{}1 m x 1 m, con una marca de salida):

{\small\begin{alltt}
1. avanza 20 cm            2. gira 90\textdegree{} a la izquierda   3. avanza 10 cm
4. gira 90\textdegree{} a la derecha   5. gira 90\textdegree{} a la izquierda   6. avanza 20 cm
7. dibuja un círculo de 40 cm de diámetro (hacia la izquierda) y para
Al terminar cada paso la pantalla muestra lo MEDIDO por el robot (encoders y giroscopio).
Medir en el piso con regla/transportador y comparar (hoja abajo).
\end{alltt}}
JOYSTICK (pruebas sueltas):

{\small\begin{alltt}
arriba = avanza 50 cm   abajo = gira 360\textdegree{} en el lugar   izquierda = gira 90\textdegree{} izq
derecha = gira 90\textdegree{} der  centro = solo el círculo de 40 cm
\end{alltt}}
BOTÓN B: detener. Con el robot quieto: rumbo en 0 (prueba a mano del giroscopio). EN REPOSO la pantalla muestra el rumbo del giroscopio: PRUEBA A MANO -> poner el robot contra

{\small\begin{alltt}
una escuadra, B (rumbo 0), girarlo A MANO una vuelta completa hasta la misma escuadra y leer
la pantalla. GYRO\_SCALE = 360 / lectura (si marca \textasciitilde{}277, GYRO\_SCALE = 1.3).
\end{alltt}}
HOJA DE CALIBRACIÓN (anotar y pasar a software/core/config.py):

\begin{itemize}\setlength\itemsep{0pt}
\item avance:  real / pedido. Si avanza 19 en vez de 20 -> RUEDA\_D * 19/20  ("sprocket\_d")
\item recto:   si se desvía siempre al mismo lado -> LR ("lr\_ratio")
\item giros:   real en el piso vs 90\textdegree{}. Si el giro real se pasa o se queda corto y la pantalla dice 90 -> el giroscopio: GYRO\_SCALE ("gyro\_scale"); si la pantalla tampoco dice 90 -> ajustar KP\_GIRO / GIRO\_MIN
\item círculo: diámetro real (medir el rastro o la marca) vs 40. Si sale más chico/grande -> SKID ("skid": derrape de las orugas al girar)
\end{itemize}
\noindent Archivo fuente: \texttt{software/robot/calibrations.py} (314 l\'ineas).

\subsubsection{\texttt{motores(izq, der)}}\label{fn:calibrations:motores}
\noindent Velocidad de cada oruga en rpm (+ = adelante; EM2 va al revés).

\par\noindent{\small\begin{tabular}{@{}p{0.2\linewidth}p{0.74\linewidth}@{}}
\textbf{Par\'ametro} & \textbf{Descripci\'on} \\
\texttt{izq} & rpm de la oruga izquierda. \\
\texttt{der} & rpm de la oruga derecha. \\
\end{tabular}}\par\smallskip
\lstinputlisting[firstline=59,lastline=61,firstnumber=59]{../../software/robot/calibrations.py}
\subsubsection{\texttt{parar()}}\label{fn:calibrations:parar}
\noindent Detiene las dos orugas.

\lstinputlisting[firstline=65,lastline=67,firstnumber=65]{../../software/robot/calibrations.py}
\subsubsection{\texttt{rumbo()}}\label{fn:calibrations:rumbo}
\noindent Rumbo del giroscopio.

\noindent\textbf{Devuelve:} Grados acumulados, + = a la DERECHA (aplica GYRO\_SCALE y el cero de la prueba a mano).\par
\lstinputlisting[firstline=72,lastline=73,firstnumber=72]{../../software/robot/calibrations.py}
\subsubsection{\texttt{orugas\_\allowbreak{}cm()}}\label{fn:calibrations:orugas-cm}
\noindent Avance de cada oruga desde el arranque.

\noindent\textbf{Devuelve:} (izquierda, derecha) en cm, + = adelante.\par
\lstinputlisting[firstline=78,lastline=79,firstnumber=78]{../../software/robot/calibrations.py}
\subsubsection{\texttt{limitar(v, lo, hi)}}\label{fn:calibrations:limitar}
\noindent Limita un valor a un rango.

\par\noindent{\small\begin{tabular}{@{}p{0.2\linewidth}p{0.74\linewidth}@{}}
\textbf{Par\'ametro} & \textbf{Descripci\'on} \\
\texttt{v} & Valor. \\
\texttt{lo} & Mínimo. \\
\texttt{hi} & Máximo. \\
\end{tabular}}\par\smallskip
\noindent\textbf{Devuelve:} v recortado a [lo, hi].\par
\lstinputlisting[firstline=87,lastline=88,firstnumber=87]{../../software/robot/calibrations.py}
\subsubsection{\texttt{anotar(texto)}}\label{fn:calibrations:anotar}
\noindent Escribe una línea de resultados en la consola del CyberPi.

\par\noindent{\small\begin{tabular}{@{}p{0.2\linewidth}p{0.74\linewidth}@{}}
\textbf{Par\'ametro} & \textbf{Descripci\'on} \\
\texttt{texto} & Texto a escribir. \\
\end{tabular}}\par\smallskip
\lstinputlisting[firstline=93,lastline=94,firstnumber=93]{../../software/robot/calibrations.py}
\subsubsection{\texttt{avanzar\_\allowbreak{}cm(cm, rpm)}}\label{fn:calibrations:avanzar-cm}
\noindent Avanza una distancia por encoders sosteniendo el rumbo con el giroscopio.

\par\noindent{\small\begin{tabular}{@{}p{0.2\linewidth}p{0.74\linewidth}@{}}
\textbf{Par\'ametro} & \textbf{Descripci\'on} \\
\texttt{cm} & cm (negativo = atrás). \\
\texttt{rpm} & rpm de avance. \\
\end{tabular}}\par\smallskip
\noindent\textbf{Devuelve:} True si recorrió la distancia.\par
\noindent\textbf{Llama a:} \hyperref[fn:calibrations:rumbo]{\texttt{rumbo}}, \hyperref[fn:calibrations:orugas-cm]{\texttt{orugas\_\allowbreak{}cm}}, \hyperref[fn:calibrations:parar]{\texttt{parar}}, \hyperref[fn:calibrations:anotar]{\texttt{anotar}}, \hyperref[fn:calibrations:limitar]{\texttt{limitar}}, \hyperref[fn:calibrations:motores]{\texttt{motores}}\par
\lstinputlisting[firstline=102,lastline=131,firstnumber=102]{../../software/robot/calibrations.py}
\subsubsection{\texttt{girar(angulo, rpm\_\allowbreak{}max, timeout)}}\label{fn:calibrations:girar}
\noindent Gira en el lugar con el giroscopio (+ = derecha).

Frena al acercarse y termina corrigiendo hasta quedar a menos de TOL\_GIRO.

\par\noindent{\small\begin{tabular}{@{}p{0.2\linewidth}p{0.74\linewidth}@{}}
\textbf{Par\'ametro} & \textbf{Descripci\'on} \\
\texttt{angulo} & Grados a girar. \\
\texttt{rpm\_\allowbreak{}max} & rpm máximas. \\
\texttt{timeout} & s máximos. \\
\end{tabular}}\par\smallskip
\noindent\textbf{Devuelve:} True si llegó al ángulo.\par
\noindent\textbf{Llama a:} \hyperref[fn:calibrations:rumbo]{\texttt{rumbo}}, \hyperref[fn:calibrations:parar]{\texttt{parar}}, \hyperref[fn:calibrations:anotar]{\texttt{anotar}}, \hyperref[fn:calibrations:limitar]{\texttt{limitar}}, \hyperref[fn:calibrations:motores]{\texttt{motores}}\par
\lstinputlisting[firstline=140,lastline=168,firstnumber=140]{../../software/robot/calibrations.py}
\subsubsection{\texttt{circulo(diametro, rpm, lado)}}\label{fn:calibrations:circulo}
\noindent Una vuelta completa sobre un círculo (prueba de derrape SKID).

Velocidades de las orugas según el radio (con el derrape SKID); el giroscopio corrige: el rumbo debe avanzar 360\textdegree{} en pi\(\cdot\)diámetro cm.

\par\noindent{\small\begin{tabular}{@{}p{0.2\linewidth}p{0.74\linewidth}@{}}
\textbf{Par\'ametro} & \textbf{Descripci\'on} \\
\texttt{diametro} & Diámetro en cm (del centro del robot). \\
\texttt{rpm} & rpm promedio. \\
\texttt{lado} & -1 = hacia la izquierda, +1 = derecha. \\
\end{tabular}}\par\smallskip
\noindent\textbf{Devuelve:} True si completó la vuelta.\par
\noindent\textbf{Llama a:} \hyperref[fn:calibrations:rumbo]{\texttt{rumbo}}, \hyperref[fn:calibrations:orugas-cm]{\texttt{orugas\_\allowbreak{}cm}}, \hyperref[fn:calibrations:parar]{\texttt{parar}}, \hyperref[fn:calibrations:anotar]{\texttt{anotar}}, \hyperref[fn:calibrations:limitar]{\texttt{limitar}}, \hyperref[fn:calibrations:motores]{\texttt{motores}}\par
\lstinputlisting[firstline=178,lastline=214,firstnumber=178]{../../software/robot/calibrations.py}
\subsubsection{\texttt{secuencia()}}\label{fn:calibrations:secuencia}
\noindent Secuencia de prueba del botón A (ver la lista de pasos al inicio del archivo).

\noindent\textbf{Devuelve:} True si terminó todos los pasos.\par
\noindent\textbf{Llama a:} \hyperref[fn:calibrations:avanzar-cm]{\texttt{avanzar\_\allowbreak{}cm}}, \hyperref[fn:calibrations:girar]{\texttt{girar}}, \hyperref[fn:calibrations:circulo]{\texttt{circulo}}\par
\lstinputlisting[firstline=220,lastline=228,firstnumber=220]{../../software/robot/calibrations.py}
\subsubsection{\texttt{correr(nombre, accion)}}\label{fn:calibrations:correr}
\noindent Corre una prueba si no hay otra corriendo (la detiene B).

\par\noindent{\small\begin{tabular}{@{}p{0.2\linewidth}p{0.74\linewidth}@{}}
\textbf{Par\'ametro} & \textbf{Descripci\'on} \\
\texttt{nombre} & Nombre que se muestra en la pantalla. \\
\texttt{accion} & Función sin argumentos que hace la prueba. \\
\end{tabular}}\par\smallskip
\noindent\textbf{Llama a:} \hyperref[fn:calibrations:anotar]{\texttt{anotar}}, \hyperref[fn:calibrations:parar]{\texttt{parar}}\par
\lstinputlisting[firstline=234,lastline=246,firstnumber=234]{../../software/robot/calibrations.py}
\subsubsection{\texttt{on\_\allowbreak{}start()}}\label{fn:calibrations:on-start}
\noindent Al encender: muestra los botones de las pruebas.

\noindent\textbf{Llama a:} \hyperref[fn:calibrations:anotar]{\texttt{anotar}}\par
\lstinputlisting[firstline=250,lastline=255,firstnumber=250]{../../software/robot/calibrations.py}
\subsubsection{\texttt{prueba\_\allowbreak{}secuencia()}}\label{fn:calibrations:prueba-secuencia}
\noindent Botón A: secuencia completa de calibración.

\noindent\textbf{Llama a:} \hyperref[fn:calibrations:correr]{\texttt{correr}}\par
\lstinputlisting[firstline=259,lastline=261,firstnumber=259]{../../software/robot/calibrations.py}
\subsubsection{\texttt{prueba\_\allowbreak{}lineal()}}\label{fn:calibrations:prueba-lineal}
\noindent Joystick arriba: avanza 50 cm.

\noindent\textbf{Llama a:} \hyperref[fn:calibrations:correr]{\texttt{correr}}, \hyperref[fn:calibrations:avanzar-cm]{\texttt{avanzar\_\allowbreak{}cm}}\par
\lstinputlisting[firstline=265,lastline=267,firstnumber=265]{../../software/robot/calibrations.py}
\subsubsection{\texttt{prueba\_\allowbreak{}360()}}\label{fn:calibrations:prueba-360}
\noindent Joystick abajo: gira 360\textdegree{} en el lugar.

\noindent\textbf{Llama a:} \hyperref[fn:calibrations:correr]{\texttt{correr}}, \hyperref[fn:calibrations:girar]{\texttt{girar}}\par
\lstinputlisting[firstline=271,lastline=273,firstnumber=271]{../../software/robot/calibrations.py}
\subsubsection{\texttt{prueba\_\allowbreak{}izq()}}\label{fn:calibrations:prueba-izq}
\noindent Joystick izquierda: gira 90\textdegree{} a la izquierda.

\noindent\textbf{Llama a:} \hyperref[fn:calibrations:correr]{\texttt{correr}}, \hyperref[fn:calibrations:girar]{\texttt{girar}}\par
\lstinputlisting[firstline=277,lastline=279,firstnumber=277]{../../software/robot/calibrations.py}
\subsubsection{\texttt{prueba\_\allowbreak{}der()}}\label{fn:calibrations:prueba-der}
\noindent Joystick derecha: gira 90\textdegree{} a la derecha.

\noindent\textbf{Llama a:} \hyperref[fn:calibrations:correr]{\texttt{correr}}, \hyperref[fn:calibrations:girar]{\texttt{girar}}\par
\lstinputlisting[firstline=283,lastline=285,firstnumber=283]{../../software/robot/calibrations.py}
\subsubsection{\texttt{prueba\_\allowbreak{}circulo()}}\label{fn:calibrations:prueba-circulo}
\noindent Joystick centro: solo el círculo de 40 cm.

\noindent\textbf{Llama a:} \hyperref[fn:calibrations:correr]{\texttt{correr}}, \hyperref[fn:calibrations:circulo]{\texttt{circulo}}\par
\lstinputlisting[firstline=289,lastline=291,firstnumber=289]{../../software/robot/calibrations.py}
\subsubsection{\texttt{detener()}}\label{fn:calibrations:detener}
\noindent Botón B: detiene la prueba en curso; con el robot quieto pone el rumbo en 0.

\noindent\textbf{Llama a:} \hyperref[fn:calibrations:parar]{\texttt{parar}}, \hyperref[fn:calibrations:anotar]{\texttt{anotar}}\par
\lstinputlisting[firstline=295,lastline=304,firstnumber=295]{../../software/robot/calibrations.py}
\subsubsection{\texttt{main()}}\label{fn:calibrations:main}
\noindent Ciclo en reposo: muestra el rumbo del giroscopio en grande (prueba a mano).

\noindent\textbf{Llama a:} \hyperref[fn:calibrations:rumbo]{\texttt{rumbo}}\par
\lstinputlisting[firstline=308,lastline=314,firstnumber=308]{../../software/robot/calibrations.py}
\subsection{\texttt{software/robot/funciones\_\allowbreak{}rover.py}}\label{mod:funciones-rover}
\textbf{Funciones de movimiento preciso del Rover X (versión mejorada de FuncionesEnPython.mblock). Archivo único: se sube directo desde mBlock.}

Mismo uso que el programa original: se sube desde mBlock (modo "subir", Python) y se prueba con los botones. Lo que cambia:

\begin{itemize}\setlength\itemsep{0pt}
\item avanzar\_cm(): avanza por DISTANCIA (encoders) y no por tiempo, sosteniendo el rumbo con el giroscopio. Usa las medidas del robot: diámetro efectivo de la oruga 4.18 cm y la oruga izquierda avanza 2.3 \% menos por grado (prueba de 1 m: 63.5 / 65 cm).
\item girar(): en vez de parar al 77 \% del ángulo ("* 1.3 por error del sensor"), FRENA al acercarse y termina corrigiendo hasta quedar a menos de 1\textdegree{}. Si el error era la inercia, esto lo resuelve; si el giroscopio de verdad mide de menos, se corrige con GYRO\_SCALE (ver la prueba "giro a mano" abajo).
\item El rumbo no se reinicia en cada giro: los ángulos se suman sin acumular error.
\end{itemize}
Botones:

{\small\begin{alltt}
A          secuencia de prueba (la misma del programa original, ahora en cm)
B          detener; con el robot quieto, pone el rumbo en 0 (para la prueba a mano)
joystick   arriba 45\textdegree{} der, abajo 90\textdegree{} izq, izquierda 30\textdegree{} der, derecha 180\textdegree{}, centro 360\textdegree{}
\end{alltt}}
En reposo la pantalla muestra el rumbo: GIRAR EL ROBOT A MANO 360\textdegree{} (contra una escuadra o una marca en el piso) y anotar lo que marca -> GYRO\_SCALE = 360 / lectura.

\noindent Archivo fuente: \texttt{software/robot/funciones\_\allowbreak{}rover.py} (309 l\'ineas).

\subsubsection{\texttt{motores(izq, der)}}\label{fn:funciones_rover:motores}
\noindent Velocidad de cada oruga en rpm (+ = adelante; EM2 va al revés).

\par\noindent{\small\begin{tabular}{@{}p{0.2\linewidth}p{0.74\linewidth}@{}}
\textbf{Par\'ametro} & \textbf{Descripci\'on} \\
\texttt{izq} & rpm de la oruga izquierda. \\
\texttt{der} & rpm de la oruga derecha. \\
\end{tabular}}\par\smallskip
\lstinputlisting[firstline=46,lastline=48,firstnumber=46]{../../software/robot/funciones_rover.py}
\subsubsection{\texttt{parar()}}\label{fn:funciones_rover:parar}
\noindent Detiene las dos orugas.

\lstinputlisting[firstline=52,lastline=54,firstnumber=52]{../../software/robot/funciones_rover.py}
\subsubsection{\texttt{rumbo()}}\label{fn:funciones_rover:rumbo}
\noindent Rumbo del giroscopio.

\noindent\textbf{Devuelve:} Grados acumulados, + = a la DERECHA (aplica GYRO\_SCALE y el cero de la prueba a mano).\par
\lstinputlisting[firstline=59,lastline=60,firstnumber=59]{../../software/robot/funciones_rover.py}
\subsubsection{\texttt{orugas\_\allowbreak{}cm()}}\label{fn:funciones_rover:orugas-cm}
\noindent Avance de cada oruga desde el arranque.

\noindent\textbf{Devuelve:} (izquierda, derecha) en cm, + = adelante.\par
\lstinputlisting[firstline=65,lastline=66,firstnumber=65]{../../software/robot/funciones_rover.py}
\subsubsection{\texttt{limitar(v, lo, hi)}}\label{fn:funciones_rover:limitar}
\noindent Limita un valor a un rango.

\par\noindent{\small\begin{tabular}{@{}p{0.2\linewidth}p{0.74\linewidth}@{}}
\textbf{Par\'ametro} & \textbf{Descripci\'on} \\
\texttt{v} & Valor. \\
\texttt{lo} & Mínimo. \\
\texttt{hi} & Máximo. \\
\end{tabular}}\par\smallskip
\noindent\textbf{Devuelve:} v recortado a [lo, hi].\par
\lstinputlisting[firstline=74,lastline=75,firstnumber=74]{../../software/robot/funciones_rover.py}
\subsubsection{\texttt{mover\_\allowbreak{}motores(vel\_\allowbreak{}m1, vel\_\allowbreak{}m2, duracion)}}\label{fn:funciones_rover:mover-motores}
\noindent Movimiento por tiempo (el del programa original; se deja para pruebas manuales).

\par\noindent{\small\begin{tabular}{@{}p{0.2\linewidth}p{0.74\linewidth}@{}}
\textbf{Par\'ametro} & \textbf{Descripci\'on} \\
\texttt{vel\_\allowbreak{}m1} & rpm de EM1 (izquierda). \\
\texttt{vel\_\allowbreak{}m2} & rpm de EM2 (derecha). \\
\texttt{duracion} & Segundos. \\
\end{tabular}}\par\smallskip
\noindent\textbf{Devuelve:} True si terminó; False si se detuvo con B.\par
\noindent\textbf{Llama a:} \hyperref[fn:funciones_rover:parar]{\texttt{parar}}\par
\lstinputlisting[firstline=84,lastline=94,firstnumber=84]{../../software/robot/funciones_rover.py}
\subsubsection{\texttt{avanzar\_\allowbreak{}cm(cm, rpm, rumbo\_\allowbreak{}obj)}}\label{fn:funciones_rover:avanzar-cm}
\noindent Avanza una distancia por encoders, sosteniendo el rumbo con el giroscopio.

Arranca y frena suave (RAMPA\_S).

\par\noindent{\small\begin{tabular}{@{}p{0.2\linewidth}p{0.74\linewidth}@{}}
\textbf{Par\'ametro} & \textbf{Descripci\'on} \\
\texttt{cm} & cm (negativo = atrás). \\
\texttt{rpm} & rpm de avance. \\
\texttt{rumbo\_\allowbreak{}obj} & Rumbo a sostener (por omisión, el de ahora). \\
\end{tabular}}\par\smallskip
\noindent\textbf{Devuelve:} True si recorrió la distancia.\par
\noindent\textbf{Llama a:} \hyperref[fn:funciones_rover:orugas-cm]{\texttt{orugas\_\allowbreak{}cm}}, \hyperref[fn:funciones_rover:parar]{\texttt{parar}}, \hyperref[fn:funciones_rover:rumbo]{\texttt{rumbo}}, \hyperref[fn:funciones_rover:limitar]{\texttt{limitar}}, \hyperref[fn:funciones_rover:motores]{\texttt{motores}}\par
\lstinputlisting[firstline=103,lastline=133,firstnumber=103]{../../software/robot/funciones_rover.py}
\subsubsection{\texttt{girar(angulo, rpm\_\allowbreak{}max, timeout)}}\label{fn:funciones_rover:girar}
\noindent Gira en el lugar con el giroscopio (+ = derecha).

Frena al acercarse y termina corrigiendo hasta quedar a menos de TOL\_GIRO.

\par\noindent{\small\begin{tabular}{@{}p{0.2\linewidth}p{0.74\linewidth}@{}}
\textbf{Par\'ametro} & \textbf{Descripci\'on} \\
\texttt{angulo} & Grados a girar. \\
\texttt{rpm\_\allowbreak{}max} & rpm máximas. \\
\texttt{timeout} & s máximos. \\
\end{tabular}}\par\smallskip
\noindent\textbf{Devuelve:} True si llegó al ángulo.\par
\noindent\textbf{Llama a:} \hyperref[fn:funciones_rover:parar]{\texttt{parar}}, \hyperref[fn:funciones_rover:rumbo]{\texttt{rumbo}}, \hyperref[fn:funciones_rover:limitar]{\texttt{limitar}}, \hyperref[fn:funciones_rover:motores]{\texttt{motores}}\par
\lstinputlisting[firstline=142,lastline=166,firstnumber=142]{../../software/robot/funciones_rover.py}
\subsubsection{\texttt{girar\_\allowbreak{}a(rumbo\_\allowbreak{}obj, rpm\_\allowbreak{}max)}}\label{fn:funciones_rover:girar-a}
\noindent Gira hasta un rumbo absoluto (0 = como se encendió o se puso en 0 con B).

\par\noindent{\small\begin{tabular}{@{}p{0.2\linewidth}p{0.74\linewidth}@{}}
\textbf{Par\'ametro} & \textbf{Descripci\'on} \\
\texttt{rumbo\_\allowbreak{}obj} & Rumbo objetivo en grados. \\
\texttt{rpm\_\allowbreak{}max} & rpm máximas. \\
\end{tabular}}\par\smallskip
\noindent\textbf{Devuelve:} True si llegó.\par
\noindent\textbf{Llama a:} \hyperref[fn:funciones_rover:girar]{\texttt{girar}}, \hyperref[fn:funciones_rover:rumbo]{\texttt{rumbo}}\par
\lstinputlisting[firstline=173,lastline=179,firstnumber=173]{../../software/robot/funciones_rover.py}
\subsubsection{\texttt{mensaje(texto, tam, espera)}}\label{fn:funciones_rover:mensaje}
\noindent Muestra un texto centrado en la pantalla.

\par\noindent{\small\begin{tabular}{@{}p{0.2\linewidth}p{0.74\linewidth}@{}}
\textbf{Par\'ametro} & \textbf{Descripci\'on} \\
\texttt{texto} & Texto. \\
\texttt{tam} & Tamaño de letra. \\
\texttt{espera} & Segundos a esperar después. \\
\end{tabular}}\par\smallskip
\lstinputlisting[firstline=187,lastline=191,firstnumber=187]{../../software/robot/funciones_rover.py}
\subsubsection{\texttt{menu()}}\label{fn:funciones_rover:menu}
\noindent Muestra los botones en la pantalla.

\lstinputlisting[firstline=195,lastline=197,firstnumber=195]{../../software/robot/funciones_rover.py}
\subsubsection{\texttt{correr(nombre, accion)}}\label{fn:funciones_rover:correr}
\noindent Corre una prueba si no hay otra corriendo (la detiene B).

\par\noindent{\small\begin{tabular}{@{}p{0.2\linewidth}p{0.74\linewidth}@{}}
\textbf{Par\'ametro} & \textbf{Descripci\'on} \\
\texttt{nombre} & Nombre que se muestra en la pantalla. \\
\texttt{accion} & Función sin argumentos que hace la prueba. \\
\end{tabular}}\par\smallskip
\noindent\textbf{Llama a:} \hyperref[fn:funciones_rover:mensaje]{\texttt{mensaje}}, \hyperref[fn:funciones_rover:parar]{\texttt{parar}}, \hyperref[fn:funciones_rover:menu]{\texttt{menu}}\par
\lstinputlisting[firstline=203,lastline=214,firstnumber=203]{../../software/robot/funciones_rover.py}
\subsubsection{\texttt{on\_\allowbreak{}start()}}\label{fn:funciones_rover:on-start}
\noindent Al encender: muestra los botones de las pruebas.

\noindent\textbf{Llama a:} \hyperref[fn:funciones_rover:mensaje]{\texttt{mensaje}}, \hyperref[fn:funciones_rover:menu]{\texttt{menu}}\par
\lstinputlisting[firstline=221,lastline=224,firstnumber=221]{../../software/robot/funciones_rover.py}
\subsubsection{\texttt{secuencia()}}\label{fn:funciones_rover:secuencia}
\noindent Secuencia de prueba del botón A (ver la lista de pasos al inicio del archivo).

\noindent\textbf{Devuelve:} True si terminó todos los pasos.\par
\noindent\textbf{Llama a:} \hyperref[fn:funciones_rover:avanzar-cm]{\texttt{avanzar\_\allowbreak{}cm}}, \hyperref[fn:funciones_rover:girar]{\texttt{girar}}\par
\lstinputlisting[firstline=234,lastline=241,firstnumber=234]{../../software/robot/funciones_rover.py}
\subsubsection{\texttt{test\_\allowbreak{}secuencia()}}\label{fn:funciones_rover:test-secuencia}
\noindent Botón A: secuencia de prueba.

\noindent\textbf{Llama a:} \hyperref[fn:funciones_rover:correr]{\texttt{correr}}\par
\lstinputlisting[firstline=245,lastline=247,firstnumber=245]{../../software/robot/funciones_rover.py}
\subsubsection{\texttt{girar\_\allowbreak{}45\_\allowbreak{}derecha()}}\label{fn:funciones_rover:girar-45-derecha}
\noindent Joystick arriba: gira 45\textdegree{} a la derecha.

\noindent\textbf{Llama a:} \hyperref[fn:funciones_rover:correr]{\texttt{correr}}, \hyperref[fn:funciones_rover:girar]{\texttt{girar}}\par
\lstinputlisting[firstline=254,lastline=256,firstnumber=254]{../../software/robot/funciones_rover.py}
\subsubsection{\texttt{girar\_\allowbreak{}90\_\allowbreak{}izquierda()}}\label{fn:funciones_rover:girar-90-izquierda}
\noindent Joystick abajo: gira 90\textdegree{} a la izquierda.

\noindent\textbf{Llama a:} \hyperref[fn:funciones_rover:correr]{\texttt{correr}}, \hyperref[fn:funciones_rover:girar]{\texttt{girar}}\par
\lstinputlisting[firstline=260,lastline=262,firstnumber=260]{../../software/robot/funciones_rover.py}
\subsubsection{\texttt{girar\_\allowbreak{}30\_\allowbreak{}derecha()}}\label{fn:funciones_rover:girar-30-derecha}
\noindent Joystick izquierda: gira 30\textdegree{} a la derecha.

\noindent\textbf{Llama a:} \hyperref[fn:funciones_rover:correr]{\texttt{correr}}, \hyperref[fn:funciones_rover:girar]{\texttt{girar}}\par
\lstinputlisting[firstline=266,lastline=268,firstnumber=266]{../../software/robot/funciones_rover.py}
\subsubsection{\texttt{girar\_\allowbreak{}180()}}\label{fn:funciones_rover:girar-180}
\noindent Joystick derecha: gira 180\textdegree{}.

\noindent\textbf{Llama a:} \hyperref[fn:funciones_rover:correr]{\texttt{correr}}, \hyperref[fn:funciones_rover:girar]{\texttt{girar}}\par
\lstinputlisting[firstline=272,lastline=274,firstnumber=272]{../../software/robot/funciones_rover.py}
\subsubsection{\texttt{girar\_\allowbreak{}360()}}\label{fn:funciones_rover:girar-360}
\noindent Joystick centro: gira 360\textdegree{}.

\noindent\textbf{Llama a:} \hyperref[fn:funciones_rover:correr]{\texttt{correr}}, \hyperref[fn:funciones_rover:girar]{\texttt{girar}}\par
\lstinputlisting[firstline=278,lastline=280,firstnumber=278]{../../software/robot/funciones_rover.py}
\subsubsection{\texttt{detener()}}\label{fn:funciones_rover:detener}
\noindent Botón B: detiene la prueba en curso; con el robot quieto pone el rumbo en 0.

\noindent\textbf{Llama a:} \hyperref[fn:funciones_rover:parar]{\texttt{parar}}, \hyperref[fn:funciones_rover:mensaje]{\texttt{mensaje}}, \hyperref[fn:funciones_rover:menu]{\texttt{menu}}\par
\lstinputlisting[firstline=287,lastline=297,firstnumber=287]{../../software/robot/funciones_rover.py}
\subsubsection{\texttt{main()}}\label{fn:funciones_rover:main}
\noindent Ciclo en reposo: muestra el rumbo del giroscopio en grande (prueba a mano).

\noindent\textbf{Llama a:} \hyperref[fn:funciones_rover:rumbo]{\texttt{rumbo}}\par
\lstinputlisting[firstline=304,lastline=309,firstnumber=304]{../../software/robot/funciones_rover.py}
